\documentclass[12pt,a4paper]{article}

% Поддержка русского языка
\usepackage[utf8]{inputenc}
\usepackage[T2A]{fontenc}
\usepackage[russian]{babel}

\usepackage{tikz}
\usetikzlibrary{hobby, arrows.meta}
\usetikzlibrary{decorations.pathreplacing}
\usetikzlibrary{calc,intersections,decorations.pathreplacing,patterns.meta}
% Современный шрифт, поддерживающий кириллицу (Charter)
\usepackage{XCharter}

% Математика
\usepackage{amsmath, amsfonts, amssymb, amsthm}

% Улучшенный межстрочный интервал
\usepackage{setspace}
\onehalfspacing

% Настройка полей страницы
\usepackage{geometry}
\geometry{top=2.5cm, bottom=2.5cm, left=2.5cm, right=2.5cm}

% --- Настройка современных пастельных цветов (HEX) ---
\usepackage{xcolor}
% Цвета для определений (сtиний)
\definecolor{defbg}{HTML}{F0F8FF}    % Очень светлый голубой фон
\definecolor{defframe}{HTML}{5DADE2} % Мягкий синий для акцента

% Цвета для теорем (зеленый)
\definecolor{thmbg}{HTML}{F4FCF4}    % Очень светлый мятный фон
\definecolor{thmframe}{HTML}{7DCEA0} % Мягкий зеленый для акцента

% Цвета для лемм (желтый/оранжевый)
\definecolor{lembg}{HTML}{FFF9F0}    % Кремовый фон
\definecolor{lemframe}{HTML}{F5B041} % Мягкий оранжевый для акцента

% Цвета для примеров (серый)
\definecolor{exbg}{HTML}{F8F9F9}     % Светло-серый фон
\definecolor{exframe}{HTML}{BFC9CA}  % Пепельно-серый для акцента

% Цвета для текста и ссылок
\definecolor{darkblue}{HTML}{154360}

% --- Настройка tcolorbox в современном стиле "Material Design" ---
\usepackage{tcolorbox}
\tcbuselibrary{skins, breakable, theorems}

% Общие настройки для всех блоков (скрываем рамку, делаем полосу слева)
\tcbset{
    modernbox/.style={
        enhanced,
        breakable,
        frame hidden,
        boxrule=0pt,
        arc=2mm,
        top=3mm, bottom=3mm, left=4mm, right=4mm,
        fonttitle=\bfseries\sffamily,
        coltitle=black!80,
        attach title to upper={\vspace{0.5em}\par}, % Заголовок внутри блока над текстом
    }
}

% Определение
\newtcbtheorem{Definition}{Определение}{
    modernbox,
    colback=defbg,
    borderline west={3pt}{0pt}{defframe},
    coltitle=defframe!60!black
}{def}

% Теорема
\newtcbtheorem{Theorem}{Теорема}{
    modernbox,
    colback=thmbg,
    borderline west={3pt}{0pt}{thmframe},
    coltitle=thmframe!60!black
}{thm}

% Лемма
\newtcbtheorem{Lemma}{Лемма}{
    modernbox,
    colback=lembg,
    borderline west={3pt}{0pt}{lemframe},
    coltitle=lemframe!60!black
}{lem}

% Пример
\newtcolorbox{Example}[1][]{
    modernbox,
    colback=exbg,
    borderline west={3pt}{0pt}{exframe},
    title=\textbf{Пример(ы) #1},
    coltitle=exframe!60!black
}

% Среда для доказательств (делаем курсивный заголовок и небольшой отступ слева)
\newenvironment{Proof}[1][Доказательство]{\vspace{0.5em}\noindent\textit{#1.}\hspace{0.5em}\ignorespaces}{\hfill$\blacksquare$\vspace{1em}}

% Красивые гиперссылки
\usepackage{hyperref}
\hypersetup{
    colorlinks=true,
    linkcolor=darkblue,
    filecolor=magenta,      
    urlcolor=darkblue,
    pdftitle={Конспект по теории меры},
}
\begin{document}
Если заметили ошибку пишите в tg: @fadmar
\tableofcontents
\section{§ Система множеств, алгебра, $\sigma$-алгебра, кольца}

\begin{Definition}{Симметричная система}{}
Система множеств $\mathcal{U}$ называется \textbf{симметричной}, если из того, что $A \in \mathcal{U}$, следует, что $A^c \in \mathcal{U}$ (где $A^c$ — дополнение множества $A$).
\end{Definition}

\begin{Example}
\begin{itemize}
    \item $\mathcal{U} = 2^X$ — булеан (множество всех подмножеств множества $X$).
    \item $\mathcal{U} = \{\emptyset, X\}$.
    \item $\mathcal{U}$ — множество всех ограниченных множеств в $\mathbb{R}$ и их дополнений.
\end{itemize}
\end{Example}

\begin{Theorem}{Эквивалентность свойств в симметричной системе}{}
Пусть $\mathcal{U}$ — симметричная система множеств. Тогда справедливы следующие эквивалентности:
$$ \sigma_0 \Leftrightarrow \delta_0, \quad \sigma \Leftrightarrow \delta $$
где свойства определяются следующим образом:
\begin{align*}
    \sigma_0 &: A, B \in \mathcal{U} \Rightarrow A \cup B \in \mathcal{U} \\
    \delta_0 &: A, B \in \mathcal{U} \Rightarrow A \cap B \in \mathcal{U} \\
    \sigma &: A_i \in \mathcal{U}, \; i \in \{1, 2, \dots\} \Rightarrow \bigcup_{i=1}^{\infty} A_i \in \mathcal{U} \\
    \delta &: A_i \in \mathcal{U}, \; i \in \{1, 2, \dots\} \Rightarrow \bigcap_{i=1}^{\infty} A_i \in \mathcal{U}
\end{align*}
\end{Theorem}

\begin{Proof}
Докажем импликацию $\sigma_0 \Rightarrow \delta_0$. Пусть $A, B \in \mathcal{U}$. 
По законам де Моргана справедливо равенство:
$$ A \cap B = (A^c \cup B^c)^c $$
Так как система $\mathcal{U}$ симметрична, множества $A^c$ и $B^c$ принадлежат $\mathcal{U}$. Из свойства $\sigma_0$ следует, что их объединение $(A^c \cup B^c)$ также принадлежит $\mathcal{U}$. Вновь применяя свойство симметричности, получаем, что дополнение этого объединения, то есть $(A^c \cup B^c)^c$, принадлежит $\mathcal{U}$. Следовательно, $A \cap B \in \mathcal{U}$.

Остальные эквивалентности доказываются аналогичным образом путём перехода к дополнениям (дальнейшее очевидно).
\end{Proof}

\begin{Definition}{Алгебра и $\sigma$-алгебра}{}
Непустая симметричная система множеств $\mathcal{U}$ называется \textbf{алгеброй} (соответственно, \textbf{$\sigma$-алгеброй}), если $\mathcal{U}$ обладает любым из равносильных свойств: $\sigma_0 \Leftrightarrow \delta_0$ (соответственно, $\sigma \Leftrightarrow \delta$).
\end{Definition}

\begin{Example}[и контрпример]
Рассмотрим систему множеств $\mathcal{U}$, состоящую из всех ограниченных подмножеств $\mathbb{R}$ и их дополнений. Это семейство является \textbf{алгеброй}, но \textbf{не является $\sigma$-алгеброй}.

\textbf{Обоснование:} Пусть $A_n = [n, n+1]$. Каждое $A_n$ ограничено, а значит $A_n \in \mathcal{U}$. Однако их счетное объединение:
$$ \bigcup_{n=1}^{\infty} [n, n+1] = [1, +\infty) $$
не принадлежит $\mathcal{U}$, так как множество $[1, +\infty)$ не является ограниченным, и его дополнение $(-\infty, 1)$ также не ограничено в $\mathbb{R}$.
\end{Example}

\begin{Lemma}{Свойства алгебры и $\sigma$-алгебры}{}
Пусть $\mathcal{U}$ — алгебра (или $\sigma$-алгебра) подмножеств базового множества $X$. Тогда выполняются следующие свойства:
\begin{enumerate}
    \item $\emptyset \in \mathcal{U}, \; X \in \mathcal{U}$.
    \item $A, B \in \mathcal{U} \Rightarrow A \setminus B \in \mathcal{U}$.
    \item Всякая $\sigma$-алгебра является алгеброй.
\end{enumerate}
\end{Lemma}

\begin{Proof} \\
1) Так как $\mathcal{U}$ непустая, то $\exists A \in \mathcal{U} \Rightarrow A^c \in \mathcal{U}$. \\
Тогда $A \cap A^c = \emptyset \in \mathcal{U}$. \\
Следовательно, $\emptyset^c = X \in \mathcal{U}$.

2) $A \setminus B = A \cap B^c \in \mathcal{U}$ (так как $B^c \in \mathcal{U}$ и алгебра замкнута относительно пересечений).

3) Очевидно (по определению $\sigma$-алгебры).
\end{Proof}

\begin{Lemma}{О пересечении семейства алгебр}{}
Пусть $\mathcal{U}_\alpha, \; \alpha \in \mathcal{A}$ — семейство алгебр (или $\sigma$-алгебр), являющихся подмножествами $2^X$. Тогда их пересечение $\bigcap\limits_{\alpha \in \mathcal{A}} \mathcal{U}_\alpha$ также является алгеброй (или $\sigma$-алгеброй).
\end{Lemma}

\begin{Proof}
1) $\emptyset \in \mathcal{U}_\alpha \Rightarrow \emptyset \in \bigcap\limits_{\alpha \in \mathcal{A}} \mathcal{U}_\alpha$.

2) ($A \in \bigcap\limits_{\alpha \in \mathcal{A}} \mathcal{U}_\alpha \Rightarrow (\forall \alpha \in \mathcal{A} : A \in \mathcal{U}_\alpha) \Rightarrow (\forall \alpha \in \mathcal{A} : A^c \in \mathcal{U}_\alpha) \Rightarrow A^c \in \bigcap\limits_{\alpha \in \mathcal{A}} \mathcal{U}_\alpha$.

3) $A, B \in \bigcap\limits_{\alpha \in \mathcal{A}} \mathcal{U}_\alpha \Rightarrow (\forall \alpha \in \mathcal{A} : A, B \in \mathcal{U}_\alpha) \Rightarrow (\forall \alpha \in \mathcal{A} : A \cup B \in \mathcal{U}_\alpha) \Rightarrow A \cup B \in \bigcap\limits_{\alpha \in \mathcal{A}} \mathcal{U}_\alpha$.
\end{Proof}

\begin{Theorem}{О существовании минимальной $\sigma$-алгебры}{}
Для любой системы подмножеств $\mathcal{E}$ множества $X$ существует минимальная $\sigma$-алгебра, содержащая все элементы $\mathcal{E}$.
\end{Theorem}

\begin{Proof}
Существует $\mathcal{U} = 2^X$ — $\sigma$-алгебра, содержащая все элементы $\mathcal{E}$. Согласно предыдущей лемме, искомая минимальная $\sigma$-алгебра — это пересечение всех $\sigma$-алгебр, содержащих все элементы $\mathcal{E}$.
\end{Proof}

\begin{Definition}{Борелевская оболочка}{}
Минимальная $\sigma$-алгебра из предыдущей теоремы называется \textbf{борелевской оболочкой} системы множеств $\mathcal{E}$ и обозначается $\mathcal{B}(\mathcal{E})$.
\end{Definition}

% ==========================================

\begin{Definition}{Кольцо и $\sigma$-кольцо}{}
\textbf{Кольцом} $\mathcal{K}$ подмножеств множества $X$ называется непустая система множеств, замкнутая относительно объединения, пересечения и разности.

Кольцо $\mathcal{K}$ называется \textbf{$\sigma$-кольцом} (или \textbf{$\delta$-кольцом}), если оно замкнуто относительно объединения (или пересечения) не более чем счётного числа множеств.
\end{Definition}

\begin{Example}
Пусть $X = \{a, b\}$. Тогда $\mathcal{K} = \{\emptyset, \{a\}\}$ — кольцо.
\end{Example}

\begin{Lemma}{Свойства кольца}{}
1) $\emptyset \in \mathcal{K}$. \\
2) Кольцо $\mathcal{K}$ является алгеброй $\Leftrightarrow X \in \mathcal{K}$.
\end{Lemma}

\begin{Proof}
1) Пусть $A \in \mathcal{K}$. Тогда $A \setminus A = \emptyset \in \mathcal{K}$. \\
2) Очевидно (из определений кольца и алгебры).
\end{Proof}

\section{§ Полукольцо и его свойства}

\begin{Definition}{Полукольцо}{}
Система подмножеств $\mathcal{P}$ множества $X$ называется \textbf{полукольцом}, если:
\begin{enumerate}
    \item $\emptyset \in \mathcal{P}$
    \item $A, B \in \mathcal{P} \Rightarrow A \cap B \in \mathcal{P}$
    \item $A, B \in \mathcal{P} \Rightarrow A \setminus B = \bigsqcup\limits_{i=1}^n Q_i$, где $Q_i \in \mathcal{P}$ (конечное дизъюнктное объединение).
\end{enumerate}
\end{Definition}

\begin{Example}
Система полуинтервалов $[a, b)$, где $a \le b$ — образует полукольцо. (Можно считать, что $a, b \in \mathbb{Q}$).
\end{Example}

\begin{Lemma}{О дизъюнктном представлении объединения («про нарезку»)}{}
Пусть $A_i$ — последовательность множеств. Тогда их объединение можно представить в виде дизъюнктного объединения:
$$ \bigcup_{i=1}^{\infty} A_i = \bigsqcup_{i=1}^{\infty} \left( A_i \setminus \bigcup_{j=0}^{i-1} A_j \right), \quad \text{где } A_0 = \emptyset $$
\end{Lemma}

\begin{Proof}
Обозначим $B_i = A_i \setminus \bigcup_{j=0}^{i-1} A_j$. Докажем, что $B_i$ дизъюнктны. \\
Пусть $\exists i > k$. Тогда $B_k \subset A_k$. По построению $B_i \cap A_k = \emptyset$, следовательно, $B_i \cap B_k = \emptyset$.

Докажем равенство множеств: включение $\bigcup B_i \subset \bigcup A_i$ очевидно. \\
В обратную сторону: пусть $x \in \bigcup_{i=1}^{\infty} A_i$. Значит, существует наименьший индекс $i_0$ такой, что $x \in A_{i_0}$. По определению множества $B_{i_0}$, $x$ не принадлежит предыдущим $A_j$, а значит $x \in B_{i_0}$. Включение доказано.
\end{Proof}

\begin{center}
\begin{tikzpicture}[scale=1.1, >=latex]
    % --- Определение пастельных цветов ---
    \definecolor{blobB1}{HTML}{5DADE2} 
    \definecolor{blobB2}{HTML}{AED6F1} 
    \definecolor{blobB3}{HTML}{EBF5FB} 
    \definecolor{darktext}{HTML}{154360} 

    % --- Слой 3 (самый большой, B3) ---
    \filldraw[fill=blobB3, draw=darktext, thick] plot [smooth cycle, tension=0.7] coordinates {
        (0, 3) (2.5, 2.5) (3.5, 0) (3, -2) (0, -2.8) (-3, -2) (-3.5, 0) (-2.5, 2.5)
    };

    % --- Слой 2 (средний, B2) ---
    \filldraw[fill=blobB2, draw=darktext, thick] plot [smooth cycle, tension=0.7] coordinates {
        (0, 2) (1.8, 1.3) (2.3, 0) (1.6, -1.4) (0, -2) (-1.6, -1.4) (-2.3, 0) (-1.8, 1.3)
    };

    % --- Слой 1 (центр, B1) ---
    \filldraw[fill=blobB1, draw=darktext, thick] plot [smooth cycle, tension=0.7] coordinates {
        (0, 1) (1, 0.7) (1.4, 0) (1, -0.7) (0, -1) (-1, -0.7) (-1.4, 0) (-1, 0.7)
    };

    % --- Подписи исходных множеств A_i (снаружи) ---
    \node[darktext, font=\bfseries] (A3) at (4.5, 1.5) {$A_3$};
    \draw[->, darktext, thick] (A3) -- (3.2, 1);

    \node[darktext, font=\bfseries] (A2) at (3.5, 2.8) {$A_2$};
    \draw[->, darktext, thick] (A2) -- (2.2, 1.5);

    \node[darktext, font=\bfseries] (A1) at (2.2, 3.2) {$A_1$};
    \draw[->, darktext, thick] (A1) -- (1.2, 0.8);

    % --- Подписи дизъюнктных кусков B_i (внутри) ---
    \node[darktext, font=\Large\bfseries] at (0, 0) {$B_1$};
    \node[darktext, font=\Large\bfseries] at (0, -1.5) {$B_2$};
    \node[darktext, font=\Large\bfseries] at (0, -2.4) {$B_3$};

    % --- Легенда снизу ---
    \node[darktext, align=left, anchor=north] at (0, -3.2) {
        $B_1 = A_1$  \\
        $B_2 = A_2 \setminus A_1$  \\
        $B_3 = A_3 \setminus (A_1 \cup A_2)$ 
    };
\end{tikzpicture}
\end{center}

\begin{Theorem}{Свойства полукольца}{}
Пусть $P, P_1, \dots, P_n \in \mathcal{P}$, где $\mathcal{P}$ — полукольцо. Тогда:
\begin{enumerate}
    \item $P \setminus \bigcup\limits_{i=1}^n P_i = \bigsqcup\limits_{i=1}^m Q_i$, где $Q_i \in \mathcal{P}$.
    \item $\bigcup\limits_{i=1}^n P_i = \bigsqcup\limits_{i=1}^n \bigsqcup\limits_{j=1}^m Q_{ij}$, где $Q_{ij} \in \mathcal{P}$ и $Q_{ij} \subset P_i$.
    \item В пункте 2 число множеств $n$ может быть равно $\infty$.
\end{enumerate}
\end{Theorem}

\begin{Proof}
\textbf{1) Индукция по $n$:} \\
\textit{База:} Для $n=1$ следует из определения полукольца. \\
\textit{Индукционный переход (от $n$ к $n+1$):}
$$ P \setminus \bigcup_{i=1}^{n+1} P_i = \left( P \setminus \bigcup_{i=1}^n P_i \right) \setminus P_{n+1} = \left( \bigsqcup_{j=1}^k Q_j \right) \setminus P_{n+1} = \bigsqcup_{j=1}^k (Q_j \setminus P_{n+1}) = \bigsqcup_{j=1}^k \bigsqcup_{l=1}^m Q_{jl} $$
где $Q_{jl} \in \mathcal{P}$.

\begin{center}
\begin{tikzpicture}[scale=1.3, >=latex]
    % --- Определение пастельных цветов ---
    \definecolor{colorP}{HTML}{EBF5FB}     % Светло-голубой для оставшейся части (Q_i)
    \definecolor{colorSub}{HTML}{FADBD8}   % Светло-красный для вырезаемых P_i
    \definecolor{darktext}{HTML}{154360}   % Темно-синий для текста и линий

    % 1. Рисуем базовый большой прямоугольник P (сначала заливаем всё голубым)
    \filldraw[fill=colorP, draw=darktext, thick] (0,0) rectangle (6,4);
    
    % 2. Вырезаем (закрашиваем красным) прямоугольники P1 и P2
    \filldraw[fill=colorSub, draw=darktext, thick] (1,1) rectangle (3,3);
    \filldraw[fill=colorSub, draw=darktext, thick] (2,2) rectangle (5,3.5);
    
    % Обводим контуры P1 и P2 еще раз, чтобы граница была четкой
    \draw[darktext, thick] (1,1) rectangle (3,3);
    \draw[darktext, thick] (2,2) rectangle (5,3.5);

    % 3. Делаем "нарезку" пунктирными линиями.
    % Продлеваем стороны P1 и P2 до краев большого прямоугольника P
    
    % Вертикальные разрезы
    \draw[dashed, darktext] (1,0) -- (1,4);
    \draw[dashed, darktext] (2,0) -- (2,4);
    \draw[dashed, darktext] (3,0) -- (3,4);
    \draw[dashed, darktext] (5,0) -- (5,4);

    % Горизонтальные разрезы
    \draw[dashed, darktext] (0,1) -- (6,1);
    \draw[dashed, darktext] (0,2) -- (6,2);
    \draw[dashed, darktext] (0,3) -- (6,3);
    \draw[dashed, darktext] (0,3.5) -- (6,3.5);

    % --- Подписи множеств ---
    
    % Подписываем вырезанные множества
    \node[darktext, font=\Large\bfseries] at (1.5, 2.5) {$P_1$};
    \node[darktext, font=\Large\bfseries] at (4, 2.75) {$P_2$};
    
    % Подписываем сам прямоугольник P
    \node[darktext, font=\Large\bfseries, above right] at (0,4) {$P$};

    % Расставляем подписи для оставшихся "кусочков" Q_j (голубая зона)
    \node[darktext] at (0.5, 0.5) {$Q_1$};
    \node[darktext] at (1.5, 0.5) {$Q_2$};
    \node[darktext] at (2.5, 0.5) {$Q_3$};
    \node[darktext] at (4.0, 0.5) {$Q_4$};
    \node[darktext] at (5.5, 0.5) {$Q_5$};
    
    \node[darktext] at (0.5, 1.5) {$Q_6$};
    \node[darktext] at (4.0, 1.5) {$Q_7$};
    \node[darktext] at (5.5, 1.5) {$Q_8$};
    
    \node[darktext] at (0.5, 2.5) {$Q_9$};
    \node[darktext] at (5.5, 2.5) {$Q_{10}$};
    
    \node[darktext] at (0.5, 3.25) {$\dots$};
    \node[darktext] at (1.5, 3.25) {$Q_k$};
    \node[darktext] at (5.5, 3.75) {$Q_m$};

    % --- Легенда снизу ---
    \node[darktext, align=left, anchor=north] at (3, -0.2) {
        $P \setminus (P_1 \cup P_2) = \bigsqcup\limits_{j=1}^m Q_j$ \\
        \small Оставшаяся площадь (голубая) разбивается на конечную сетку \\
        \small из непересекающихся прямоугольников $Q_j \in \mathcal{P}$.
    };
\end{tikzpicture}
\end{center}

\textbf{2) Применяем лемму «про нарезку»:}
$$ \bigcup_{i=1}^n P_i = \bigsqcup_{i=1}^n \left( P_i \setminus \bigcup_{j=0}^{i-1} P_j \right) $$
где $P_0 = \emptyset$. По доказанному пункту 1, разность в скобках можно представить как дизъюнктное объединение элементов полукольца:
$$ P_i \setminus \bigcup_{j=0}^{i-1} P_j = \bigsqcup_{j=1}^{m_i} Q_{ij}, \quad \text{где } Q_{ij} \in \mathcal{P}, \; Q_{ij} \subset P_i $$
Отсюда следует требуемое равенство.
\end{Proof}

\begin{Lemma}{Об упрощении структуры объединения}{}
Пункт 2 предыдущей теоремы можно упростить до вида: $\bigcup\limits_{i=1}^n P_i = \bigsqcup\limits_{j=1}^m Q_j$, причем для любой пары $(i, j)$ справедливо: либо $Q_j \subset P_i$, либо $Q_j \cap P_i = \emptyset$.
\end{Lemma}

\begin{Proof}
Доказывается по индукции. Для двух множеств:
$$ P_1 \cup P_2 = (P_1 \setminus P_2) \cup (P_1 \cap P_2) \cup (P_2 \setminus P_1) = \left( \bigsqcup\limits_{j=1}^{n_1} Q_j \right) \sqcup (P_1 \cap P_2) \sqcup \left( \bigsqcup\limits_{j=1}^{n_2} Q'_j \right) $$
Все эти элементы лежат в полукольце и попарно не пересекаются.
\end{Proof}

\section{§ Объём и его свойства}

\begin{Definition}{Объём}{}
Функция $\varphi : \mathcal{P} \to [0, +\infty]$ называется \textbf{объёмом}, если:
\begin{enumerate}
    \item $\varphi(\emptyset) = 0$
    \item Если $A, B \in \mathcal{P}$ и их дизъюнктное объединение $A \sqcup B \in \mathcal{P}$, то $\varphi(A \sqcup B) = \varphi(A) + \varphi(B)$.
\end{enumerate}
\end{Definition}

\begin{Example}
\begin{enumerate}
    \item Пусть $\mathcal{P}$ — полукольцо полуинтервалов $[a, b)$. Функция $\varphi([a, b)) = b - a$ является объёмом (\textbf{мера Лебега}).
    \item На том же полукольце $\mathcal{P}$, пусть $f$ — возрастающая функция. Тогда $\varphi([a, b)) = f(b) - f(a)$ является объёмом (\textbf{мера Лебега-Стилтьеса}).
    \item Пусть $\Omega = \{\omega_1, \dots, \omega_n, \dots\}$ — пространство элементарных исходов (не более чем счётное), $\mathcal{U} = 2^\Omega$. Каждому исходу $\omega_i$ сопоставим вероятность $p_i \ge 0$, причем $\sum p_i = 1$. Тогда функция
    $$ \mathbb{P}(A) = \sum_{i: \omega_i \in A} p_i $$
    является объёмом (вероятностной мерой).
\end{enumerate}
\end{Example}

\setcounter{section}{3} % Продолжаем нумерацию параграфов, если нужно

\begin{Example}
Пусть $\varphi: \mathcal{P} \to [0, +\infty]$, $\varphi(\emptyset) = 0$, $\varphi(A \sqcup B) = \varphi(A) + \varphi(B)$.
Рассмотрим дискретное пространство элементарных исходов $\Omega = \{\omega_1, \dots, \omega_n, \dots\}$. Каждому исходу сопоставим вес $p_i \ge 0$. Пусть $\mathcal{U} = 2^\Omega$. Для $A \in \mathcal{U}$ определим:
$$ \varphi(A) = \sum\limits_{i: \omega_i \in A} p_i $$
\begin{enumerate}
    \item Считающая мера: если $p_i = 1$, то $\varphi(A)$ — количество элементов в $A$.
    \item Вероятностный объём: требуется, чтобы $\sum\limits_{i} p_i = 1$.
    \begin{itemize}
        \item $\Omega = \{\text{ребро, решка, орёл}\}$.
        \item \textbf{Распределение Пуассона:} $p_i = \frac{e^{-\lambda}\lambda^i}{i!}$, где $\lambda > 0$, $i \in \mathbb{N} \cup \{0\}$. Сумма $\sum\limits_{i=0}^{\infty} p_i = 1$.
        \item \textbf{Геометрическое распределение:} $p_i = p(1-p)^{i-1}$, $i \in \{1, 2, \dots\}$, $p \in (0, 1)$.
        \item \textbf{Биномиальное распределение} (вероятности Бернулли): $\Omega = \{\omega_1, \dots, \omega_n\}$, $p_i = C_n^i p^i (1-p)^{n-i}$, $p \in (0, 1)$.
    \end{itemize}
\end{enumerate}
\end{Example}

\begin{center}
\begin{tikzpicture}[>=latex, x=0.75cm, y=8cm] % Настройка масштаба: 1 единица X = 0.75 см, 1 единица Y (100%) = 8 см

    % --- Определение пастельных цветов ---
    \definecolor{darktext}{HTML}{154360}
    \definecolor{colorPoi}{HTML}{7DCEA0}  % Мятно-зеленый
    \definecolor{colorGeo}{HTML}{F5B041}  % Оранжево-желтый
    \definecolor{colorBin}{HTML}{5DADE2}  % Небесно-голубой

    % ==========================================
    % 1. Распределение Пуассона (СЛЕВА СВЕРХУ)
    % ==========================================
    \begin{scope}[shift={(0cm, 0cm)}] % Позиция левого верхнего графика
        % Оси
        \draw[->, thick, darktext] (-0.5, 0) -- (9, 0) node[right] {$i$};
        \draw[->, thick, darktext] (0, -0.05) -- (0, 0.55) node[above] {$p_i$};
        
        % Разметка оси X (от 0 до 8)
        \foreach \x in {0,1,2,3,4,5,6,7,8} {
            \draw[darktext] (\x, 0) -- (\x, -0.015) node[below=2pt] {\scriptsize $\x$};
        }
        % Разметка оси Y (от 0.1 до 0.5)
        \foreach \y in {0.1, 0.2, 0.3, 0.4, 0.5} {
            \draw[darktext] (0, \y) -- (-0.15, \y) node[left=1pt] {\scriptsize $\y$};
        }
        
        % Линии и точки (точные значения для lambda = 3)
        \foreach \x/\y in {0/0.050, 1/0.149, 2/0.224, 3/0.224, 4/0.168, 5/0.101, 6/0.050, 7/0.021, 8/0.008} {
            \draw[thick, colorPoi] (\x, 0) -- (\x, \y);
            \fill[colorPoi] (\x, \y) circle (2.5pt);
        }
        
        % Заголовки
        \node[darktext, font=\bfseries] at (4, 0.65) {Распределение Пуассона};
        \node[darktext] at (4, 0.58) {$p_i = \frac{e^{-\lambda}\lambda^i}{i!}$ \quad $(\lambda=3)$};
    \end{scope}

    % ==========================================
    % 2. Геометрическое распределение (СПРАВА СВЕРХУ)
    % ==========================================
    \begin{scope}[shift={(8.5cm, 0cm)}] % Сдвиг вправо на 8.5 см
        % Оси
        \draw[->, thick, darktext] (-0.5, 0) -- (9, 0) node[right] {$i$};
        \draw[->, thick, darktext] (0, -0.05) -- (0, 0.55) node[above] {$p_i$};
        
        % Разметка оси X
        \foreach \x in {0,1,2,3,4,5,6,7,8} {
            \draw[darktext] (\x, 0) -- (\x, -0.015) node[below=2pt] {\scriptsize $\x$};
        }
        % Разметка оси Y
        \foreach \y in {0.1, 0.2, 0.3, 0.4, 0.5} {
            \draw[darktext] (0, \y) -- (-0.15, \y) node[left=1pt] {\scriptsize $\y$};
        }
        
        % Линии и точки (точные значения для p = 0.4)
        % Начинается с 1, так как геометрическое распределение считает число попыток ДО первого успеха
        \foreach \x/\y in {1/0.400, 2/0.240, 3/0.144, 4/0.086, 5/0.052, 6/0.031, 7/0.018, 8/0.011} {
            \draw[thick, colorGeo] (\x, 0) -- (\x, \y);
            \fill[colorGeo] (\x, \y) circle (2.5pt);
        }
        
        % Заголовки
        \node[darktext, font=\bfseries] at (4, 0.65) {Геометрическое распр.};
        \node[darktext] at (4, 0.58) {$p_i = p(1-p)^{i-1}$ \quad $(p=0.4)$};
    \end{scope}

    % ==========================================
    % 3. Биномиальное распределение (ПО ЦЕНТРУ СНИЗУ)
    % ==========================================
    \begin{scope}[shift={(4.25cm, -6cm)}] % Сдвиг в центр и вниз на 6 см
        % Оси
        \draw[->, thick, darktext] (-0.5, 0) -- (9, 0) node[right] {$i$};
        \draw[->, thick, darktext] (0, -0.05) -- (0, 0.55) node[above] {$p_i$};
        
        % Разметка оси X
        \foreach \x in {0,1,2,3,4,5,6,7,8} {
            \draw[darktext] (\x, 0) -- (\x, -0.015) node[below=2pt] {\scriptsize $\x$};
        }
        % Разметка оси Y
        \foreach \y in {0.1, 0.2, 0.3, 0.4, 0.5} {
            \draw[darktext] (0, \y) -- (-0.15, \y) node[left=1pt] {\scriptsize $\y$};
        }
        
        % Линии и точки (точные значения для n = 8, p = 0.5)
        \foreach \x/\y in {0/0.004, 1/0.031, 2/0.109, 3/0.219, 4/0.273, 5/0.219, 6/0.109, 7/0.031, 8/0.004} {
            \draw[thick, colorBin] (\x, 0) -- (\x, \y);
            \fill[colorBin] (\x, \y) circle (2.5pt);
        }
        
        % Заголовки
        \node[darktext, font=\bfseries] at (4, 0.65) {Биномиальное распр.};
        \node[darktext] at (4, 0.58) {$p_i = C_n^i p^i (1-p)^{n-i}$ \quad $(p=0.5)$};
    \end{scope}

\end{tikzpicture}
\end{center}

\begin{Definition}{Конечный и $\sigma$-компактный объём}{}
\begin{itemize}
    \item Объём $\varphi$ называется \textbf{конечным}, если базовое множество $X \in \mathcal{P}$ и $\varphi(X) < +\infty$.
    \item Объём $\varphi$ называется \textbf{$\sigma$-компактным} (или $\sigma$-конечным), если множество $X$ можно представить как $X = \bigcup\limits_{i=1}^{\infty} X_i$, где для каждого $i$ выполняется $\varphi(X_i) < +\infty$.
\end{itemize}
\end{Definition}

\begin{Example}
На полукольце $\mathcal{P} = \{[a, b)\} : a \le b, \; a, b \in \mathbb{R}$ задан объём $\varphi([a, b)) = b - a$. Пространство $X = \mathbb{R}$ можно представить как $\mathbb{R} = \bigcup\limits_{n=-\infty}^{\infty} [n, n+1)$, где объём каждого полуинтервала равен 1. Значит, это $\sigma$-конечный объём.
\end{Example}

\begin{Theorem}{Свойства объёма}{}
Пусть $\varphi: \mathcal{P} \to [0, +\infty]$ — объём, и $P_1, P_2, \dots, P_n \in \mathcal{P}$. Тогда:
\begin{enumerate}
    \item \textbf{Монотонность:} Если $P \subset P_1$, то $\varphi(P) \le \varphi(P_1)$.
    \item \textbf{Усиленная монотонность:} Если $\bigsqcup\limits_{i=1}^n P_i \subset P$, то $\varphi(P) \ge \sum\limits_{i=1}^n \varphi(P_i)$.
    \item \textbf{Конечная полуаддитивность:} Если $P \subset \bigcup\limits_{i=1}^n P_i$, то $\varphi(P) \le \sum\limits_{i=1}^n \varphi(P_i)$.
\end{enumerate}
\end{Theorem}

\begin{Proof}[Доказательство]
\textbf{Пункт 2:} По свойству полукольца, разность можно представить как конечное дизъюнктное объединение элементов полукольца:
$$ P \setminus \bigsqcup\limits_{i=1}^n P_i = \bigsqcup\limits_{j=1}^m Q_j, \quad Q_j \in \mathcal{P} $$
Тогда множество $P$ представимо в виде:
$$ P = \left( \bigsqcup\limits_{i=1}^n P_i \right) \sqcup \left( \bigsqcup\limits_{j=1}^m Q_j \right) $$
В силу аддитивности объёма:
$$ \varphi(P) = \sum\limits_{i=1}^n \varphi(P_i) + \sum\limits_{j=1}^m \varphi(Q_j) \ge \sum\limits_{i=1}^n \varphi(P_i) $$
(что автоматически доказывает и пункт 1, если взять $n=1$).

\textbf{Пункт 3:} Обозначим $P_i' = P \cap P_i \in \mathcal{P}$. Тогда $P = \bigcup\limits_{i=1}^n P_i'$. По лемме об упрощении структуры объединения (свойство полукольца), это объединение можно нарезать на дизъюнктные куски:
$$ P = \bigsqcup\limits_{i=1}^n \bigsqcup\limits_{j=1}^{m_i} Q_{ij} $$
Причем $Q_{ij} \subset P_i'$ (а значит, и $Q_{ij} \subset P_i$).
Используя аддитивность и монотонность (из п.2):
$$ \varphi(P) = \sum\limits_{i=1}^n \sum\limits_{j=1}^{m_i} \varphi(Q_{ij}) \le \sum\limits_{i=1}^n \varphi(P_i') \le \sum\limits_{i=1}^n \varphi(P_i) $$
\end{Proof}

\begin{center}
\begin{tikzpicture}[scale=1.2, >=latex]

    % --- Определение цветов ---
    \definecolor{darktext}{HTML}{154360}
    \definecolor{colorP1}{HTML}{7DCEA0} % Зеленый для P_1
    \definecolor{colorP2}{HTML}{5DADE2} % Синий для P_2

    % ==========================================
    % 1. Покрывающие множества P1 и P2 (с полупрозрачностью)
    % ==========================================
    
    % Рисуем P1 (слева)
    \filldraw[fill=colorP1, fill opacity=0.4, draw=colorP1, thick] (0.5, 0.5) rectangle (4.5, 4.5);
    \node[colorP1!80!black, font=\Large\bfseries] at (1, 4.1) {$P_1$};

    % Рисуем P2 (справа). Там, где они наложатся, цвет станет темнее.
    \filldraw[fill=colorP2, fill opacity=0.4, draw=colorP2, thick] (2.5, 1) rectangle (6.5, 4);
    \node[colorP2!80!black, font=\Large\bfseries] at (6, 3.6) {$P_2$};


    % ==========================================
    % 2. Наше исходное множество P
    % ==========================================
    
    % Рисуем P толстой темно-синей линией (оно полностью внутри P1 U P2)
    \draw[draw=darktext, ultra thick] (1.5, 1.5) rectangle (5.5, 3.5);
    \node[darktext, font=\Large\bfseries] at (1.8, 3.2) {$P$};


    % ==========================================
    % 3. Дизъюнктная нарезка внутри P на куски Q_j
    % ==========================================
    
    % Граница, где начинается P2 (x = 2.5)
    \draw[dashed, darktext, thick] (2.5, 1.5) -- (2.5, 3.5);
    % Граница, где заканчивается P1 (x = 4.5)
    \draw[dashed, darktext, thick] (4.5, 1.5) -- (4.5, 3.5);

    % Подписываем куски Q_j
    \node[darktext, font=\Large\bfseries] at (2, 2.5) {$Q_1$};
    \node[darktext, font=\Large\bfseries] at (3.5, 2.5) {$Q_2$};
    \node[darktext, font=\Large\bfseries] at (5, 2.5) {$Q_3$};


    % ==========================================
    % 4. Легенда снизу с математической логикой
    % ==========================================
    \node[darktext, align=left, anchor=north] at (3.5, 0) {
        Множество $P$ нарезано на дизъюнктные части: $P = Q_1 \sqcup Q_2 \sqcup Q_3$ \\
        Тогда его объём: $\varphi(P) = \varphi(Q_1) + \varphi(Q_2) + \varphi(Q_3)$ \\
        \\
        Заметим, что $Q_1, Q_2 \subset P_1$, а значит $\varphi(P_1) \ge \varphi(Q_1) + \varphi(Q_2)$ \\
        Также $Q_2, Q_3 \subset P_2$, а значит $\varphi(P_2) \ge \varphi(Q_2) + \varphi(Q_3)$ \\
        \\
        Суммируя, получаем: $\varphi(P_1) + \varphi(P_2) \ge \varphi(Q_1) + 2\varphi(Q_2) + \varphi(Q_3) \ge \varphi(P)$ \\
        \small(Центральный кусок $Q_2$ считается дважды, поэтому возникает неравенство $\ge$)
    };

\end{tikzpicture}
\end{center}

\section{§ Произведение полуколец и объёмов}

\begin{Definition}{Произведение полуколец}{}
Пусть $\mathcal{P}$ и $\mathcal{Q}$ — полукольца. \textbf{Произведением} полуколец $\mathcal{P}$ и $\mathcal{Q}$ назовём систему множеств:
$$ \mathcal{P} \odot \mathcal{Q} = \{ P \times Q \mid P \in \mathcal{P}, \; Q \in \mathcal{Q} \} $$
где $\times$ обозначает декартово произведение множеств.
\end{Definition}

\begin{Lemma}{О структуре произведения полуколец}{}
Система $\mathcal{P} \odot \mathcal{Q}$ является полукольцом.
\end{Lemma}

\begin{Proof}
Проверим аксиомы полукольца:
1) $\emptyset = \emptyset \times \emptyset \in \mathcal{P} \odot \mathcal{Q}$.

2) Пусть $A, B \in \mathcal{P} \odot \mathcal{Q}$, где $A = P_1 \times Q_1$ и $B = P_2 \times Q_2$. Тогда:
$$ A \cap B = (P_1 \cap P_2) \times (Q_1 \cap Q_2) \in \mathcal{P} \odot \mathcal{Q} $$

3) Пусть $A, B \in \mathcal{P} \odot \mathcal{Q}$. Можно считать, что $B \subset A$ (иначе $A \setminus B = A \setminus (A \cap B)$).
Тогда $P_2 \subset P_1$ и $Q_2 \subset Q_1$. Разложим разности $P_1 \setminus P_2$ и $Q_1 \setminus Q_2$ в дизъюнктные объединения:
$$ P_1 \setminus P_2 = \bigsqcup\limits_{i=1}^n P_i', \quad Q_1 \setminus Q_2 = \bigsqcup\limits_{j=1}^m Q_j' $$
Тогда $P_1 = P_2 \sqcup \bigsqcup\limits_{i=1}^n P_i'$ и $Q_1 = Q_2 \sqcup \bigsqcup\limits_{j=1}^m Q_j'$. 
Подставляя это в декартово произведение $P_1 \times Q_1$, получим:
$$ P_1 \times Q_1 = (P_2 \times Q_2) \sqcup \left( \bigsqcup\limits_{j=1}^m P_2 \times Q_j' \right) \sqcup \left( \bigsqcup\limits_{i=1}^n P_i' \times Q_2 \right) \sqcup \left( \bigsqcup\limits_{i=1}^n \bigsqcup\limits_{j=1}^m P_i' \times Q_j' \right) $$
Это доказывает, что разность $A \setminus B$ представима как конечное дизъюнктное объединение элементов из $\mathcal{P} \odot \mathcal{Q}$.
\end{Proof}

\begin{center}
\begin{tikzpicture}[scale=1.4, >=latex]

    % --- Настройка цветов ---
    \definecolor{colorB}{HTML}{FADBD8}      % Бледно-красный (вычитаемое множество B)
    \definecolor{colorQ1}{HTML}{D4E6F1}     % Светло-синий (верхний левый кусок)
    \definecolor{colorQ2}{HTML}{A9CCE3}     % Синий (нижний правый кусок)
    \definecolor{colorQ3}{HTML}{7FB3D5}     % Темно-синий (верхний правый кусок)
    \definecolor{darktext}{HTML}{154360}    % Темно-синий для текста и линий

    % ==========================================
    % 1. Отрисовка цветных блоков
    % ==========================================

    % Множество B = P_2 x Q_2 (вычитаем его)
    \filldraw[fill=colorB, draw=darktext, thick] (0,0) rectangle (3,2);
    
    % Часть 1 разности: P_2 x (Q_1 \setminus Q_2)
    \filldraw[fill=colorQ1, draw=darktext, thick] (0,2) rectangle (3,4);
    
    % Часть 2 разности: (P_1 \setminus P_2) x Q_2
    \filldraw[fill=colorQ2, draw=darktext, thick] (3,0) rectangle (5,2);
    
    % Часть 3 разности: (P_1 \setminus P_2) x (Q_1 \setminus Q_2)
    \filldraw[fill=colorQ3, draw=darktext, thick] (3,2) rectangle (5,4);

    % ==========================================
    % 2. Подписи внутри прямоугольников
    % ==========================================
    
    \node[darktext, font=\Large\bfseries] at (1.5, 1) {$B$};
    \node[darktext, font=\small] at (1.5, 0.5) {$P_2 \times Q_2$};

    \node[darktext, font=\bfseries, align=center] at (1.5, 3) {Часть 1 \\ $P_2 \times Q'_1$};
    \node[darktext, font=\bfseries, align=center] at (4, 1) {Часть 2 \\ $P'_1 \times Q_2$};
    \node[darktext, font=\bfseries, align=center] at (4, 3) {Часть 3 \\ $P'_1 \times Q'_1$};

    % ==========================================
    % 3. Разметка осей X и Y (проекции множеств)
    % ==========================================

    % Ось X (множества P)
    \draw[|<->|, darktext, thick] (0, -0.4) -- (3, -0.4) node[midway, below] {$P_2$};
    \draw[|<->|, darktext, thick] (3, -0.4) -- (5, -0.4) node[midway, below] {$P'_1 = P_1 \setminus P_2$};
    \draw[|<->|, darktext, thick] (0, -1) -- (5, -1) node[midway, below] {$P_1$};

    % Ось Y (множества Q)
    \draw[|<->|, darktext, thick] (-0.4, 0) -- (-0.4, 2) node[midway, left] {$Q_2$};
    \draw[|<->|, darktext, thick] (-0.4, 2) -- (-0.4, 4) node[midway, left] {$Q'_1 = Q_1 \setminus Q_2$};
    \draw[|<->|, darktext, thick] (-1, 0) -- (-1, 4) node[midway, left] {$Q_1$};

    % ==========================================
    % 4. Общая граница большого прямоугольника A
    % ==========================================
    \draw[darktext, ultra thick] (0,0) rectangle (5,4);
    \node[darktext, font=\Large\bfseries, above right] at (0,4) {$A = P_1 \times Q_1$};

    % ==========================================
    % 5. Легенда снизу
    % ==========================================
    \node[darktext, align=left, anchor=north] at (2.5, -1.5) {
        Разность $A \setminus B$ представляет собой «уголок» (синяя зона). \\
        Мы разбиваем проекции на попарно непересекающиеся части: \\
        $P_1 = P_2 \sqcup P'_1$ \quad и \quad $Q_1 = Q_2 \sqcup Q'_1$ \\
        В результате разность $A \setminus B$ разбивается на 3 прямоугольника, \\
        каждый из которых принадлежит полукольцу $\mathcal{P} \odot \mathcal{Q}$.
    };

\end{tikzpicture}
\end{center}

\begin{Example}
Множество $\mathcal{P}^n = \{[a, b)\} $ является полукольцом, где $a = (a_1, \dots, a_n)$, $b = (b_1, \dots, b_n)$ и $a_i \le b_i$.
Многомерный полуинтервал определяется как:
$$ [a, b) = \{ x \in \mathbb{R}^n \mid a_i \le x_i < b_i \} $$
Это есть декартово произведение $n$ одномерных полуколец.
\end{Example}

\begin{Definition}{Произведение объёмов}{}
Пусть $\mathcal{P}$ и $\mathcal{Q}$ — полукольца, $\varphi$ — объём на $\mathcal{P}$, $\psi$ — объём на $\mathcal{Q}$.
Тогда функция $\nu(P \times Q) = \varphi(P) \cdot \psi(Q)$ называется \textbf{произведением объёмов} (считается, что $0 \cdot (+\infty) = 0$ и $(+\infty) \cdot 0 = 0$).
\end{Definition}

\begin{Theorem}{Об объёме на произведении}{}
Определенная выше функция $\nu$ является объёмом на $\mathcal{P} \odot \mathcal{Q}$.
\end{Theorem}

\begin{Proof}
Покажем, что функция $\nu(P \times Q) = \varphi(P) \cdot \psi(Q)$ является объёмом на $\mathcal{P} \odot \mathcal{Q}$.

\textbf{1) Свойство нуля:} \\
Так как $\emptyset \in \mathcal{P}$ и $\emptyset \in \mathcal{Q}$, пустое множество в произведении полуколец можно представить как $\emptyset \times \emptyset$.
Тогда по определению:
$$ \nu(\emptyset) = \nu(\emptyset \times \emptyset) = \varphi(\emptyset) \cdot \psi(\emptyset) = 0 \cdot 0 = 0 $$
(Аналогично, если хотя бы один из множителей $P$ или $Q$ является пустым, произведение даст $0$).

\textbf{2) Конечная аддитивность:} \\
Пусть прямоугольник $P \times Q \in \mathcal{P} \odot \mathcal{Q}$ разбит на конечное дизъюнктное объединение прямоугольников из того же полукольца:
$$ P \times Q = \bigsqcup\limits_{k=1}^n (P_k \times Q_k) $$
где $P, P_k \in \mathcal{P}$ и $Q, Q_k \in \mathcal{Q}$.
Воспользуемся свойством полукольца об упрощении структуры объединения. Спроецируем все стороны прямоугольников на оси. Существует такое разбиение $P$ на дизъюнктные множества $\tilde{P}_i \in \mathcal{P}$ и $Q$ на множества $\tilde{Q}_j \in \mathcal{Q}$, что:
$$ P = \bigsqcup\limits_{i=1}^{I} \tilde{P}_i \quad \text{и} \quad Q = \bigsqcup\limits_{j=1}^{J} \tilde{Q}_j $$
При этом каждое $P_k$ является точным объединением какой-то части множеств $\tilde{P}_i$, а каждое $Q_k$ — части множеств $\tilde{Q}_j$.
В силу конечной аддитивности самих объёмов $\varphi$ и $\psi$:
$$ \nu(P \times Q) = \varphi(P)\psi(Q) = \left( \sum\limits_{i=1}^{I} \varphi(\tilde{P}_i) \right) \left( \sum\limits_{j=1}^{J} \psi(\tilde{Q}_j) \right) = \sum\limits_{i=1}^{I} \sum\limits_{j=1}^{J} \varphi(\tilde{P}_i)\psi(\tilde{Q}_j) $$
С другой стороны, каждый из меньших прямоугольников $P_k \times Q_k$ состоит в точности из тех «клеток» $\tilde{P}_i \times \tilde{Q}_j$, которые в него попадают. Запишем сумму их объёмов:
$$ \sum\limits_{k=1}^n \nu(P_k \times Q_k) = \sum\limits_{k=1}^n \varphi(P_k)\psi(Q_k) = \sum\limits_{k=1}^n \sum\limits_{\text{клетки } (i, j) \in P_k \times Q_k} \varphi(\tilde{P}_i)\psi(\tilde{Q}_j) $$
Поскольку исходное объединение прямоугольников было дизъюнктным и в точности покрывало всё множество $P \times Q$, каждая пара $(i, j)$ учитывается в правой сумме ровно один раз. Значит, суммы абсолютно идентичны. Аддитивность доказана.
\end{Proof}

\section{§ Мера и её свойства}

\begin{Definition}{Мера}{}
Объём $\varphi: \mathcal{P} \to [0, +\infty]$ называется \textbf{мерой}, если для любой последовательности непересекающихся множеств $P, P_1, P_2, \dots \in \mathcal{P}$ выполняется свойство \textbf{счётной аддитивности}:
$$ P = \bigsqcup\limits_{i=1}^{\infty} P_i \implies \varphi(P) = \sum\limits_{i=1}^{\infty} \varphi(P_i) $$
\end{Definition}

\begin{Example}
\begin{enumerate}
    \item Мера Лебега на полукольце полуинтервалов в $\mathbb{R}$: 
    $$ \lambda([a, b)) = b - a $$
    Она же в многомерном случае $\mathbb{R}^n$: 
    $$ \lambda([a, b)) = \prod\limits_{i=1}^n (b_i - a_i) $$
    
    \item Пусть $f$ — возрастающая функция. Зададим объём (меру Лебега-Стилтьеса):
    $$ \nu([a, b)) = f(b) - f(a) $$
    Представим полуинтервал в виде счётного объединения:
    $$ [a, b) = \bigcup\limits_{i=1}^{\infty} [a_{i-1}, a_i) $$
    где $a_0 = a$, $a_i < a_{i+1}$ и $\lim\limits_{n \to \infty} a_n = b$. 
    Если счётная аддитивность выполняется, то:
    $$ \nu([a, b)) = \sum\limits_{i=1}^{\infty} \nu([a_{i-1}, a_i)) = \sum\limits_{i=1}^{\infty} \left( f(a_i) - f(a_{i-1}) \right) = \lim\limits_{n \to \infty} (f(a_n) - f(a_0)) $$
    Отсюда следует, что $\lim\limits_{n \to \infty} f(a_n) = f(b)$, то есть функция $f$ должна быть \textbf{непрерывна слева}.
    
    \item Пусть $\Omega$ — не более чем счётное пространство элементарных исходов, $\mathcal{U} = 2^\Omega$. Каждому $\omega_i \leftrightarrow p_i \ge 0$.
    Для любого $A \in \mathcal{U}$ задан объём:
    $$ \varphi(A) = \sum\limits_{i: \omega_i \in A} p_i $$
    Покажем, что это мера. Пусть $A = \bigsqcup\limits_{i=1}^{\infty} A_i$. Нам нужно доказать, что $\varphi(A) = \sum\limits_{i=1}^{\infty} \varphi(A_i)$.
    Для любого конечного $N \in \mathbb{N}$ верно, что $\bigsqcup\limits_{i=1}^N A_i \subset A$. В силу усиленной монотонности объёма:
    $$ \sum\limits_{i=1}^N \varphi(A_i) \le \varphi(A) $$
    Переходя к пределу при $N \to \infty$, получаем $\sum\limits_{i=1}^{\infty} \varphi(A_i) \le \varphi(A)$. \\
    С другой стороны, из свойства полуаддитивности: 
    $$ \varphi(A) \le \sum\limits_{i=1}^N \varphi(A_i) \le \sum\limits_{i=1}^{\infty} \varphi(A_i) $$
    Следовательно, $\varphi(A) = \sum\limits_{i=1}^{\infty} \varphi(A_i)$, и объём $\varphi$ является мерой.
\end{enumerate}
\end{Example}

\begin{Theorem}{Критерий меры}{}
Пусть задан объём $\varphi: \mathcal{P} \to [0, +\infty]$. Для того чтобы $\varphi$ был мерой, необходимо и достаточно, чтобы выполнялось свойство \textbf{счётной полуаддитивности}:
$$ P, P_1, P_2, \dots \in \mathcal{P} \quad \text{и} \quad P \subset \bigcup\limits_{i=1}^{\infty} P_i \implies \varphi(P) \le \sum\limits_{i=1}^{\infty} \varphi(P_i) $$
\end{Theorem}

\begin{Proof}
$\Rightarrow$ Пусть $\varphi$ — мера. Обозначим $P_i' = P \cap P_i \in \mathcal{P}$. Тогда $P = \bigcup\limits_{i=1}^{\infty} P_i'$. 
По свойству полукольца (лемма про нарезку), это объединение можно представить как дизъюнктное:
$$ \bigcup\limits_{i=1}^{\infty} P_i' = \bigsqcup\limits_{i=1}^{\infty} \bigsqcup\limits_{j=1}^{m_i} Q_{ij} $$
где $Q_{ij} \in \mathcal{P}$ и $Q_{ij} \subset P_i' \subset P_i$. Из усиленной монотонности для каждого $i$ следует:
$$ \sum\limits_{j=1}^{m_i} \varphi(Q_{ij}) \le \varphi(P_i') \le \varphi(P_i) $$
Так как $\varphi$ мера (счётно аддитивна на дизъюнктных объединениях), то:
$$ \varphi(P) = \sum\limits_{i=1}^{\infty} \sum\limits_{j=1}^{m_i} \varphi(Q_{ij}) \le \sum\limits_{i=1}^{\infty} \varphi(P_i) $$

$\Leftarrow$ Пусть выполняется счётная полуаддитивность. Рассмотрим $P = \bigsqcup\limits_{i=1}^{\infty} P_i$, где $P, P_i \in \mathcal{P}$. 
Тогда из счётной полуаддитивности имеем: $\varphi(P) \le \sum\limits_{i=1}^{\infty} \varphi(P_i)$. \\
С другой стороны, для любого $N$ верно $\bigsqcup\limits_{i=1}^N P_i \subset P$. Из усиленной монотонности объёма:
$$ \sum\limits_{i=1}^N \varphi(P_i) \le \varphi(P) \implies \lim\limits_{N \to \infty} \sum\limits_{i=1}^N \varphi(P_i) = \sum\limits_{i=1}^{\infty} \varphi(P_i) \le \varphi(P) $$
Объединяя неравенства $\le$ и $\ge$, получаем равенство $\varphi(P) = \sum\limits_{i=1}^{\infty} \varphi(P_i)$.
\end{Proof}

\begin{Theorem}{Непрерывность меры снизу}{}
Пусть $\varphi$ — объём, заданный на алгебре $\mathcal{U}$. Объём $\varphi$ является мерой тогда и только тогда, когда $\varphi$ \textbf{непрерывен снизу}, то есть:
$$ A_i \subset A_{i+1}, \quad A = \bigcup\limits_{i=1}^{\infty} A_i \quad (A_i, A \in \mathcal{U}) \implies \lim\limits_{i \to \infty} \varphi(A_i) = \varphi(A) $$
\end{Theorem}

\begin{Proof}
\textbf{Лемма (вспомогательная):} Если $A_{i-1} \subset A_i$, то $A_i = A_{i-1} \sqcup (A_i \setminus A_{i-1})$. По аддитивности объёма $\varphi(A_i \setminus A_{i-1}) = \varphi(A_i) - \varphi(A_{i-1})$.

$\Rightarrow$ Пусть $\varphi$ — мера. Представим $A$ в виде дизъюнктного объединения:
$$ A = \bigsqcup\limits_{i=1}^{\infty} (A_i \setminus A_{i-1}), \quad \text{где } A_0 = \emptyset $$
Тогда:
$$ \varphi(A) = \sum\limits_{i=1}^{\infty} \varphi(A_i \setminus A_{i-1}) = \sum\limits_{i=1}^{\infty} (\varphi(A_i) - \varphi(A_{i-1})) = \lim\limits_{N \to \infty} (\varphi(A_N) - \varphi(A_0)) = \lim\limits_{N \to \infty} \varphi(A_N) $$

$\Leftarrow$ Пусть $A = \bigsqcup\limits_{i=1}^{\infty} A_i$, где $A, A_i \in \mathcal{U}$. Обозначим $B_N = \bigsqcup\limits_{i=1}^N A_i$.
Очевидно, что $B_{N+1} \supset B_N$ и $A = \bigcup\limits_{N=1}^{\infty} B_N$. В силу непрерывности снизу:
$$ \lim\limits_{N \to \infty} \varphi(B_N) = \varphi(A) $$
Так как $\varphi(B_N) = \sum\limits_{i=1}^N \varphi(A_i)$, то $\lim\limits_{N \to \infty} \sum\limits_{i=1}^N \varphi(A_i) = \sum\limits_{i=1}^{\infty} \varphi(A_i) = \varphi(A)$. Счётная аддитивность доказана.
\end{Proof}

\begin{center}
\begin{tikzpicture}[scale=1]
    % Самый большой овал (A3)
    \draw[thick, defframe, fill=defbg!30] (0,0) ellipse (4cm and 2.5cm);
    
    % Средний овал (A2)
    \draw[thick, defframe, fill=defbg!60] (-0.5,-0.2) ellipse (3cm and 1.8cm);
    
    % Маленький овал (A1)
    \draw[thick, defframe, fill=defbg!100] (-1,-0.4) ellipse (1.5cm and 1cm);
    
    % Подписи
    \node at (-1,-0.4) {$A_1$};
    \node at (1.2, 0.5) {$A_2 \setminus A_1$};
    \node at (2.5, 1.2) {$A_3 \setminus A_2$};
    \node[above] at (0, 2.6) {Обобщение: $A = \bigsqcup\limits_{i=1}^{\infty} (A_i \setminus A_{i-1})$};
\end{tikzpicture}
\end{center}

\begin{Theorem}{Непрерывность меры сверху}{}
Пусть $\varphi$ — \textbf{конечный} объём, заданный на алгебре $\mathcal{U}$. Тогда следующие утверждения эквивалентны:
\begin{enumerate}
    \item $\varphi$ — мера.
    \item $\varphi$ \textbf{непрерывна сверху}, то есть: $A_{i+1} \subset A_i$, $A = \bigcap\limits_{i=1}^{\infty} A_i \quad (A, A_i \in \mathcal{U}) \implies \varphi(A) = \lim\limits_{N \to \infty} \varphi(A_N)$.
    \item $\varphi$ непрерывна сверху на пустом множестве (то есть при $A = \emptyset$).
\end{enumerate}
\end{Theorem}

\begin{Proof}
\textbf{1 $\Rightarrow$ 2:} Пусть $\varphi$ — мера. Представим $A_1$ в виде дизъюнктного объединения:
$$ A_1 = A \sqcup \bigsqcup\limits_{i=1}^{\infty} (A_i \setminus A_{i+1}) $$
По свойству меры:
$$ \varphi(A_1) = \varphi(A) + \sum\limits_{i=1}^{\infty} (\varphi(A_i) - \varphi(A_{i+1})) = \varphi(A) + \lim\limits_{N \to \infty} \sum\limits_{i=1}^N (\varphi(A_i) - \varphi(A_{i+1})) $$
Эта сумма телескопическая, поэтому:
$$ \varphi(A_1) = \varphi(A) + \lim\limits_{N \to \infty} (\varphi(A_1) - \varphi(A_{N+1})) = \varphi(A) + \varphi(A_1) - \lim\limits_{N \to \infty} \varphi(A_{N+1}) $$
Так как мера \textbf{конечна} (т.е. $\varphi(A_1) < \infty$), мы можем вычесть $\varphi(A_1)$ из обеих частей, получая $\varphi(A) = \lim\limits_{N \to \infty} \varphi(A_{N+1})$.

\textbf{2 $\Rightarrow$ 3:} Очевидно (достаточно подставить $A = \emptyset$, тогда $\varphi(\emptyset) = 0$).

\textbf{3 $\Rightarrow$ 1:} Пусть $A = \bigsqcup\limits_{i=1}^{\infty} A_i$, где $A, A_i \in \mathcal{U}$. Рассмотрим остаток множества $B_n$:
$$ B_n = A \setminus \bigsqcup\limits_{i=1}^n A_i $$
Множества $B_n$ убывают ($B_{n+1} \subset B_n$) и $\bigcap\limits_{n=1}^{\infty} B_n = \emptyset$. По условию 3, $\lim\limits_{n \to \infty} \varphi(B_n) = \varphi(\emptyset) = 0$.
С другой стороны:
$$ 0 = \lim\limits_{n \to \infty} \varphi(B_n) = \varphi(A) - \lim\limits_{n \to \infty} \varphi \left( \bigsqcup\limits_{i=1}^n A_i \right) = \varphi(A) - \lim\limits_{n \to \infty} \sum\limits_{i=1}^n \varphi(A_i) = \varphi(A) - \sum\limits_{i=1}^{\infty} \varphi(A_i) $$
Отсюда $\varphi(A) = \sum\limits_{i=1}^{\infty} \varphi(A_i)$, что означает, что $\varphi$ — мера.
\end{Proof}


\setcounter{section}{4}

\section{§ Внешняя мера и её свойства}

Глобальная идея: погрузим множество в не более чем счётное семейство множеств, которое можно измерить. 
$$ \text{Мера на скудной системе} \xrightarrow{\text{строим}} \text{Внешнюю меру} \xrightarrow{\text{сужаем}} \text{Мера на большей системе множеств} $$

\begin{Definition}{Внешняя мера}{}
Пусть $X$ — множество. Функция $\tau: 2^X \to [0, +\infty]$ называется \textbf{внешней мерой}, если:
\begin{enumerate}
    \item $\tau(\emptyset) = 0$.
    \item $\tau$ счётно-полуаддитивна: 
    $$ A \subset \bigcup\limits_{i=1}^{\infty} A_i \implies \tau(A) \le \sum\limits_{i=1}^{\infty} \tau(A_i) $$
\end{enumerate}
\end{Definition}

\noindent \textbf{Свойства внешней меры:}
\begin{enumerate}
    \item Внешняя мера \textbf{конечно-полуаддитивна}: $ A \subset \bigcup\limits_{i=1}^n A_i \implies \tau(A) \le \sum\limits_{i=1}^n \tau(A_i) $
    \item Внешняя мера \textbf{монотонна}: $ A \subset B \implies \tau(A) \le \tau(B) $
\end{enumerate}

\begin{Proof}[Доказательство свойств]
(Из 1 следует 2). Пусть $A \subset B$. Множество $A$ можно покрыть следующим образом: $A \subset B \cup \bigcup\limits_{k=1}^{\infty} \emptyset$. Тогда из счётной полуаддитивности:
$$ \tau(A) \le \tau(B) + \sum\limits_{k=1}^{\infty} \tau(\emptyset) = \tau(B) + 0 = \tau(B) $$
\end{Proof}

\begin{Definition}{$\tau$-измеримое множество}{}
Множество $A$ называется \textbf{$\tau$-измеримым}, если для любого множества $E \subset X$ выполняется равенство:
$$ \tau(E) = \tau(E \cap A) + \tau(E \setminus A) $$
(Интуитивно: множество $A$ разбивает произвольное множество $E$ аддитивно).
\end{Definition}

\begin{center}
\begin{tikzpicture}[scale=0.9, every node/.style={scale=0.9}]
    % Множество A (правая полуплоскость)
    \fill[defbg!80!white] (1.5, -2) rectangle (5, 3);
    \draw[thick, defframe] (1.5, -2) -- (1.5, 3) node[above right, text=black] {$A$ (измеримое множество)};
    
    % Множество E (овал)
    \draw[thick, fill=lembg, opacity=0.9] (0,0) ellipse (2.5cm and 1.5cm);
    \node[left] at (-2.6, 0) {$E$};
    
    % Пересечение E \cap A
    \begin{scope}
        \clip (1.5, -2) rectangle (5, 3);
        \fill[lemframe!40!white] (0,0) ellipse (2.5cm and 1.5cm);
    \end{scope}
    
    % Обводка E поверх всего, чтобы граница была четкой
    \draw[thick, lemframe] (0,0) ellipse (2.5cm and 1.5cm);
    \draw[thick, dashed] (1.5, -1.2) -- (1.5, 1.2); % Линия разреза внутри E
    
    % Подписи
    \node at (0.2, 0) {$E \setminus A$};
    \node at (2.6, 0) {$E \cap A$};
\end{tikzpicture}
\end{center}

\noindent \textbf{NB!} Автоматически всегда имеем неравенство $\tau(E) \le \tau(E \cap A) + \tau(E \setminus A)$. \\
\textit{Доказательство:} Множество $E$ можно представить как $E = (E \cap A) \sqcup (E \setminus A)$. Значит, $E \subset (E \cap A) \cup (E \setminus A)$. Так как внешняя мера конечно-полуаддитивна, то $\tau(E) \le \tau(E \cap A) + \tau(E \setminus A)$.

Следовательно, для проверки $\tau$-измеримости достаточно доказать обратное неравенство: 
$$ \tau(E) \ge \tau(E \cap A) + \tau(E \setminus A) $$

\noindent \textbf{Замечание:} Пустое множество $\emptyset$ является $\tau$-измеримым. \\
Действительно, $E \cap \emptyset = \emptyset \subset A \implies \tau(E \cap \emptyset) = 0$. Проверяем неравенство:
$$ \tau(E) \ge \tau(E \setminus \emptyset) = \tau(E) = \tau(E \cap \emptyset) + \tau(E \setminus \emptyset) \implies \emptyset \text{ — } \tau\text{-измеримо.} $$

\begin{Theorem}{О $\sigma$-алгебре измеримых множеств}{}
Пусть $\tau$ — внешняя мера. Система $\mathcal{U}_\tau$ всех $\tau$-измеримых множеств является $\sigma$-алгеброй, а сужение функции $\tau$ на $\mathcal{U}_\tau$ является мерой.
\end{Theorem}

\begin{Proof}
\textbf{Пункт 1: Покажем, что $\mathcal{U}_\tau$ — алгебра.} \\
Мы уже показали, что $\emptyset \in \mathcal{U}_\tau$. \\
Пусть $\exists A \in \mathcal{U}_\tau \implies \forall E \in 2^X: \tau(E) = \tau(E \setminus A) + \tau(E \cap A)$. \\
Заметим, что $E \cap A = E \setminus A^c$ и $E \setminus A = E \cap A^c$. Подставляя, получаем:
$$ \tau(E) = \tau(E \setminus A^c) + \tau(E \cap A^c) \implies A^c \in \mathcal{U}_\tau $$
Пусть $A, B \in \mathcal{U}_\tau$. Проверим $A \cup B$:
\begin{align*}
    \tau(E) &= \tau(E \cap A) + \tau(E \setminus A) \\
    &= \tau(E \cap A) + \tau((E \setminus A) \cap B) + \tau((E \setminus A) \setminus B) \\
    &\ge \tau((E \cap A) \cup ((E \setminus A) \cap B)) + \tau(E \setminus (A \cup B)) \\
    &= \tau(E \cap (A \cup B)) + \tau(E \setminus (A \cup B))
\end{align*}
Получили неравенство $\tau(E) \ge \tau(E \cap (A \cup B)) + \tau(E \setminus (A \cup B))$, что влечёт $A \cup B \in \mathcal{U}_\tau$.

\textbf{Пункт 2: Покажем, что $\tau$ — объём на $\mathcal{U}_\tau$.} \\
Пусть $\exists A, B \in \mathcal{U}_\tau$ и $A \cap B = \emptyset$. Возьмём в качестве $E = A \sqcup B$:
$$ \tau(A \sqcup B) = \tau((A \sqcup B) \cap A) + \tau((A \sqcup B) \setminus A) = \tau(A) + \tau(B) $$
Следовательно, есть конечная аддитивность.

\textbf{Пункт 3: Покажем, что $\mathcal{U}_\tau$ — $\sigma$-алгебра.} \\
Пусть $A = \bigsqcup\limits_{i=1}^{\infty} A_i$, где $A_i \in \mathcal{U}_\tau$. \\
Для любого $N \in \mathbb{N}$ по индукции (используя измеримость):
$$ \tau(E) = \tau \left( E \cap \bigsqcup\limits_{i=1}^N A_i \right) + \tau \left( E \setminus \bigsqcup\limits_{i=1}^N A_i \right) = \sum\limits_{i=1}^N \tau(E \cap A_i) + \tau \left( E \setminus \bigsqcup\limits_{i=1}^N A_i \right) $$
Так как $\bigsqcup\limits_{i=1}^N A_i \subset A$, то $E \setminus A \subset E \setminus \bigsqcup\limits_{i=1}^N A_i$. В силу монотонности:
$$ \tau(E) \ge \sum\limits_{i=1}^N \tau(E \cap A_i) + \tau(E \setminus A) $$
Перейдём к пределу при $N \to \infty$:
$$ \tau(E) \ge \lim\limits_{N \to \infty} \sum\limits_{i=1}^N \tau(E \cap A_i) + \tau(E \setminus A) = \sum\limits_{i=1}^{\infty} \tau(E \cap A_i) + \tau(E \setminus A) $$
Используя счётную полуаддитивность внешней меры:
$$ \sum\limits_{i=1}^{\infty} \tau(E \cap A_i) \ge \tau \left( \bigcup\limits_{i=1}^{\infty} (E \cap A_i) \right) = \tau(E \cap A) $$
Следовательно, $\tau(E) \ge \tau(E \cap A) + \tau(E \setminus A) \implies A \in \mathcal{U}_\tau$. \\
Любое объединение $A = \bigsqcup\limits_{i=1}^{\infty} A_i$ можно записать как дизъюнктное: $\bigsqcup\limits_{i=1}^{\infty} \left( A_i \setminus \bigcup\limits_{j=0}^{i-1} A_j \right)$. Значит, $\mathcal{U}_\tau$ — $\sigma$-алгебра.

\textbf{Пункт 4.} По критерию меры $\tau$ является мерой, так как это объём + есть счётная полуаддитивность (по определению внешней меры).
\end{Proof}

\begin{Definition}{Полная мера}{}
Мера $\mu$ на полукольце $\mathcal{P}$ называется \textbf{полной}, если из того, что $A \in \mathcal{P}$ и $\mu(A) = 0$, следует, что для любого $B \subset A$:
\begin{enumerate}
    \item $B \in \mathcal{P}$
    \item $\mu(B) = 0$
\end{enumerate}
\end{Definition}

\noindent \textbf{Мотивирующий пример (Мера Лебега):} Пусть $\mathcal{U}$ — система подмножеств $\mathbb{R}$, содержащая одноточечные множества, и $\mu(\{\alpha\}) = 0$. Пусть $A \notin \mathcal{U}$. Рассмотрим прямое произведение $B = B_1 \times B_2$ и меру $\mu_2(B_1 \times B_2) = \mu_1(B_1)\mu_1(B_2)$.
Тогда $\mu_2(\{\alpha\} \times \mathbb{R}) = \mu_1(\{\alpha\}) \cdot \mu_1(\mathbb{R}) = 0 \cdot (+\infty) = 0$.
Однако множество $\{\alpha\} \times A \subset \{\alpha\} \times \mathbb{R}$, но при этом не факт, что оно измеримо, если мера не полная. В полной мере такое подмножество автоматически измеримо и имеет меру 0.

\begin{Lemma}{О полноте меры на $\mathcal{U}_\tau$}{}
Сужение $\tau$ на $\mathcal{U}_\tau$ является полной мерой.
\end{Lemma}

\begin{Proof}
Пусть $A \in \mathcal{U}_\tau$ и $\mu(A) = \tau(A) = 0$. Пусть $\forall B \subset A$. \\
В силу монотонности внешней меры: $\tau(B) \le \tau(A) = 0 \implies \mu(B) = 0$. \\
Покажем, что $B \in \mathcal{U}_\tau$ (т.е. $\tau$-измеримое). Проверим неравенство:
$$ \tau(E \cap B) + \tau(E \setminus B) \le \tau(E) $$
Заметим, что $\tau(E \cap B) \le \tau(B) = 0 \implies \tau(E \cap B) = 0$. \\
Также $E \setminus B \subset E$, поэтому $\tau(E \setminus B) \le \tau(E)$. \\
Итого: $0 + \tau(E \setminus B) \le \tau(E)$. Неравенство доказано, $B \in \mathcal{U}_\tau$.
\end{Proof}

\begin{Example}
Пусть $X = \mathbb{R}$. Зададим $\tau$: $\tau(A) = 1$ если $A \neq \emptyset$, и $\tau(\emptyset) = 0$. Это внешняя мера. \\
\textbf{Факт:} $\tau$-измеримы только $\emptyset$ и $\mathbb{R}$. То есть $\mathcal{U}_\tau = \{\emptyset, \mathbb{R}\}$. \\
\textit{Доказательство:} Пусть $\exists A \neq \emptyset$ и $A \neq \mathbb{R}$. Для любого $E$ должно быть $\tau(E) = \tau(E \cap A) + \tau(E \setminus A)$. Возьмём $E = \mathbb{R}$.
Тогда $1 = \tau(\mathbb{R} \cap A) + \tau(\mathbb{R} \setminus A) \implies 1 = 1 + 1 \implies 1 = 2$. Противоречие. Значит, $A \notin \mathcal{U}_\tau$.
\end{Example}

\section{§ Процесс Каратеодори (Продолжение меры)}

\begin{Definition}{Внешняя мера, порождённая мерой}{}
Пусть $\mu_0$ — мера на полукольце $\mathcal{P}$. \textbf{Внешней мерой, порождённой $\mu_0$}, назовём функцию:
$$ \mu^*(A) = \inf \left\{ \sum\limits_{i=1}^{\infty} \mu_0(P_i) : A \subset \bigcup\limits_{i=1}^{\infty} P_i, \quad P_i \in \mathcal{P} \right\} $$
(Полная аналогия с суммами Дарбу для интеграла Римана). Если множество $A$ не может быть покрыто элементами полукольца, то полагаем $\mu^*(A) = +\infty$.
\end{Definition}

\begin{Theorem}{О свойствах порождённой внешней меры}{}
Функция $\mu^*$ является внешней мерой, совпадающей с $\mu_0$ на полукольце $\mathcal{P}$.
\end{Theorem}

\begin{Proof}
\textbf{1) Покажем, что $\forall P \in \mathcal{P}, \;\; \mu^*(P) = \mu_0(P)$.} \\
Так как $P \subset P \cup \emptyset \cup \emptyset \dots$, очевидно, что $\mu^*(P) \le \mu_0(P)$. \\
С другой стороны, пусть $P \subset \bigcup\limits_{i=1}^{\infty} P_i$. Так как $\mu_0$ — мера (а значит, обладает свойством счётной полуаддитивности на полукольце):
$$ \mu_0(P) \le \sum\limits_{i=1}^{\infty} \mu_0(P_i) $$
Так как это верно для любого покрытия, переходя к точной нижней грани ($\inf$) по всем покрытиям, получим $\mu_0(P) \le \mu^*(P)$. Следовательно, $\mu^*(P) = \mu_0(P)$.

\textbf{2) Покажем, что $\mu^*$ — внешняя мера.} \\
$\mu^* \ge 0$ в силу определения. \\
$\mu^*(\emptyset) = \mu_0(\emptyset) = 0$ (так как $\mu_0$ — мера). \\
Осталось доказать счётную полуаддитивность: $A \subset \bigcup\limits_{i=1}^{\infty} A_i \implies \mu^*(A) \le \sum\limits_{i=1}^{\infty} \mu^*(A_i)$.
Будем считать, что $\forall i \in \mathbb{N}: \mu^*(A_i) < +\infty$, иначе неравенство выполняется тривиально.

Зафиксируем произвольное $\varepsilon > 0$. По определению точной нижней грани, для каждого $A_i$ существует покрытие элементами полукольца $P_i^n \in \mathcal{P}$ такое, что $A_i \subset \bigcup\limits_{n=1}^{\infty} P_i^n$ и выполняется неравенство:
$$ \sum\limits_{n=1}^{\infty} \mu_0(P_i^n) \le \mu^*(A_i) + \frac{\varepsilon}{2^i} $$

Множество $A$ содержится в объединении всех $A_i$, следовательно:
$$ A \subset \bigcup\limits_{i=1}^{\infty} A_i \subset \bigcup\limits_{i=1}^{\infty} \bigcup\limits_{n=1}^{\infty} P_i^n $$
Это счётное объединение элементов из полукольца $\mathcal{P}$. По определению внешней меры:
$$ \mu^*(A) \le \sum\limits_{i=1}^{\infty} \sum\limits_{n=1}^{\infty} \mu_0(P_i^n) \le \sum\limits_{i=1}^{\infty} \left( \mu^*(A_i) + \frac{\varepsilon}{2^i} \right) = \sum\limits_{i=1}^{\infty} \mu^*(A_i) + \varepsilon \sum\limits_{i=1}^{\infty} \frac{1}{2^i} = \sum\limits_{i=1}^{\infty} \mu^*(A_i) + \varepsilon $$
Так как $\varepsilon > 0$ было выбрано произвольно, то $\mu^*(A) \le \sum\limits_{i=1}^{\infty} \mu^*(A_i)$. Счётная полуаддитивность доказана.
\end{Proof}

\begin{Theorem}{О продолжении меры (Теорема Каратеодори)}{}
Сужение внешней меры $\mu^*$ на систему $\mathcal{U}_{\mu^*}$ всех $\mu^*$-измеримых множеств является мерой. Это сужение является \textbf{продолжением} меры $\mu_0$ с полукольца $\mathcal{P}$. В частности, $\mathcal{P} \subset \mathcal{U}_{\mu^*}$.
\end{Theorem}

\begin{Proof}
Из предыдущих теорем мы уже знаем, что $\mu^*$ — внешняя мера, а её сужение на систему измеримых множеств $\mathcal{U}_{\mu^*}$ является мерой. Осталось показать, что элементы полукольца $\mathcal{P}$ измеримы, то есть $\mathcal{P} \subset \mathcal{U}_{\mu^*}$.

Пусть $P \in \mathcal{P}$. Нам нужно доказать, что для любого множества $E \subset X$ выполняется:
$$ \mu^*(E) \ge \mu^*(E \cap P) + \mu^*(E \setminus P) $$

\textbf{Случай а):} $E \in \mathcal{P}$ (множество $E$ из полукольца). \\
Тогда разность $E \setminus P$ можно представить как конечное дизъюнктное объединение: $E \setminus P = \bigsqcup\limits_{i=1}^n Q_i$, где $Q_i \in \mathcal{P}$. 
Следовательно, $E = (E \cap P) \sqcup \bigsqcup\limits_{i=1}^n Q_i$.
Так как на элементах полукольца $\mu^*$ совпадает с $\mu_0$, а $\mu_0$ конечно аддитивна:
$$ \mu^*(E) = \mu_0(E) = \mu_0(E \cap P) + \sum\limits_{i=1}^n \mu_0(Q_i) = \mu^*(E \cap P) + \sum\limits_{i=1}^n \mu^*(Q_i) $$
Применяя конечную полуаддитивность внешней меры:
$$ \mu^*(E) \ge \mu^*(E \cap P) + \mu^* \left( \bigcup\limits_{i=1}^n Q_i \right) = \mu^*(E \cap P) + \mu^*(E \setminus P) $$

\textbf{Случай б):} $E \notin \mathcal{P}$ (множество $E$ не из полукольца). \\
Если $\mu^*(E) = +\infty$, то неравенство выполняется тривиально. 
Если $\mu^*(E) < +\infty$, то для любого $\varepsilon > 0$ существует покрытие $E \subset \bigcup\limits_{i=1}^{\infty} P_i$ ($P_i \in \mathcal{P}$), такое что:
$$ \mu^*(E) + \varepsilon > \sum\limits_{i=1}^{\infty} \mu_0(P_i) $$
Так как $P_i \in \mathcal{P}$, для каждого из них выполняется доказанный \textbf{Случай а} (подставляя $P_i$ вместо $E$):
$$ \sum\limits_{i=1}^{\infty} \mu_0(P_i) = \sum\limits_{i=1}^{\infty} \mu^*(P_i) \ge \sum\limits_{i=1}^{\infty} \left( \mu^*(P_i \cap P) + \mu^*(P_i \setminus P) \right) $$
Используя счётную полуаддитивность внешней меры:
$$ \sum\limits_{i=1}^{\infty} \mu^*(P_i \cap P) + \sum\limits_{i=1}^{\infty} \mu^*(P_i \setminus P) \ge \mu^* \left( \bigcup\limits_{i=1}^{\infty} (P_i \cap P) \right) + \mu^* \left( \bigcup\limits_{i=1}^{\infty} (P_i \setminus P) \right) $$
Заметим, что $E \cap P \subset \bigcup\limits_{i=1}^{\infty} (P_i \cap P)$ и $E \setminus P \subset \bigcup\limits_{i=1}^{\infty} (P_i \setminus P)$. В силу монотонности:
$$ \mu^*(E) + \varepsilon > \mu^*(E \cap P) + \mu^*(E \setminus P) $$
Устремляя $\varepsilon \to 0$, получаем требуемое неравенство. Следовательно, $P \in \mathcal{U}_{\mu^*}$.
\end{Proof}

\begin{Definition}{Пространство с мерой}{}
Тройка $(X, \mathcal{U}, \mu)$, где $X$ — базовое множество, $\mathcal{U}$ — $\sigma$-алгебра измеримых множеств, а $\mu$ — мера, заданная на $\mathcal{U}$, называется \textbf{пространством с мерой} (или измеримым пространством).
\end{Definition}

\begin{Example}
Пусть $X = \{1, 2, 3\}$. Дано полукольцо $\mathcal{P} = \{\emptyset, \{2, 3\}\}$. Задана мера $\mu_0(\emptyset) = 0, \mu_0(\{2, 3\}) = 0$.
Выполним процесс Каратеодори для построения $\mu^*$:
\begin{itemize}
    \item $\mu^*(\{2\}) = 0$ и $\mu^*(\{3\}) = 0$, так как они покрываются множеством $\{2, 3\}$ (монотонность).
    \item Множества, содержащие элемент $1$, невозможно покрыть элементами из $\mathcal{P}$, поэтому их внешняя мера равна $+\infty$:
    $$ \mu^*(\{1\}) = \mu^*(\{1, 2\}) = \mu^*(\{1, 3\}) = \mu^*(\{1, 2, 3\}) = +\infty $$
\end{itemize}
В данном случае множество $X$ нельзя покрыть элементами конечной меры, поэтому эта мера \textbf{не является $\sigma$-конечной}. Система измеримых множеств $\mathcal{U}_{\mu^*} = 2^X$.
\end{Example}

\section{§ Свойства стандартного продолжения}

\begin{Theorem}{Свойства продолжения меры}{}
Пусть $\mu_0$ — мера на полукольце $\mathcal{P}$. Пусть $\mu$ — её стандартное продолжение (по Каратеодори) на $\sigma$-алгебру $\mathcal{U}$, а $\nu$ — какое-либо другое продолжение $\mu_0$ на некоторую $\sigma$-алгебру $\mathcal{U}_0$. Тогда:
\begin{enumerate}
    \item Для любого $A \in \mathcal{U} \cap \mathcal{U}_0$ выполняется $\nu(A) \le \mu(A)$. Более того, если $\mu(A) < +\infty$, то $\mu(A) = \nu(A)$.
    \item Если мера $\mu_0$ является $\textbf{$\sigma$-конечной}$ (то есть $X = \bigcup\limits_{i=1}^{\infty} P_i$, где $\mu_0(P_i) < +\infty$), то $\nu = \mu$ на пересечении $\mathcal{U} \cap \mathcal{U}_0$. (Теорема о единственности продолжения).
\end{enumerate}
\end{Theorem}

\begin{Proof}
\textbf{Доказательство пункта 1:} \\
Если $\mu(A) = +\infty$, то неравенство $\nu(A) \le +\infty$ выполняется очевидным образом.

Пусть $\mu(A) < +\infty$. По определению внешней меры, для любого $\varepsilon > 0$ существует покрытие $A \subset \bigcup\limits_{i=1}^{\infty} P_i$ ($P_i \in \mathcal{P}$), такое что:
$$ \sum\limits_{i=1}^{\infty} \mu_0(P_i) \le \mu(A) + \varepsilon $$
Так как $\nu$ — продолжение меры $\mu_0$, то $\nu(P_i) = \mu_0(P_i)$. В силу счётной полуаддитивности меры $\nu$:
$$ \nu(A) \le \sum\limits_{i=1}^{\infty} \nu(P_i) = \sum\limits_{i=1}^{\infty} \mu_0(P_i) \le \mu(A) + \varepsilon $$
Переходя к точной нижней грани по всем покрытиям ($\varepsilon \to 0$), получаем $\nu(A) \le \mu(A)$.

\textbf{Покажем равенство при $\mu(A) < +\infty$:} \\
Сначала докажем вспомогательный факт: $\forall P \in \mathcal{P}$, $\nu(A \cap P) = \mu(A \cap P)$. \\
Предположим противное: пусть $\nu(A \cap P) < \mu(A \cap P)$. Так как $P = (P \cap A) \sqcup (P \setminus A)$, то:
$$ \nu(P) = \nu(P \cap A) + \nu(P \setminus A) < \mu(P \cap A) + \mu(P \setminus A) = \mu(P) $$
Но $P \in \mathcal{P}$, а на полукольце обе меры совпадают, то есть $\nu(P) = \mu_0(P) = \mu(P)$. Получили противоречие: $\mu(P) < \mu(P)$. Следовательно, $\nu(A \cap P) = \mu(A \cap P)$.

Теперь рассмотрим покрытие $A \subset \bigcup\limits_{i=1}^{\infty} P_i$. По свойствам полукольца, объединение $\bigcup P_i$ можно представить как дизъюнктное объединение элементов полукольца: $\bigsqcup\limits_{i=1}^{\infty} \bigsqcup\limits_{j=1}^{m_i} Q_{ij}$, где $Q_{ij} \in \mathcal{P}$.
Тогда само множество $A$ можно записать как дизъюнктное объединение: $A = \bigsqcup\limits_{i, j} (A \cap Q_{ij})$.
Используя счётную аддитивность и доказанный факт $\nu(A \cap Q_{ij}) = \mu(A \cap Q_{ij})$:
$$ \nu(A) = \sum\limits_{i, j} \nu(A \cap Q_{ij}) = \sum\limits_{i, j} \mu(A \cap Q_{ij}) = \mu(A) $$

\textbf{Доказательство пункта 2 (Единственность продолжения):} \\
По условию мера $\mu_0$ является $\sigma$-конечной, значит, пространство можно представить как $X = \bigcup\limits_{n=1}^{\infty} P_n$, где $P_n \in \mathcal{P}$ и $\mu_0(P_n) < +\infty$.
Обозначим $S_n = \bigcup\limits_{k=1}^n P_k$. По лемме о нарезке, это конечное объединение можно представить как дизъюнктное объединение элементов полукольца $Q_j \in \mathcal{P}$:
$$ S_n = \bigsqcup\limits_{j=1}^{m_n} Q_j $$
Так как $\mu_0(Q_j) < +\infty$, то $\mu(Q_j) < +\infty$.
Рассмотрим произвольное измеримое множество $A \in \mathcal{U} \cap \mathcal{U}_0$. Пересечем его с $Q_j$:
Множество $A \cap Q_j \subset Q_j$, значит, $\mu(A \cap Q_j) \le \mu(Q_j) < +\infty$. 
В силу доказанного пункта 1, для множеств конечной меры продолжения совпадают, поэтому $\nu(A \cap Q_j) = \mu(A \cap Q_j)$.
Суммируя по дизъюнктным частям, получаем совпадение мер на пересечении с $S_n$:
$$ \nu(A \cap S_n) = \sum\limits_{j=1}^{m_n} \nu(A \cap Q_j) = \sum\limits_{j=1}^{m_n} \mu(A \cap Q_j) = \mu(A \cap S_n) $$
Заметим, что $S_1 \subset S_2 \subset \dots$, причем $\bigcup\limits_{n=1}^{\infty} S_n = X$. Следовательно, последовательность множеств $(A \cap S_n)$ монотонно возрастает и $A \cap S_n \nearrow A$. 
Поскольку $\mu$ и $\nu$ являются мерами на $\sigma$-алгебре, они обладают свойством непрерывности снизу. Переходя к пределу при $n \to \infty$, получаем:
$$ \nu(A) = \lim\limits_{n \to \infty} \nu(A \cap S_n) = \lim\limits_{n \to \infty} \mu(A \cap S_n) = \mu(A) $$
Теорема о единственности продолжения полностью доказана.
\end{Proof}


\section{§ Классы множеств и аппроксимация меры}

\begin{Definition}{Множества типа $\mathcal{U}_\sigma$ и $\mathcal{U}_\delta$}{}
Пусть $\mathcal{U}$ — некоторая система множеств.
\begin{itemize}
    \item Множество $A$ называется множеством типа $\mathbf{\mathcal{U}_\sigma}$, если оно представимо в виде не более чем счётного объединения элементов из $\mathcal{U}$:
    $$ A = \bigcup\limits_{i=1}^{\infty} A_i, \quad A_i \in \mathcal{U} $$
    \item Множество $A$ называется множеством типа $\mathbf{\mathcal{U}_\delta}$, если оно представимо в виде не более чем счётного пересечения элементов из $\mathcal{U}$:
    $$ A = \bigcap\limits_{i=1}^{\infty} A_i, \quad A_i \in \mathcal{U} $$
\end{itemize}
\end{Definition}

\noindent \textbf{NB:} Справедливы следующие обозначения для итераций: $\mathcal{U}_{\sigma\delta} = (\mathcal{U}_\sigma)_\delta$, $\quad \mathcal{U}_{\delta\sigma} = (\mathcal{U}_\delta)_\sigma$.

\begin{Lemma}{О структуре открытых множеств в $\mathbb{R}^n$}{}
Любое непустое открытое в $\mathbb{R}^n$ множество является множеством типа $(\mathcal{P}^n)_\sigma$. Более того, оно есть объединение не более чем счётного числа попарно дизъюнктных $n$-мерных ячеек из полукольца $\mathcal{P}^n$.
\end{Lemma}

\begin{Proof}
Пусть $G \subset \mathbb{R}^n$ — непустое открытое множество.
Так как $G$ открыто, для любой точки $x \in G$ существует окрестность, целиком лежащая в $G$. Значит, существует $n$-мерный полуинтервал $[a_x, b_x) \in \mathcal{P}^n$ с рациональными координатами вершин, такой что:
$$ x \in [a_x, b_x) \quad \text{и} \quad [a_x, b_x) \subset G $$
Очевидно, что $G = \bigcup\limits_{x \in G} [a_x, b_x)$. 
Так как множество всех полуинтервалов с рациональными концами счётно, то это объединение содержит не более чем счётное число различных элементов полукольца $\mathcal{P}^n$. Значит, $G = \bigcup\limits_{i=1}^{\infty} P_i$, где $P_i \in \mathcal{P}^n$. \\
По свойству полукольца (лемма «про нарезку»), любое счётное объединение элементов полукольца можно представить в виде дизъюнктного объединения элементов того же полукольца.
\end{Proof}

\begin{Definition}{Борелевская оболочка}{}
Минимальная $\sigma$-алгебра, содержащая заданную систему множеств $\mathcal{A}$, называется \textbf{борелевской оболочкой} системы $\mathcal{A}$ и обозначается $\mathcal{B}(\mathcal{A})$.
\end{Definition}

\begin{Theorem}{Об аппроксимации внешней меры}{}
Пусть $\mu$ — стандартное продолжение меры $\mu_0$ с полукольца $\mathcal{P}$. Если $\mu^*(E) < +\infty$, то существует множество $C \in \mathcal{P}_{\sigma\delta} \subset \mathcal{B}(\mathcal{P})$ такое, что:
$$ E \subset C \quad \text{и} \quad \mu^*(E) = \mu(C) $$
\end{Theorem}

\begin{Proof}
По определению внешней меры, порождённой мерой $\mu_0$:
$$ \mu^*(E) = \inf \left\{ \sum\limits_{i=1}^{\infty} \mu_0(P_i) : E \subset \bigcup\limits_{i=1}^{\infty} P_i, \quad P_i \in \mathcal{P} \right\} $$
По свойству инфимума, для любого $n \in \mathbb{N}$ существует покрывающий набор элементов $P_i^n \in \mathcal{P}$, такой что $E \subset \bigcup\limits_{i=1}^{\infty} P_i^n$ и:
$$ \mu^*(E) + \frac{1}{n} > \sum\limits_{i=1}^{\infty} \mu_0(P_i^n) $$
Обозначим объединение этого набора за $C_n = \bigcup\limits_{i=1}^{\infty} P_i^n$. Очевидно, что $C_n \in \mathcal{P}_\sigma$ и $E \subset C_n$. 
Из счётной полуаддитивности внешней меры $\mu^*$ и совпадения $\mu^*$ с $\mu_0$ на полукольце следует:
$$ \mu^*(E) \le \mu^*(C_n) \le \sum\limits_{i=1}^{\infty} \mu^*(P_i^n) = \sum\limits_{i=1}^{\infty} \mu_0(P_i^n) < \mu^*(E) + \frac{1}{n} $$
Определим множество $C = \bigcap\limits_{n=1}^{\infty} C_n$. Так как это счётное пересечение множеств типа $\mathcal{P}_\sigma$, то $C \in \mathcal{P}_{\sigma\delta} \subset \mathcal{B}(\mathcal{P})$.
Поскольку $E \subset C_n$ для всех $n$, то $E \subset C$. 
Тогда из монотонности меры: $\mu^*(E) \le \mu(C) \le \mu(C_n) < \mu^*(E) + \frac{1}{n}$ для любого $n$. Устремляя $n \to \infty$, получаем $\mu(C) = \mu^*(E)$.
\end{Proof}

\begin{Theorem}{Следствие (Аппроксимация измеримого множества)}{}
Пусть $\mu$ — стандартное продолжение меры $\mu_0$ с полукольца $\mathcal{P}$, и множество $A$ измеримо ($\mu(A) < +\infty$). Тогда существуют борелевские множества $B, C \in \mathcal{B}(\mathcal{P})$, такие что:
$$ B \subset A \subset C \quad \text{и} \quad \mu(C \setminus B) = 0 $$
\end{Theorem}

\begin{Proof}
\textit{}

\textbf{1) Построение $C$ (оценка снаружи):} \\
Так как $\mu(A) < +\infty$, по теореме об аппроксимации существует множество $C \in \mathcal{P}_{\sigma\delta} \subset \mathcal{B}(\mathcal{P})$ такое, что $A \subset C$ и $\mu(C) = \mu(A)$. \\
Так как $A$ измеримо и $\mu(A)$ конечна, разность $C \setminus A$ также измерима, и её мера равна:
$$ \mu(C \setminus A) = \mu(C) - \mu(A) = 0 $$

\textbf{2) Построение $B$ (оценка изнутри):} \\
Множество $C \setminus A$ имеет меру ноль. Применим теорему об аппроксимации к этому множеству: существует множество $N \in \mathcal{B}(\mathcal{P})$ такое, что $(C \setminus A) \subset N$ и $\mu(N) = \mu(C \setminus A) = 0$.

Определим множество $B = C \setminus N$. Так как $C$ и $N$ — борелевские множества, то и $B \in \mathcal{B}(\mathcal{P})$. \\
Покажем, что $B \subset A$. Заметим, что $(C \setminus A) \subset N$, следовательно $C \setminus N \subset C \setminus (C \setminus A) = A$. Таким образом, $B \subset A \subset C$.

\textbf{3) Проверка разности:} \\
Мера разности $\mu(C \setminus B) = \mu(C \setminus (C \setminus N)) = \mu(C \cap N)$. 
Поскольку $C \cap N \subset N$, в силу монотонности меры:
$$ \mu(C \setminus B) \le \mu(N) = 0 \implies \mu(C \setminus B) = 0 $$
Что и требовалось доказать.
\end{Proof}

\section{§ Мера Лебега и её свойства}

\begin{Definition}{Мера Лебега в $\mathbb{R}^n$}{}
\textbf{Мерой Лебега} в $\mathbb{R}^n$ называется стандартное продолжение (по Каратеодори) $\lambda_n$ объёма с полукольца многомерных ячеек $\mathcal{P}^n$.
\end{Definition}

Напомним конструкцию:
1) На $\mathcal{P}^1 = \{[a, b)\}$ мера $\lambda_1([a, b)) = b - a$.
2) $\mathcal{P}^n = \mathcal{P}^1 \odot \dots \odot \mathcal{P}^1$ (декартово произведение $n$ раз) является полукольцом. Объём на нём задаётся как $\lambda_n = \lambda_1 \cdot \dots \cdot \lambda_1$, то есть $\lambda_n([a, b)) = \prod\limits_{k=1}^n (b_k - a_k)$.

\begin{Theorem}{Счётная аддитивность объёма на $\mathcal{P}^n$}{}
Функция объёма $\lambda_n$ является мерой (то есть счётно аддитивна) на полукольце $\mathcal{P}^n$.
\end{Theorem}

\begin{Proof}
Достаточно доказать свойство счётной полуаддитивности: если $[a, b) \subset \bigcup\limits_{i=1}^{\infty} [a_i, b_i)$, то $\lambda_n([a, b)) \le \sum\limits_{i=1}^{\infty} \lambda_n([a_i, b_i))$.

Зафиксируем произвольное $\varepsilon > 0$. Из-за того, что объём непрерывен как функция от координат, мы можем слегка "раздуть" каждый элемент покрытия. Существуют векторы $a_i^\varepsilon < a_i$ (покоординатно), такие что:
$$ \lambda_n([a_i^\varepsilon, b_i)) < \lambda_n([a_i, b_i)) + \frac{\varepsilon}{2^i} $$
Рассмотрим любую точку $t \in [a, b)$ (то есть $a \le t < b$). Замкнутый параллелепипед $[a, t]$ является компактом в $\mathbb{R}^n$. 
Справедлива цепочка включений:
$$ [a, t] \subset [a, b) \subset \bigcup\limits_{i=1}^{\infty} [a_i^\varepsilon, b_i) \subset \bigcup\limits_{i=1}^{\infty} (a_i^\varepsilon, b_i) $$
Мы покрыли компакт $[a, t]$ открытыми множествами $(a_i^\varepsilon, b_i)$. По лемме Гейне-Бореля из этого покрытия можно выделить конечное подпокрытие. Значит, существует такой номер $N$, что:
$$ [a, t] \subset \bigcup\limits_{i=1}^N (a_i^\varepsilon, b_i) $$
Отсюда следует, что:
$$ [a, t) \subset [a, t] \subset \bigcup\limits_{i=1}^N (a_i^\varepsilon, b_i) \subset \bigcup\limits_{i=1}^N [a_i^\varepsilon, b_i) $$
Так как для конечного объединения в полукольце действует конечная полуаддитивность объёма, то:
$$ \lambda_n([a, t)) \le \sum\limits_{i=1}^N \lambda_n([a_i^\varepsilon, b_i)) \le \sum\limits_{i=1}^{\infty} \lambda_n([a_i^\varepsilon, b_i)) < \sum\limits_{i=1}^{\infty} \lambda_n([a_i, b_i)) + \sum\limits_{i=1}^{\infty} \frac{\varepsilon}{2^i} $$
То есть $\lambda_n([a, t)) < \sum\limits_{i=1}^{\infty} \lambda_n([a_i, b_i)) + \varepsilon$.
Это неравенство справедливо для всех $t < b$. Устремим $t \to b$. В силу непрерывности объёма получим:
$$ \lambda_n([a, b)) \le \sum\limits_{i=1}^{\infty} \lambda_n([a_i, b_i)) + \varepsilon $$
Так как $\varepsilon > 0$ было выбрано произвольно ($\varepsilon \downarrow 0$), получаем $\lambda_n([a, b)) \le \sum\limits_{i=1}^{\infty} \lambda_n([a_i, b_i))$. Счётная полуаддитивность доказана, значит объём является мерой.
\end{Proof}

\section{§ Свойства меры Лебега и измеримых множеств}

\begin{Theorem}{Базовые свойства измеримых множеств}{}
\begin{enumerate}
    \item Открытые множества измеримы. Мера непустого открытого множества строго положительна.
    \item Замкнутые множества измеримы.
    \item Мера одноточечного множества равна 0.
    \item Мера любого не более чем счётного множества равна 0.
    \item Мера измеримого ограниченного множества конечна. Всякое измеримое множество — это объединение не более чем счётного числа множеств конечной меры.
    \item Гиперплоскости имеют меру нуль. Если $H_{n-1} = \{ x \in \mathbb{R}^n : w_1 x_1 + \dots + w_n x_n + w_0 = 0 \}$, то $\lambda_m(H_{n-1}) = 0$ при $m \ge n$.
    \item Пусть $D \subset \mathbb{R}^m$ и $\Phi: D \to \mathbb{R}^n$, причем $\Phi \in C^1(D)$. Если $A \subset D$ измеримо относительно $\lambda_m$, то $\Phi(A)$ измеримо относительно $\lambda_n$.
    \item Измеримость и мера Лебега сохраняются при движениях пространства (сдвиг + ортогональное преобразование).
\end{enumerate}
\end{Theorem}

\begin{Proof}[Развёрнутые доказательства (упражнения с лекции)]
\textbf{1) Открытые множества:} Любое открытое в $\mathbb{R}^n$ множество $G$ можно представить как объединение не более чем счётного числа $n$-мерных ячеек (прямоугольников) из полукольца $\mathcal{P}^n$. Так как полукольцо лежит в $\sigma$-алгебре измеримых множеств $\mathcal{U}_\lambda$, а $\sigma$-алгебра замкнута относительно счётных объединений, то $G \in \mathcal{U}_\lambda$. Если $G \neq \emptyset$, оно содержит хотя бы одну ячейку со строго положительными длинами сторон, значит $\lambda(G) > 0$.

\textbf{2) Замкнутые множества:} Любое замкнутое множество $F$ является дополнением к некоторому открытому множеству $G$ (то есть $F = \mathbb{R}^n \setminus G$). Так как открытое $G$ измеримо, а $\sigma$-алгебра замкнута относительно перехода к дополнению, то $F$ измеримо.

\textbf{3) Одноточечное множество:} Точку $x = (x_1, \dots, x_n)$ можно представить как счётное пересечение вложенных полуинтервалов: 
$$ \{x\} = \bigcap\limits_{k=1}^{\infty} \left[ x_1, x_1 + \frac{1}{k} \right) \times \dots \times \left[ x_n, x_n + \frac{1}{k} \right) $$
Мера каждого такого полуинтервала равна $\frac{1}{k^n}$. В силу непрерывности меры Лебега сверху: $\lambda(\{x\}) = \lim\limits_{k \to \infty} \frac{1}{k^n} = 0$.

\textbf{4) Счётное множество:} Пусть $A = \{x_1, x_2, \dots\}$. Представим его как счётное объединение одноточечных множеств: $A = \bigcup\limits_{i=1}^{\infty} \{x_i\}$. В силу счётной полуаддитивности (или аддитивности, так как точки не пересекаются) меры: 
$$ \lambda(A) \le \sum\limits_{i=1}^{\infty} \lambda(\{x_i\}) = \sum\limits_{i=1}^{\infty} 0 = 0 $$

\textbf{5) Ограниченные множества:} Если множество $A$ ограничено, оно целиком помещается внутри некоторого большого куба $K = [-R, R]^n$. В силу монотонности меры: $\lambda(A) \le \lambda(K) = (2R)^n < +\infty$.
Любое пространство $\mathbb{R}^n$ можно разбить на счётное объединение кубов $K_m = [-m, m]^n$. Тогда любое измеримое множество $A = \bigcup\limits_{m=1}^{\infty} (A \cap K_m)$, где мера каждого куска конечна.

\textbf{6) Гиперплоскость:} Ортогональным преобразованием (которое сохраняет меру Лебега, см. п. 8) любую гиперплоскость можно перевести в координатную плоскость $x_n = 0$. Тогда она представляет собой прямое произведение $\mathbb{R}^{n-1} \times \{0\}$. Мера этого множества равна произведению $(+\infty) \cdot 0 = 0$.

\textbf{7) Гладкие отображения:} Гладкая функция $\Phi \in C^1$ на компакте удовлетворяет условию Липшица. Известно, что липшицевы отображения переводят множества лебеговой меры нуль в множества меры нуль, а так как любое измеримое множество есть объединение борелевского (которое переходит в борелевское) и множества меры нуль, образ остаётся измеримым.

\textbf{8) Инвариантность при движениях:} Сдвиг не меняет длин рёбер ячеек из полукольца, поэтому мера Лебега трансляционно инвариантна по определению. При ортогональном преобразовании (повороте/отражении) матрица Якоби имеет определитель $\pm 1$. При замене переменных в кратном интеграле объём множества умножается на модуль определителя Якоби, то есть на $| \pm 1 | = 1$. Следовательно, объём (и мера) сохраняется.
\end{Proof}

\begin{Example}[ Множество нулевой меры мощности континуум]
Рассмотрим \textbf{множество Кантора}. Построим его итеративно на отрезке $[0, 1]$:
\begin{itemize}
    \item $C_0 = [0, 1]$
    \item $C_1 = \left[0, \frac{1}{3}\right] \cup \left[\frac{2}{3}, 1\right]$ (сегменты первого ранга, удаляем среднюю треть)
    \item $C_2 = \left[0, \frac{1}{9}\right] \cup \left[\frac{2}{9}, \frac{3}{9}\right] \cup \left[\frac{6}{9}, \frac{7}{9}\right] \cup \left[\frac{8}{9}, 1\right]$
\end{itemize}
Само множество Кантора определяется как бесконечное пересечение: $C = \bigcap\limits_{i=1}^{\infty} C_i$.

\begin{center}
\begin{tikzpicture}[scale=1.2, >=latex]
    \definecolor{darktext}{HTML}{154360}
    \definecolor{colorPoi}{HTML}{7DCEA0}
    
    % C_0
    \draw[darktext, ultra thick] (0, 0) -- (9, 0);
    \node[left, darktext] at (-0.5, 0) {$C_0$};
    
    % C_1
    \draw[darktext, ultra thick] (0, -0.6) -- (3, -0.6);
    \draw[darktext, ultra thick] (6, -0.6) -- (9, -0.6);
    \node[left, darktext] at (-0.5, -0.6) {$C_1$};
    
    % C_2
    \draw[darktext, ultra thick] (0, -1.2) -- (1, -1.2);
    \draw[darktext, ultra thick] (2, -1.2) -- (3, -1.2);
    \draw[darktext, ultra thick] (6, -1.2) -- (7, -1.2);
    \draw[darktext, ultra thick] (8, -1.2) -- (9, -1.2);
    \node[left, darktext] at (-0.5, -1.2) {$C_2$};
    
    % C_3 (набросок точек)
    \foreach \x in {0, 0.33, 0.66, 1, 2, 2.33, 2.66, 3, 6, 6.33, 6.66, 7, 8, 8.33, 8.66, 9} {
        \fill[colorPoi] (\x, -1.8) circle (1.5pt);
    }
    \node[left, darktext] at (-0.5, -1.8) {$C_\infty$};
\end{tikzpicture}
\end{center}
\end{Example}

\begin{Lemma}{Свойства множества Кантора}{}
1) Множество $C$ измеримо относительно $\lambda_1$. \\
2) $\lambda_1(C) = 0$. \\
3) Мощность множества $C$ — континуум.
\end{Lemma}

\begin{Proof}
1) Измеримо, так как состоит из пересечения счётного числа измеримых множеств (каждое $C_i$ — конечное объединение отрезков). \\
2) На каждом шаге $i$ длина оставшихся отрезков составляет $2^i \cdot \frac{1}{3^i} = \left(\frac{2}{3}\right)^i$. Так как $C \subset C_i$, то:
$$ \lambda_1(C) \le \lambda_1(C_i) = \left(\frac{2}{3}\right)^i \xrightarrow{i \to \infty} 0 \implies \lambda_1(C) = 0 $$
3) Существует биекция между элементами множества Кантора и множеством всех двоичных последовательностей бесконечной длины, мощность которого равна континууму.
\end{Proof}

\section{§ Неизмеримые множества}

\begin{Theorem}{Множество Витали}{}
Любое множество положительной меры Лебега содержит неизмеримое подмножество.
\end{Theorem}

\begin{Proof}
Пусть существует измеримое множество $A$, такое что $0 < \lambda_n(A) < +\infty$. \\
Будем считать элементы эквивалентными $x \sim y$ для $x, y \in A$, если их разность $x - y \in \mathbb{Q}^n$ (вектор с рациональными координатами). Это задаёт отношение эквивалентности. \\
Следовательно, множество $A$ разбивается на дизъюнктные классы эквивалентности. \\
В силу аксиомы выбора, мы можем создать множество $E$, содержащее \textbf{ровно по одному представителю} из каждого класса эквивалентности.

Так как $A$ имеет конечную меру, оно ограничено. Пусть $\forall x \in A : \|x\| \le R$. \\
Следовательно, для любых $x, y \in A$ расстояние $\|x - y\| \le 2R$. \\
Рассмотрим счётное объединение сдвигов множества $E$ на рациональные векторы $r$:
$$ W = \bigcup\limits_{r \in \mathbb{Q}^n, \|r\| \le 2R} (r + E) $$
По построению множества представителей, $A \subset W$. 
Предположим от противного, что $E$ \textbf{измеримо}. Тогда каждый сдвиг $(r + E)$ также измерим (мера инвариантна относительно сдвига), а значит и объединение $W$ измеримо.
Так как сдвиги $(r + E)$ попарно не пересекаются, мера их счётного объединения равна сумме мер:
$$ 0 < \lambda_n(A) \le \lambda_n(W) = \sum\limits_{k=1}^{\infty} \lambda_n(r_k + E) = \sum\limits_{k=1}^{\infty} \lambda_n(E) $$
\begin{itemize}
    \item Если $\lambda_n(E) = 0$, то $\lambda_n(W) = \sum 0 = 0$, что противоречит $0 < \lambda_n(A) \le \lambda_n(W)$.
    \item Если $\lambda_n(E) > 0$, то сумма бесконечного числа одинаковых положительных чисел равна $+\infty$, то есть $\lambda_n(W) = +\infty$. Однако множество $W$ ограничено радиусом $3R$, поэтому его мера должна быть конечна. Противоречие.
\end{itemize}
Следовательно, допущение об измеримости неверно. Множество $E \subset A$ является неизмеримым.
\end{Proof}

\noindent \textbf{Замечание (О гладких преобразованиях):} Отказаться совсем от гладкости $\Phi$ как достаточного условия сохранения измеримости (см. свойство 7) нельзя. Классический контрпример строится на базе функции Кантора (лестницы Кантора) $\varphi$. Можно построить непрерывную (но не абсолютно непрерывную) функцию, которая отображает множество меры нуль на множество положительной меры (содержащее неизмеримое подмножество $E$). Тогда $e = \varphi^{-1}(E)$ является подмножеством множества меры нуль, следовательно $\lambda(e) = 0$ (измеримо), но его образ $E$ — неизмерим. Вывод: простой непрерывности недостаточно.

\section{§ Измеримые функции и их свойства}

Идея интеграла Лебега отличается от интеграла Римана тем, что мы разбиваем не ось $X$, а ось $Y$. 

\begin{center}
\begin{center}
\begin{center}
\begin{tikzpicture}[scale=1.2, >=latex]

    % --- Определение цветов ---
    \definecolor{darktext}{HTML}{154360}     % Темно-синий
    \definecolor{colorY}{HTML}{D4E6F1}       % Светло-голубой для полосы Y
    \definecolor{colorHighlight}{HTML}{2874A6} % Насыщенный синий для выделения множеств
    \definecolor{colorCurve}{HTML}{85929E}   % Серый для невыделенной части графика

    % ==========================================
    % 1. Горизонтальная разбивка по оси Y (Интеграл Лебега)
    % ==========================================
    
    % Закрашиваем исследуемый коридор по Y
    \fill[colorY, opacity=0.6] (-0.2, 1.8) rectangle (8.2, 2.5);
    
    % Линии уровней y_{i-1} и y_i
    \draw[dashed, colorHighlight, thick] (-0.2, 1.8) -- (8.2, 1.8) node[right] {$y_{i-1}$};
    \draw[dashed, colorHighlight, thick] (-0.2, 2.5) -- (8.2, 2.5) node[right] {$y_i$};

    % Оси координат
    \draw[->, thick, darktext] (-0.5, 0) -- (8.8, 0) node[right] {$x$};
    \draw[->, thick, darktext] (0, -0.5) -- (0, 4) node[above] {$f(x)$};

    % ==========================================
    % 2. Отрисовка функции f(x) кусками для точного попадания
    % ==========================================
    
    \draw[thick, colorCurve] (0,1) to[out=60, in=225] (0.5, 1.8);
    \draw[ultra thick, colorHighlight] (0.5, 1.8) to[out=45, in=225] (1.0, 2.5); 
    \draw[thick, colorCurve] (1.0, 2.5) to[out=45, in=180] (1.5, 3.5) to[out=0, in=135] (2.0, 2.5);
    \draw[ultra thick, colorHighlight] (2.0, 2.5) to[out=-45, in=135] (2.6, 1.8);
    \draw[thick, colorCurve] (2.6, 1.8) to[out=-45, in=180] (3.5, 1.2) to[out=0, in=225] (4.4, 1.8);
    \draw[ultra thick, colorHighlight] (4.4, 1.8) to[out=45, in=225] (4.9, 2.5);
    \draw[thick, colorCurve] (4.9, 2.5) to[out=45, in=180] (5.5, 3.2) to[out=0, in=135] (6.1, 2.5);
    \draw[ultra thick, colorHighlight] (6.1, 2.5) to[out=-45, in=135] (6.8, 1.8);
    \draw[thick, colorCurve] (6.8, 1.8) to[out=-45, in=135] (7.8, 0.5);

    % ==========================================
    % 3. Проекция на ось X (формирование множества e_i)
    % ==========================================
    
    \foreach \x in {0.5, 1.0, 2.0, 2.6, 4.4, 4.9, 6.1, 6.8} {
        \draw[dashed, colorHighlight!60!white] (\x, 1.8) -- (\x, 0);
    }
    
    \draw[dashed, colorHighlight!60!white] (1.0, 2.5) -- (1.0, 1.8);
    \draw[dashed, colorHighlight!60!white] (2.0, 2.5) -- (2.0, 1.8);
    \draw[dashed, colorHighlight!60!white] (4.9, 2.5) -- (4.9, 1.8);
    \draw[dashed, colorHighlight!60!white] (6.1, 2.5) -- (6.1, 1.8);

    \draw[line width=4pt, colorHighlight] (0.5, 0) -- (1.0, 0);
    \draw[line width=4pt, colorHighlight] (2.0, 0) -- (2.6, 0);
    \draw[line width=4pt, colorHighlight] (4.4, 0) -- (4.9, 0);
    \draw[line width=4pt, colorHighlight] (6.1, 0) -- (6.8, 0);

    % ==========================================
    % 4. Подписи и Легенда (перенесено ВНИЗ, чтобы ничего не перекрывать)
    % ==========================================
    
    % Квадратная "скобка", объединяющая отрезки
    \draw[darktext, thick] (0.5, -0.1) -- (0.5, -0.3) -- (6.8, -0.3) -- (6.8, -0.1);
    \node[darktext, below=2pt] at (3.65, -0.3) {
        $\mathbf{e_i = f^{-1}([y_{i-1}, y_i))}$ --- измеримое лебегово множество
    };

    % Пояснительный текст (строго под графиком)
    \node[darktext, align=center] at (3.65, -3) {
        \textbf{Идея интеграла Лебега:} \\
        Мы разбиваем не ось аргументов $X$, а ось значений $Y$. \\
        Одному шагу по $Y$ может соответствовать множество $e_i$, \\
        состоящее из множества разрозненных кусков на оси $X$. \\
        \vspace{0.5em} \\
        Сумма Лебега: $\sum\limits_{i=1}^n y_i \cdot \mu(e_i)$
    };

\end{tikzpicture}
\end{center}
\end{center}

\begin{Definition}{Лебеговы множества}{}
Пусть дана функция $f: X \to \overline{\mathbb{R}}$. \textbf{Лебеговыми множествами} функции $f$, соответствующими значению $a \in \mathbb{R}$, называются следующие подмножества $X$:
\begin{enumerate}
    \item[a)] $X(f > a) = \{ x \in X : f(x) > a \}$
    \item[б)] $X(f \ge a) = \{ x \in X : f(x) \ge a \}$
    \item[в)] $X(f < a) = \{ x \in X : f(x) < a \}$
    \item[г)] $X(f \le a) = \{ x \in X : f(x) \le a \}$
\end{enumerate}
\end{Definition}

\begin{Lemma}{Эквивалентность измеримости}{}
Пусть $X$ — множество, $\mathcal{U}$ — $\sigma$-алгебра. Измеримость (принадлежность $\mathcal{U}$) для любого $a \in \mathbb{R}$ любого из четырех типов лебеговых множеств равносильна.
\end{Lemma}

\begin{Proof}
Докажем по цепочке переходов, используя свойства $\sigma$-алгебры:
\begin{itemize}
    \item $a \Rightarrow \text{б}$: Заметим, что нестрогое неравенство можно представить как счётное пересечение строгих:
    $$ X(f \ge a) = \bigcap\limits_{n=1}^{\infty} X\left(f > a - \frac{1}{n}\right) $$
    Счётное пересечение измеримых множеств измеримо.
    \item $\text{б} \Rightarrow \text{в}$: Переход к дополнению: $X(f < a) = X \setminus X(f \ge a)$. Дополнение измеримо.
    \item $\text{в} \Rightarrow \text{г}$: $X(f \le a) = \bigcap\limits_{n=1}^{\infty} X\left(f < a + \frac{1}{n}\right)$. Измеримо.
    \item $\text{г} \Rightarrow \text{a}$: $X(f > a) = X \setminus X(f \le a)$. Измеримо.
\end{itemize}
Все определения эквивалентны.
\end{Proof}

\begin{Definition}{Измеримая функция}{}
Функция $f$ называется \textbf{измеримой}, если все её лебеговы множества $X(f > a)$ (а следовательно, и все остальные типы) являются измеримыми для любого $a \in \mathbb{R}$.
\end{Definition}

\begin{Lemma}{Свойства измеримых функций}{}
Пусть $f$ — измеримая функция. Тогда:
\begin{enumerate}
    \item Прообразы одноточечных множеств $f^{-1}(\{a\})$ измеримы.
    \item Прообраз любого промежутка измерим.
    \item Функция модуля $|f|$ измерима.
    \item Прообраз любого открытого множества измерим.
    \item Сужение функции $f$ на измеримое подмножество $E \subset X$ измеримо.
    \item Если $f$ измерима на каждом $E_i$, то $f$ измерима на их объединении $\bigcup\limits_{i=1}^{\infty} E_i$.
\end{enumerate}
\end{Lemma}
\end{center}
\begin{Proof}
\textbf{1)} Прообраз точки — это пересечение двух лебеговых множеств: 
$$ f^{-1}(\{a\}) = X(f \le a) \setminus X(f < a) $$

\textbf{2)} Для полуинтервала $[a, b)$:
$$ f^{-1}([a, b)) = X(f < b) \setminus X(f < a) $$

\textbf{3)} По определению лебегова множества для модуля:
$$ X(|f| < a) = X(-a < f < a) = X(f < a) \cap X(f > -a) $$
Пересечение измеримых множеств измеримо.

\textbf{4)} Пусть $G$ — открытое множество. По теореме о структуре открытых множеств, $G$ представимо в виде счётного объединения интервалов: $G = \bigcup\limits_{i=1}^{\infty} (a_i, b_i)$.
Прообраз объединения равен объединению прообразов:
$$ f^{-1}(G) = f^{-1}\left(\bigcup\limits_{i=1}^{\infty} (a_i, b_i)\right) = \bigcup\limits_{i=1}^{\infty} f^{-1}((a_i, b_i)) $$
Так как прообраз интервала измерим (свойство 2), а счётное объединение измеримых множеств измеримо, то $f^{-1}(G)$ измеримо.

\textbf{5) (Доказательство упражнения):} Пусть $f$ измерима на $X$, и задано измеримое подмножество $M \subset X$. Рассмотрим сужение $f|_M$. 
Лебегово множество для сужения имеет вид:
$$ M(f|_M > a) = \{ x \in M : f(x) > a \} $$
Очевидно, что это множество является пересечением:
$$ M(f|_M > a) = M \cap X(f > a) $$
Множество $M$ измеримо по условию. Множество $X(f > a)$ измеримо, так как $f$ измерима на $X$. Пересечение двух измеримых множеств измеримо. Значит, функция $f|_M$ измерима.

\textbf{6) (Доказательство упражнения):} Пусть $E = \bigcup\limits_{i=1}^{\infty} E_i$, и на каждом $E_i$ функция $f$ измерима. 
Рассмотрим лебегово множество на объединении $E$:
$$ E(f > a) = \{ x \in E : f(x) > a \} = \bigcup\limits_{i=1}^{\infty} \{ x \in E_i : f(x) > a \} = \bigcup\limits_{i=1}^{\infty} E_i(f > a) $$
По условию функция $f$ измерима на каждом $E_i$, следовательно, каждое множество $E_i(f > a)$ измеримо. Так как $\sigma$-алгебра замкнута относительно счётных объединений, итоговое множество $E(f > a)$ измеримо.
\end{Proof}

\section{§ Замкнутость измеримых функций}

\begin{Theorem}{Замкнутость относительно предельного перехода}{}
Пусть $\varphi_n$ — последовательность измеримых функций. Тогда:
\begin{enumerate}
    \item Функции $g(x) = \inf\limits_n \varphi_n(x)$ и $h(x) = \sup\limits_n \varphi_n(x)$ измеримы.
    \item Верхний предел $\varlimsup \varphi_n$ и нижний предел $\varliminf \varphi_n$ измеримы.
\end{enumerate}
В частности, если существует поточечный предел $\lim\limits_{n \to \infty} \varphi_n(x)$, то он является измеримой функцией.
\end{Theorem}

\begin{Proof}
Докажем пункт 1, проверив лебеговы множества для точных граней.
Для инфимума: значение инфимума строго меньше $a$ тогда и только тогда, когда хотя бы один элемент последовательности строго меньше $a$. 
$$ X(g < a) = \bigcup\limits_{n=1}^{\infty} X(\varphi_n < a) $$
Счётное объединение измеримых множеств измеримо.

Для супремума: значение супремума не превосходит $a$ тогда и только тогда, когда \textbf{все} элементы последовательности не превосходят $a$.
$$ X(h \le a) = \bigcap\limits_{n=1}^{\infty} X(\varphi_n \le a) $$
Счётное пересечение измеримых множеств измеримо.
\end{Proof}

\begin{Example}[ к предельному переходу]
Рассмотрим последовательность $\varphi_n(x) = x^n$ на отрезке $X = [0, 1]$. Очевидно, что каждая функция непрерывна, а значит измерима.
Найдём точные грани:
$$ g(x) = \inf\limits_n (x^n) = \begin{cases} 0, & x < 1 \\ 1, & x = 1 \end{cases} \quad \text{и} \quad h(x) = \sup\limits_n (x^n) = x $$
Проверим лебеговы множества для $a = 1/2$:
$$ X\left(g < \frac{1}{2}\right) = [0, 1) \quad \text{и} \quad X\left(h < \frac{1}{2}\right) = \left[0, \frac{1}{2}\right) $$
Оба полученных множества измеримы (принадлежат борелевской оболочке), что подтверждает теорему.
\end{Example}

\section{§ Арифметические свойства измеримых функций}

\begin{Theorem}{Суперпозиция измеримых и непрерывной функции}{}
Пусть $f_1, \dots, f_n : X \to \mathbb{R}$ — измеримые функции, а функция $\varphi : Y \to \mathbb{R}$ (где $Y \subset \mathbb{R}^n$) является непрерывной ($\varphi \in C(Y)$). 
Тогда сложная функция $\varphi(f_1(x), \dots, f_n(x))$ также является измеримой на $X$.
\end{Theorem}

\begin{Proof}
Рассмотрим вектор-функцию $F(x) = (f_1(x), \dots, f_n(x))$. 
Сначала докажем, что прообраз любого $n$-мерного полуинтервала $[a, b) = [a_1, b_1) \times \dots \times [a_n, b_n)$ при отображении $F$ является измеримым множеством. Действительно:
$$ F^{-1}([a, b)) = \bigcap\limits_{i=1}^n f_i^{-1}([a_i, b_i)) $$
Так как каждая $f_i$ измерима, прообразы $[a_i, b_i)$ измеримы. Конечное пересечение измеримых множеств также измеримо.

Теперь пусть $G$ — произвольное открытое множество в $\mathbb{R}^n$. Согласно теореме о структуре открытого множества, $G$ можно представить как не более чем счётное объединение ячеек: $G = \bigcup\limits_{k=1}^{\infty} [a^{(k)}, b^{(k)})$.
Тогда прообраз открытого множества:
$$ F^{-1}(G) = \bigcup\limits_{k=1}^{\infty} F^{-1}([a^{(k)}, b^{(k)})) $$
является измеримым как счётное объединение измеримых множеств.

Перейдем к функции $\varphi$. Рассмотрим её лебегово множество для значения $a$. Множество тех $y \in Y$, где $\varphi(y) < a$, есть:
$$ \{ y \in Y : \varphi(y) < a \} = \varphi^{-1}((-\infty, a)) = G_a \cap Y $$
где $G_a$ — некоторое открытое множество в $\mathbb{R}^n$ (в силу непрерывности $\varphi$).
Проверим измеримость лебегова множества сложной функции:
$$ X(\varphi(F(x)) < a) = \{ x \in X : F(x) \in G_a \cap Y \} = X(F^{-1}(G_a)) $$
Как мы доказали выше, прообраз открытого множества $G_a$ измерим. Следовательно, сложная функция измерима.
\end{Proof}

\noindent \textbf{NB (Важное замечание):} В расширенной числовой прямой $\overline{\mathbb{R}}$ мы доопределяем операции следующим образом:
$$ 0 \cdot (\pm \infty) = (\pm \infty) \cdot 0 = 0, \quad \frac{x}{\pm \infty} = 0 \text{ при } x \in \mathbb{R} $$
Операции вида $(+\infty) - (+ \infty)$ не определены (полагаем равными 0 по договоренности или избегаем их).

\begin{Theorem}{Арифметические свойства измеримых функций}{}
Пусть $f, g$ — измеримые функции. Тогда:
\begin{enumerate}
    \item $f + g$ — измерима.
    \item $f \cdot g$ — измерима.
    \item $c \cdot f$ — измерима (при $c \in \mathbb{R}$).
    \item $\frac{f}{g}$ — измерима (на множестве, где $g \neq 0$).
\end{enumerate}
\end{Theorem}

\begin{Proof}
Будем считать, что значения $f$ и $g$ конечны. Тогда достаточно применить предыдущую теорему о суперпозиции:
1) Функция $\varphi(x, y) = x + y \in C(\mathbb{R}^2)$ непрерывна, значит $f(x) + g(x)$ измерима.
2) Функция $\varphi(x, y) = x \cdot y \in C(\mathbb{R}^2)$ непрерывна, значит $f(x) \cdot g(x)$ измерима.
Остальные пункты доказываются аналогично.
\end{Proof}

\section{§ Простые функции и аппроксимация}

\begin{Definition}{Простая функция}{}
Измеримая функция называется \textbf{простой}, если множество её значений конечно.
\end{Definition}

\begin{Definition}{Допустимое разбиение}{}
Разбиение области определения $X$ простой функции $f$ называется \textbf{допустимым}, если:
\begin{enumerate}
    \item Элементы разбиения измеримы.
    \item На каждом элементе разбиения функция $f$ постоянна (не меняет значение).
\end{enumerate}
\end{Definition}

\begin{Lemma}{Общее допустимое разбиение}{}
Пусть $f, g$ — простые функции. Тогда для них существует общее допустимое разбиение.
\end{Lemma}

\begin{Proof}
Пусть $f$ принимает значения $a_1, \dots, a_n$. Тогда множества $E_i = f^{-1}(\{a_i\})$ ($i \in \{1 \dots n\}$) образуют допустимое разбиение для $f$. \\
Пусть $g$ принимает значения $b_1, \dots, b_k$. Тогда множества $G_j = g^{-1}(\{b_j\})$ ($j \in \{1 \dots k\}$) образуют допустимое разбиение для $g$. \\
Рассмотрим систему попарных пересечений $\{ E_i \cap G_j \}_{i, j}$. Это множество измеримо, и на нём обе функции постоянны. Следовательно, это и есть общее допустимое разбиение для $f$ и $g$ (некоторые пересечения могут быть пустыми, это нормально).
\end{Proof}

\noindent \textbf{Следствие:} Сумма и произведение простых функций — тоже простые функции (доказывается применением леммы об общем разбиении).

\begin{Theorem}{Об аппроксимации измеримой функции}{}
Любая неотрицательная измеримая функция является поточечным пределом последовательности простых функций. Причем последовательность можно считать монотонно возрастающей ($f_1 \le f_2 \le \dots \le f$). \\
Если функция $f$ ограничена, то сходимость можно считать равномерной.
\end{Theorem}

\begin{Proof}
Зафиксируем $n \in \mathbb{N}$. Разобьём ось $Y$ от $0$ до $n$ на отрезки длины $1/n$, а всё, что больше $n$, сгруппируем в один кусок.
Определим полуинтервалы:
$$ \Delta_i^n = \left[ \frac{i}{n}, \frac{i+1}{n} \right), \quad i \in \{0, 1, \dots, n^2-1\} $$
И последний элемент для неограниченного хвоста: $\Delta_{n^2}^n = [n, +\infty)$.
Построим прообразы этих интервалов: $e_i^n = f^{-1}(\Delta_i^n)$. Так как $f$ измерима, все $e_i^n$ измеримы.

Построим простую функцию $g_n(x)$, задав её значения на каждом $e_i^n$ равными нижней границе интервала:
$$ g_n(x) = \frac{i}{n} \quad \text{при } x \in e_i^n $$
\begin{center}
\begin{tikzpicture}[scale=1.5, >=latex]
    \definecolor{darktext}{HTML}{154360}
    \definecolor{colorPoi}{HTML}{7DCEA0}
    \definecolor{colorArea}{HTML}{D4E6F1}

    % Оси
    \draw[->, thick, darktext] (-0.5, 0) -- (6, 0) node[right] {$x$};
    \draw[->, thick, darktext] (0, -0.5) -- (0, 3) node[above] {$f(x)$};

    % Кривая f(x)
    \draw[thick, darktext] plot [smooth, tension=0.7] coordinates { (0, 2.5) (1, 0.5) (2.5, 2.2) (4, 1.2) (5.5, 0.5) };
    \node[darktext] at (2.8, 2.5) {$f(x)$};

    % Разметка по Y
    \draw[dashed, gray] (0, 0.5) node[left, darktext] {$\frac{i}{n}$} -- (5.5, 0.5);
    \draw[dashed, gray] (0, 1.0) node[left, darktext] {$\frac{i+1}{n}$} -- (4.5, 1.0);
    \draw[dashed, gray] (0, 2.0) node[left, darktext] {$n$} -- (2.8, 2.0);

    % Ступеньки g_n(x)
    \draw[ultra thick, colorPoi] (0, 2.0) -- (0.3, 2.0);
    \draw[ultra thick, colorPoi] (0.3, 1.0) -- (0.7, 1.0);
    \draw[ultra thick, colorPoi] (0.7, 0.5) -- (1.6, 0.5);
    \draw[ultra thick, colorPoi] (1.6, 1.0) -- (1.9, 1.0);
    \draw[ultra thick, colorPoi] (1.9, 2.0) -- (3.0, 2.0);
    \draw[ultra thick, colorPoi] (3.0, 1.0) -- (4.2, 1.0);
    \draw[ultra thick, colorPoi] (4.2, 0.5) -- (5.5, 0.5);

    % Пунктиры вниз
    \draw[dashed, colorPoi] (0.7, 0.5) -- (0.7, 0);
    \draw[dashed, colorPoi] (1.6, 0.5) -- (1.6, 0);
    \draw[decorate, decoration={brace, amplitude=4pt, mirror}, darktext] (0.7, -0.1) -- (1.6, -0.1) node[midway, below=4pt] {$e_i^n$};
    
    \node[colorPoi, font=\bfseries] at (4.5, 2.5) {Простая $g_n(x)$};
\end{tikzpicture}
\end{center}

Для любого $x \in e_i^n$ (где $i < n^2$) выполняется оценка:
$$ g_n(x) \le f(x) < g_n(x) + \frac{1}{n} $$
Отсюда очевидно, что $\lim\limits_{n \to \infty} g_n(x) = f(x)$ поточечно. 
Для $x \in e_{n^2}^n$ (хвост): $f(x) \ge n = g_n(x)$ (здесь $f$ может быть неограничена).

Чтобы гарантировать монотонное возрастание приближений, переопределим функцию: 
$$ h_n(x) = \max(g_1(x), \dots, g_n(x)) $$
Она также простая, возрастает и сходится к $f$.

\textbf{Равномерная сходимость:} Если $f \le C$ (ограничена), то при $n > C$ последний "бесконечный" интервал $\Delta_{n^2}^n$ будет пустым. Тогда для всех $x \in X$ выполнено $0 \le f(x) - g_n(x) < \frac{1}{n}$. Это означает равномерную сходимость.
\end{Proof}

\begin{Theorem}{Следствие}{}
Для любой измеримой (в том числе знакопеременной) функции $f$ существует поточечный предел последовательности простых функций $f_n$, такой что $|f_n| \le |f|$.
\end{Theorem}

\begin{Proof}
Разложим функцию на положительную и отрицательную части: $f = f_+ - f_-$, где $f_+ = \max(f, 0)$ и $f_- = \max(-f, 0)$. Применим теорему об аппроксимации к каждой из них по отдельности, а затем вычтем приближения.
\end{Proof}

\section{§ Сходимость почти всюду и по мере}

Пусть $(X, \mathcal{U}, \mu)$ — пространство с мерой.
Обозначим $L^0(X, \mu)$ — пространство измеримых \textbf{почти всюду конечных} функций.
Свойство выполняется \textbf{почти всюду (п.в.)}, если оно выполняется для всех точек $X$, кроме множества меры нуль.

\begin{Example}
Пусть $A$ — множество меры нуль ($\lambda(A) = 0$). Индикаторная функция $\chi_A(x) = \begin{cases} 1, & x \in A \\ 0, & x \notin A \end{cases}$. \\
Можно сказать, что $\chi_A = 0$ почти всюду.
\end{Example}

\begin{Definition}{Сходимость почти всюду}{}
Говорят, что последовательность $f_n \in L^0(X, \mu)$ \textbf{сходится почти всюду} к $f$ (обозначается $f_n \xrightarrow{\text{п.в.}} f$), если существует множество $E \subset X$ меры нуль ($\mu(E) = 0$), такое что для всех $x \in X \setminus E$ выполняется поточечная сходимость: $\lim\limits_{n \to \infty} f_n(x) = f(x)$.
\end{Definition}

\begin{Example}
Пусть $f_n(x) = x^n$ на $[0, 1]$. \\
Поточечный предел: $f(x) = \begin{cases} 1, & x = 1 \\ 0, & x \in [0, 1) \end{cases}$. \\
Так как множество $\{1\}$ имеет меру Лебега 0, то $f_n \xrightarrow{\text{п.в.}} 0$.
\end{Example}

\begin{Definition}{Сходимость по мере}{}
Говорят, что $f_n$ \textbf{сходится к $f$ по мере $\mu$} (обозначается $f_n \xrightarrow{\mu} f$), если для любого $\varepsilon > 0$:
$$ \lim\limits_{n \to \infty} \mu(X(|f_n - f| \ge \varepsilon)) = 0 $$
\end{Definition}

\begin{Example}[ Контрпример ]
Из сходимости почти всюду \textbf{не всегда} следует сходимость по мере.
Рассмотрим $(\mathbb{R}, \mathcal{U}, \lambda)$. Пусть $f_n(x) = \chi_{[n, n+1]}(x)$ (индикатор отрезка $[n, n+1]$).
Для любой фиксированной точки $x \in \mathbb{R}$, при достаточно большом $n$ блок уедет правее, и $f_n(x)$ станет равным 0. Следовательно, $f_n \xrightarrow{\text{п.в.}} 0$.
Однако сходимости по мере нет! Для $\varepsilon = 1/2$:
$$ \lambda(\mathbb{R}(|f_n - 0| \ge 1/2)) = \lambda([n, n+1]) = 1 \not\to 0 $$
\end{Example}

\begin{center}
\begin{tikzpicture}[scale=0.8, >=latex]
    \definecolor{darktext}{HTML}{154360}
    \definecolor{colorHighlight}{HTML}{2874A6}
    
    \draw[->, thick, darktext] (0, 0) -- (10, 0) node[right] {$x$};
    \draw[->, thick, darktext] (0, -0.5) -- (0, 2) node[above] {$f_n(x)$};
    
    \draw[thick, colorHighlight, fill=colorHighlight!20] (2, 0) rectangle (3, 1);
    \node[above, darktext] at (2.5, 1) {$n=1$};
    
    \draw[thick, colorHighlight, fill=colorHighlight!20] (5, 0) rectangle (6, 1);
    \node[above, darktext] at (5.5, 1) {$n=2$};
    
    \draw[thick, colorHighlight, fill=colorHighlight!20] (8, 0) rectangle (9, 1);
    \node[above, darktext] at (8.5, 1) {$n=3 \to \infty$};
    
    \draw[->, ultra thick, darktext, dashed] (3.5, 0.5) -- (4.5, 0.5);
    \draw[->, ultra thick, darktext, dashed] (6.5, 0.5) -- (7.5, 0.5);
\end{tikzpicture}
\end{center}

\begin{Theorem}{Связь сходимостей (Теорема Лебега)}{}
Если пространство имеет конечную меру ($\mu(X) < +\infty$), то из сходимости почти всюду следует сходимость по мере.
\end{Theorem}

\begin{Proof}
Рассмотрим случай, когда $f_n \xrightarrow{X} f$ сходится всюду и монотонно убывая ($f_1 \ge f_2 \ge \dots$).
Зафиксируем $\varepsilon > 0$. Обозначим $X_n = X(|f_n - f| \ge \varepsilon)$.
В силу монотонного убывания $f_n$, множества $X_n$ вложены: $X_1 \supset X_2 \supset \dots$
Так как $f_n \to f$ всюду, ни одна точка не может остаться в $X_n$ навсегда. Значит, $\bigcap\limits_{n=1}^{\infty} X_n = \emptyset$.
Так как мера всего пространства конечна ($\mu(X) < +\infty$), мы можем применить свойство \textbf{непрерывности меры сверху}:
$$ \lim\limits_{n \to \infty} \mu(X_n) = \mu\left( \bigcap\limits_{n=1}^{\infty} X_n \right) = \mu(\emptyset) = 0 $$
Следовательно, $f_n \xrightarrow{\mu} f$. 
\end{Proof}

\begin{Example}[ Контрпример ]
Из сходимости по мере \textbf{не следует} сходимость почти всюду. Рассмотрим классический контрпример на отрезке $X = [0, 1]$ с мерой Лебега. 

Построим последовательность функций-индикаторов, которые «скользят» по отрезку, при этом их ширина постоянно уменьшается. 
Пусть $f_1 = \chi_{[0, 1]}$. 
Затем разделим отрезок пополам: $f_2 = \chi_{[0, 1/2]}$, $f_3 = \chi_{[1/2, 1]}$. 
Затем на три части: $f_4 = \chi_{[0, 1/3]}$, $f_5 = \chi_{[1/3, 2/3]}$, $f_6 = \chi_{[2/3, 1]}$, и так далее.

\begin{center}
\begin{tikzpicture}[scale=1.1, >=latex]
    \definecolor{darktext}{HTML}{154360}
    \definecolor{colorHighlight}{HTML}{2874A6}
    
    % n=1
    \begin{scope}[shift={(0,0)}]
        \draw[->, thick, darktext] (-0.2, 0) -- (2.5, 0) node[right] {$x$};
        \draw[->, thick, darktext] (0, -0.2) -- (0, 1.5) node[above] {$f_n(x)$};
        \draw[darktext] (2, 0) -- (2, -0.1) node[below] {$1$};
        \draw[darktext] (0, 1) -- (-0.1, 1) node[left] {$1$};
        
        \filldraw[fill=colorHighlight!20, draw=colorHighlight, ultra thick] (0, 0) rectangle (2, 1);
        \node[darktext, font=\bfseries] at (1, 1.3) {$n=1$};
    \end{scope}

    % n=2, 3
    \begin{scope}[shift={(3.5,0)}]
        \draw[->, thick, darktext] (-0.2, 0) -- (2.5, 0) node[right] {$x$};
        \draw[->, thick, darktext] (0, -0.2) -- (0, 1.5);
        \draw[darktext] (2, 0) -- (2, -0.1) node[below] {$1$};
        
        \draw[colorHighlight, ultra thick, dashed] (0, 0) rectangle (1, 1);
        \filldraw[fill=colorHighlight!20, draw=colorHighlight, ultra thick] (1, 0) rectangle (2, 1);
        
        \node[darktext, font=\bfseries] at (0.5, 1.3) {$f_2$};
        \node[darktext, font=\bfseries] at (1.5, 1.3) {$f_3$};
        \node[darktext, font=\bfseries, color=colorHighlight] at (1, -0.6) {Ширина $1/2$};
    \end{scope}

    % n=4, 5, 6
    \begin{scope}[shift={(7,0)}]
        \draw[->, thick, darktext] (-0.2, 0) -- (2.5, 0) node[right] {$x$};
        \draw[->, thick, darktext] (0, -0.2) -- (0, 1.5);
        \draw[darktext] (2, 0) -- (2, -0.1) node[below] {$1$};
        
        \draw[colorHighlight, ultra thick, dashed] (0, 0) rectangle (0.66, 1);
        \filldraw[fill=colorHighlight!20, draw=colorHighlight, ultra thick] (0.66, 0) rectangle (1.33, 1);
        \draw[colorHighlight, ultra thick, dashed] (1.33, 0) rectangle (2, 1);
        
        \node[darktext, font=\bfseries] at (0.33, 1.3) {$f_4$};
        \node[darktext, font=\bfseries] at (1.0, 1.3) {$f_5$};
        \node[darktext, font=\bfseries] at (1.66, 1.3) {$f_6$};
        \node[darktext, font=\bfseries, color=colorHighlight] at (1, -0.6) {Ширина $1/3$};
    \end{scope}
\end{tikzpicture}
\end{center}

Заметим, что мера носителя функции $\lambda(X(|f_n| > 0))$ стремится к 0, следовательно, $f_n \xrightarrow{\lambda} 0$ (сходится по мере к нулю).
Однако для \textbf{любой} точки $x \in [0, 1]$ этот «блок» будет проходить через неё бесконечное число раз. Значит, $f_n(x)$ бесконечно часто принимает значение 1 и бесконечно часто значение 0. Поточечного предела не существует нигде!
\end{Example}

\begin{Theorem}{Рисса (О выборе сходящейся подпоследовательности)}{}
Из любой последовательности $f_n$, сходящейся по мере к $f$, можно выбрать подпоследовательность $f_{n_k}$, которая сходится к $f$ почти всюду.
\end{Theorem}

\begin{Proof}
Пусть $f_n \xrightarrow{\mu} f$. По определению сходимости по мере, для любого $k \in \mathbb{N}$ (где $\varepsilon = \frac{1}{k}$):
$$ \lim\limits_{n \to \infty} \mu\left( X\left(|f_n - f| \ge \frac{1}{k}\right) \right) = 0 $$
Следовательно, по определению предела, существует такой номер $n_0(k)$, что для всех $n \ge n_0(k)$ мера этого множества становится меньше $\frac{1}{2^k}$:
$$ \mu\left( X\left(|f_n - f| \ge \frac{1}{k}\right) \right) < \frac{1}{2^k} $$
Будем выбирать возрастающую последовательность индексов $n_k \ge n_0(k)$. 
По лемме Бореля-Кантелли (так как сумма мер $\sum\limits_{k=1}^{\infty} \frac{1}{2^k} < +\infty$), множество точек, которые попадают в эти исключительные множества бесконечно часто, имеет меру нуль. Вне этого множества меры нуль $|f_{n_k} - f| < \frac{1}{k}$ для всех достаточно больших $k$, что означает поточечную сходимость почти всюду.
\end{Proof}

\begin{Theorem}{Следствия из теоремы Рисса}{}
\begin{enumerate}
    \item Если $f_n \xrightarrow{\mu} f$ и $f_n \xrightarrow{\mu} g$, то $f$ и $g$ совпадают почти всюду ($f \stackrel{\text{п.в.}}{=} g$).
    \item Если $f_n \xrightarrow{\mu} f$ и $f_n \le g$ почти всюду, то $f \le g$ почти всюду.
\end{enumerate}
\end{Theorem}

\begin{Proof}
\textbf{1)} По теореме Рисса из последовательности $f_n$ можно выделить подпоследовательность $f_{n_k}$, такую что $f_{n_k} \xrightarrow{\text{п.в.}} f$. 
С другой стороны, любая подпоследовательность сходящейся по мере последовательности сама сходится по мере к тому же пределу. Значит, $f_{n_k} \xrightarrow{\mu} g$. Применим теорему Рисса к $f_{n_k}$: из неё можно выделить подподпоследовательность $f_{n_{k_m}} \xrightarrow{\text{п.в.}} g$. 
Так как $f_{n_{k_m}}$ является подпоследовательностью сходящейся почти всюду к $f$, то она также сходится к $f$. В силу единственности предела числовой последовательности, $f \stackrel{\text{п.в.}}{=} g$.

\textbf{2)} Аналогично, переходя к подпоследовательности, сходящейся почти всюду. Предельный переход сохраняет нестрогие неравенства для числовых последовательностей.
\end{Proof}


\section{§ Интеграл по мере Лебега и его свойства}

\begin{Definition}{Интеграл от простой функции}{}
Пусть задана измеримая \textbf{неотрицательная простая функция} $g$, и пусть $\{A_i\}_{i=1}^n$ — некоторое допустимое разбиение для $g$, где $a_i$ — значения функции $g$ на множестве $A_i$. Тогда интеграл от функции $g$ по мере $\mu$ определяется как:
$$ \int\limits_X g \, d\mu = \sum\limits_{i=1}^n a_i \cdot \mu(A_i) $$
\end{Definition}

\begin{Example}[ Интеграл от функции Дирихле]
Рассмотрим пространство $X = \mathbb{R}$ с классической мерой Лебега $\lambda$.
Функция Дирихле $d(x) = \begin{cases} 1, & x \in \mathbb{Q} \\ 0, & x \notin \mathbb{Q} \end{cases}$.
Она является простой. Допустимое разбиение: $A_1 = \mathbb{Q}$ (значение $a_1 = 1$) и $A_2 = \mathbb{R} \setminus \mathbb{Q}$ (значение $a_2 = 0$).
$$ \int\limits_{\mathbb{R}} d(x) \, d\lambda = 1 \cdot \lambda(\mathbb{Q}) + 0 \cdot \lambda(\mathbb{R} \setminus \mathbb{Q}) = 1 \cdot 0 + 0 \cdot (+\infty) = 0 $$
\end{Example}

\begin{Lemma}{Корректность определения}{}
Определение интеграла от простой функции не зависит от выбора допустимого разбиения.
\end{Lemma}

\begin{Proof}
Пусть $\{A_i\}_{i=1}^n$ и $\{B_j\}_{j=1}^m$ — два различных допустимых разбиения для функции $g$ со значениями $a_i$ и $b_j$ соответственно.
Заметим, что система попарных пересечений $\{A_i \cap B_j\}$ является их общим измельчением. Каждое $A_i$ представимо как дизъюнктное объединение $A_i = \bigsqcup\limits_{j=1}^m (A_i \cap B_j)$.
Запишем сумму для первого разбиения, используя конечную аддитивность меры:
$$ \sum\limits_{i=1}^n a_i \mu(A_i) = \sum\limits_{i=1}^n a_i \left( \sum\limits_{j=1}^m \mu(A_i \cap B_j) \right) = \sum\limits_{i=1}^n \sum\limits_{j=1}^m a_i \mu(A_i \cap B_j) $$
Так как пересечение $A_i \cap B_j$ не пусто только если функция принимает на нём одно и то же значение, то на этом пересечении $a_i = b_j$. Заменим $a_i$ на $b_j$:
$$ \dots = \sum\limits_{i=1}^n \sum\limits_{j=1}^m b_j \mu(A_i \cap B_j) = \sum\limits_{j=1}^m b_j \left( \sum\limits_{i=1}^n \mu(A_i \cap B_j) \right) = \sum\limits_{j=1}^m b_j \mu(B_j) $$
Суммы совпали, следовательно, значение интеграла корректно.
\end{Proof}

\begin{Lemma}{Свойства интеграла для простых функций}{}
\begin{enumerate}
    \item Если $g(x) = C$ на измеримом множестве $E$, то $\int\limits_E g \, d\mu = C \cdot \mu(E)$.
    \item Если $g_1, g_2$ — простые функции и $g_1 \le g_2$, то $\int\limits_X g_1 \, d\mu \le \int\limits_X g_2 \, d\mu$.
\end{enumerate}
\end{Lemma}

\begin{Proof}
\textbf{1)} Допустимым разбиением для $X$ является множество $E$ (где $g=C$) и $X \setminus E$ (доопределим там функцию нулём). 
Тогда $\int\limits_X \hat{g} \, d\mu = 0 \cdot \mu(X \setminus E) + C \cdot \mu(E) = C \cdot \mu(E)$.

\textbf{2)} По лемме об общем допустимом разбиении, существует разбиение $\{A_i\}$, на котором обе функции постоянны. Пусть их значения равны $a_i$ и $b_i$. Из условия $g_1 \le g_2$ следует, что $\forall i: a_i \le b_i$.
$$ \int\limits_X g_1 \, d\mu = \sum\limits_{i} a_i \mu(A_i) \le \sum\limits_{i} b_i \mu(A_i) = \int\limits_X g_2 \, d\mu $$
\end{Proof}

\begin{Definition}{Интеграл от измеримой функции}{}
Пусть $f$ — неотрицательная измеримая функция. Её \textbf{интеграл Лебега} определяется как супремум интегралов от всех простых функций, не превосходящих $f$:
$$ \int\limits_X f \, d\mu = \sup \left\{ \int\limits_X g \, d\mu \;\middle|\; g \text{ — простая, } g \le f \right\} $$

Если $f$ — произвольная знакопеременная измеримая функция, её представляют в виде разности $f = f_+ - f_-$, где $f_+ = \max(0, f)$ и $f_- = \max(0, -f)$. Интеграл равен:
$$ \int\limits_X f \, d\mu = \int\limits_X f_+ \, d\mu - \int\limits_X f_- \, d\mu $$
(при условии, что эта разность определена в $\overline{\mathbb{R}}$, то есть не возникает неопределенности $\infty - \infty$).
\end{Definition}

\begin{Definition}{Суммируемость}{}
Если оба интеграла $\int\limits_X f_+ \, d\mu$ и $\int\limits_X f_- \, d\mu$ \textbf{конечны}, то функция $f$ называется \textbf{суммируемой} по мере $\mu$. Пространство суммируемых функций обозначается $L^1(X, \mu)$.
(Если конечен только один из них, говорят, что функция интегрируема, но её интеграл равен $+\infty$ или $-\infty$).
\end{Definition}

\begin{Example}
Рассмотрим функцию $\frac{\sin x}{x}$ на $[0, +\infty)$. 
В смысле несобственного интеграла Римана: $\int\limits_0^{+\infty} \frac{\sin x}{x} dx = \frac{\pi}{2}$ (сходится условно).
Однако интеграла Лебега $\int\limits_{[0, +\infty)} \frac{\sin x}{x} \, d\lambda$ \textbf{не существует}, так как площади и положительной ($f_+$), и отрицательной ($f_-$) частей бесконечны, что даёт неопределённость $\infty - \infty$. (По Лебегу функция интегрируема тогда и только тогда, когда интегрируем её модуль).
\end{Example}

\begin{Theorem}{Свойства интеграла Лебега}{}
Пусть $f, g$ — неотрицательные измеримые функции.
\begin{enumerate}
    \item \textbf{Монотонность:} Если $f \le g$ на множестве $E$, то $\int\limits_E f \, d\mu \le \int\limits_E g \, d\mu$.
    \item \textbf{Интегрирование по множеству меры нуль:} Если $\mu(E) = 0$, то $\int\limits_E f \, d\mu = 0$.
\end{enumerate}
\end{Theorem}

\begin{Proof}
\textbf{1)} По определению, интеграл — это супремум по простым функциям. Если простая функция $h \le f$, то в силу $f \le g$, она также $h \le g$. Следовательно, множество простых функций, аппроксимирующих $f$, является подмножеством простых функций, аппроксимирующих $g$. Супремум по меньшему множеству не превосходит супремума по большему.

\textbf{2)} Если функция $f$ ограничена ($f \le C$), то по доказанному ранее для простых функций (свойство 1):
$$ \int\limits_E |f| \, d\mu \le \int\limits_E C \, d\mu = C \cdot \mu(E) = C \cdot 0 = 0 $$
Если функция неограничена, то любая простая функция $h \le f$ всё равно принимает лишь конечное число значений, а значит она ограничена. Тогда $\int\limits_E h \, d\mu = 0$. Так как супремум множества нулей равен нулю, интеграл равен нулю.
\end{Proof}

\section{§ Предельный переход под знаком интеграла}

\begin{Theorem}{Леви (О монотонной сходимости)}{}
Пусть $f_n$ — последовательность \textbf{неотрицательных} измеримых функций на $X$, которая монотонно возрастает ($f_{n+1} \ge f_n$) и поточечно сходится к функции $f$. 
Тогда:
$$ \lim\limits_{n \to \infty} \int\limits_X f_n \, d\mu = \int\limits_X f \, d\mu $$
\end{Theorem}

\begin{Proof}
В силу монотонности функций ($f_n \le f_{n+1} \le f$), их интегралы образуют неубывающую числовую последовательность, ограниченную сверху числом $\int\limits_X f \, d\mu$. Значит, существует предел:
$$ L = \lim\limits_{n \to \infty} \int\limits_X f_n \, d\mu \le \int\limits_X f \, d\mu $$
(где $L \in \overline{\mathbb{R}}$). Осталось доказать обратное неравенство: $L \ge \int\limits_X f \, d\mu$.

Возьмём произвольную простую функцию $g$, такую что $0 \le g \le f$. 
Зафиксируем число $\theta \in (0, 1)$. Рассмотрим множества:
$$ X_n = X(f_n \ge \theta g) = \{ x \in X : f_n(x) \ge \theta g(x) \} $$
Так как $f_n$ возрастает, множества вложены: $X_1 \subset X_2 \subset \dots \subset X_n \subset X_{n+1}$. \\
Покажем, что $\bigcup\limits_{n=1}^{\infty} X_n = X$. Действительно, для любого $x \in X$, если $f(x) = 0$, то $x \in X_1$. Если $f(x) > 0$, то $\theta g(x) < g(x) \le f(x) = \lim\limits_{n \to \infty} f_n(x)$. Значит, начиная с некоторого номера $N$, значение $f_N(x)$ превысит $\theta g(x)$, и точка попадёт в $X_N$.

Используя монотонность интеграла:
$$ \int\limits_X f_n \, d\mu \ge \int\limits_{X_n} f_n \, d\mu \ge \int\limits_{X_n} \theta g \, d\mu = \theta \sum\limits_{i=1}^k a_i \mu(A_i \cap X_n) $$
где $A_i$ — разбиение для простой функции $g$ со значениями $a_i$. \\
Устремим $n \to \infty$. В левой части получим предел $L$. В правой части, в силу \textbf{непрерывности меры снизу} (так как множества $A_i \cap X_n$ расширяются до $A_i$), мера $\mu(A_i \cap X_n) \to \mu(A_i)$.
$$ L \ge \theta \sum\limits_{i=1}^k a_i \mu(A_i) = \theta \int\limits_X g \, d\mu $$
Это неравенство верно для любого $\theta \in (0, 1)$. Переходя к пределу при $\theta \to 1-0$, получаем $L \ge \int\limits_X g \, d\mu$. \\
Наконец, так как неравенство верно для любой простой функции $g \le f$, возьмём супремум по всем таким $g$. По определению интеграла Лебега получим $L \ge \int\limits_X f \, d\mu$. Теорема доказана.
\end{Proof}

\begin{Theorem}{Линейность интеграла Лебега}{}
Пусть $f, g$ — суммируемые функции ($f, g \in L^1(X, \mu)$), а $\alpha, \beta \in \mathbb{R}$. Тогда:
$$ \int\limits_X (\alpha f + \beta g) \, d\mu = \alpha \int\limits_X f \, d\mu + \beta \int\limits_X g \, d\mu $$
\end{Theorem}

\begin{Proof}
\textit{(Развёрнутое доказательство)} \\
Сначала докажем для случая $\alpha, \beta \ge 0$ и неотрицательных функций $f, g \ge 0$.
Если $f$ и $g$ — \textbf{простые функции}, то у них есть общее допустимое разбиение $\{A_i\}$ со значениями $a_i$ и $b_i$. Комбинация $\alpha f + \beta g$ также является простой функцией со значениями $\alpha a_i + \beta b_i$. По определению:
$$ \int\limits_X (\alpha f + \beta g) \, d\mu = \sum\limits_{i} (\alpha a_i + \beta b_i) \mu(A_i) = \alpha \sum\limits_{i} a_i \mu(A_i) + \beta \sum\limits_{i} b_i \mu(A_i) = \alpha \int\limits_X f \, d\mu + \beta \int\limits_X g \, d\mu $$

Перейдем к \textbf{произвольным неотрицательным измеримым} функциям $f$ и $g$. По теореме об аппроксимации, существуют возрастающие последовательности простых функций $f_n \nearrow f$ и $g_n \nearrow g$. 
Тогда линейная комбинация $h_n = \alpha f_n + \beta g_n$ является возрастающей последовательностью простых функций, которая поточечно сходится к $\alpha f + \beta g$. \\
Применим теорему Леви (о монотонной сходимости) к последовательности $h_n$:
\begin{align*}
    \int\limits_X (\alpha f + \beta g) \, d\mu &= \lim\limits_{n \to \infty} \int\limits_X (\alpha f_n + \beta g_n) \, d\mu \\
    &= \lim\limits_{n \to \infty} \left( \alpha \int\limits_X f_n \, d\mu + \beta \int\limits_X g_n \, d\mu \right) \\
    &= \alpha \lim\limits_{n \to \infty} \int\limits_X f_n \, d\mu + \beta \lim\limits_{n \to \infty} \int\limits_X g_n \, d\mu \\
    &= \alpha \int\limits_X f \, d\mu + \beta \int\limits_X g \, d\mu
\end{align*}
Для знакопеременных функций утверждение доказывается разложением $f = f_+ - f_-$ и $g = g_+ - g_-$ и группировкой положительных слагаемых.
\end{Proof}
\begin{Lemma}{Счётная аддитивность интеграла Лебега}{}
Пусть $f \ge 0$ — измеримая функция, и $A = \bigsqcup\limits_{i=1}^{\infty} A_i$ (дизъюнктное объединение измеримых множеств). Тогда:
$$ \int\limits_A f \, d\mu = \sum\limits_{i=1}^{\infty} \int\limits_{A_i} f \, d\mu $$
\end{Lemma}

\begin{Proof}
Введём последовательность функций $f_n(x) = \sum\limits_{i=1}^n f(x) \cdot I_{A_i}(x)$, где $I_{A_i}$ — индикатор множества $A_i$.
Заметим, что:
1) $f_n(x)$ монотонно возрастает (так как $f \ge 0$ и добавляются новые неотрицательные слагаемые).
2) $f_n(x)$ поточечно сходится к $f(x) \cdot I_A(x)$.
По теореме Леви о монотонной сходимости предел интегралов равен интегралу от предела:
$$ \int\limits_A f \, d\mu = \lim\limits_{n \to \infty} \int\limits_X f_n \, d\mu = \lim\limits_{n \to \infty} \sum\limits_{i=1}^n \int\limits_{A_i} f \, d\mu = \sum\limits_{i=1}^{\infty} \int\limits_{A_i} f \, d\mu $$
\end{Proof}

\noindent \textbf{NB:} Пусть $f \ge 0$, задано пространство $(\mathbb{R}, \mathcal{U}, \lambda)$. Функция множества $\mu(A) = \int\limits_A f \, d\lambda$ является мерой для любого $A \in \mathcal{U}$. При этом функция $f$ называется \textbf{плотностью} меры $\mu$.

\begin{Lemma}{О хвостах интегрируемой функции}{}
Пусть $f \in L^1(X, \mu)$ и мера пространства бесконечна ($\mu(X) = +\infty$). Тогда для любого $\varepsilon > 0$ существует измеримое множество $A$, такое что:
$$ \mu(A) < +\infty \quad \text{и} \quad \int\limits_{X \setminus A} |f| \, d\mu < \varepsilon $$
\end{Lemma}

\begin{Proof}
Рассмотрим последовательность множеств $A_n = \left\{ x \in X : |f(x)| < \frac{1}{n} \right\}$. 
Очевидно, что эти множества вложены ($A_1 \supset A_2 \supset A_3 \dots$), то есть последовательность убывает. Пересечение всех этих множеств: $\bigcap\limits_{n=1}^{\infty} A_n = \{ x \in X : f(x) = 0 \}$.
По теореме о предельном переходе (в силу интегрируемости $|f|$):
$$ \lim\limits_{n \to \infty} \int\limits_{A_n} |f| \, d\mu = \int\limits_{\{f=0\}} |f| \, d\mu = 0 $$
По определению предела, для заданного $\varepsilon > 0$ существует номер $n_0$, такой что $\int\limits_{A_{n_0}} |f| \, d\mu < \varepsilon$.
Положим искомое множество $A = X \setminus A_{n_0} = \left\{ x \in X : |f(x)| \ge \frac{1}{n_0} \right\}$.
Тогда $\int\limits_{X \setminus A} |f| \, d\mu = \int\limits_{A_{n_0}} |f| \, d\mu < \varepsilon$.
Проверим конечность меры $A$:
$$ \int\limits_X |f| \, d\mu \ge \int\limits_A |f| \, d\mu \ge \int\limits_A \frac{1}{n_0} \, d\mu = \frac{1}{n_0} \mu(A) $$
Так как интеграл от $|f|$ конечен ($f \in L^1$), то $\mu(A) \le n_0 \int\limits_X |f| \, d\mu < +\infty$. Лемма доказана.
\end{Proof}

\begin{Example}[ О «дельта-функции»]
В классическом анализе часто пишут, что $\delta(x) = \begin{cases} +\infty, & x=0 \\ 0, & x \ne 0 \end{cases}$. 
В рамках теории меры Лебега это \textbf{совершенно неверно}, так как функция, равная нулю почти везде, имеет интеграл Лебега равный $0$, а не $1$.

В реальности дельта-функция Дирака — это обобщённая функция, которая порождается дельта-образной последовательностью функций $\varphi_\delta(x)$, сходящейся к ней в смысле интегралов:
\begin{center}
\begin{tikzpicture}[scale=1, >=latex]
    \definecolor{darktext}{HTML}{154360}
    \definecolor{colorHighlight}{HTML}{2874A6}
    
    \draw[->, thick, darktext] (-3, 0) -- (3, 0) node[right] {$x$};
    \draw[->, thick, darktext] (0, -0.5) -- (0, 3) node[above] {$y$};
    \draw[colorHighlight, ultra thick, fill=colorHighlight!20] (-1, 0) rectangle (1, 2);
    
    \draw[darktext] (-1, 0) -- (-1, -0.1) node[below] {$-\delta$};
    \draw[darktext] (1, 0) -- (1, -0.1) node[below] {$\delta$};
    \draw[darktext] (0, 2) -- (-0.1, 2) node[left] {$\frac{1}{2\delta}$};
    
    \node[darktext, font=\bfseries] at (2, 2) {$\varphi_\delta(x)$};
\end{tikzpicture}
\end{center}
$$ \varphi_\delta(x) = \begin{cases} \frac{1}{2\delta}, & x \in [-\delta, \delta] \\ 0, & \text{иначе} \end{cases} \implies \int\limits_{\mathbb{R}} \varphi_\delta \, d\lambda = 1 $$
Тогда $\int\limits_{\mathbb{R}} \varphi_\delta(x) \psi(x-x_0) \, d\lambda \xrightarrow{\delta \to 0} \psi(x_0)$. 

Связь с дискретной вероятностью: 
$$ P(x \in A) = \sum\limits_{i: a_i \in A} p_i = \int\limits_A \sum\limits_{i} p_i \delta(x - a_i) \, d\lambda $$
\end{Example}

\begin{Lemma}{Абсолютная непрерывность интеграла Лебега}{}
Пусть $f \in L^1(X, \mu)$. Тогда для любого $\varepsilon > 0$ существует такое $\delta > 0$, что для любого измеримого множества $e \subset X$, мера которого $\mu(e) < \delta$, выполняется неравенство:
$$ \int\limits_e |f| \, d\mu < \varepsilon $$
\end{Lemma}

\begin{Proof}
Зафиксируем $\varepsilon > 0$. По теореме об аппроксимации существует простая неотрицательная функция $g_\varepsilon$, такая что $0 \le g_\varepsilon \le |f|$ и:
$$ \int\limits_X |f| \, d\mu < \int\limits_X g_\varepsilon \, d\mu + \varepsilon $$
Интеграл по множеству $e$ можно переписать и оценить:
$$ \int\limits_e |f| \, d\mu = \int\limits_e (|f| - g_\varepsilon) \, d\mu + \int\limits_e g_\varepsilon \, d\mu \le \int\limits_X (|f| - g_\varepsilon) \, d\mu + \int\limits_e g_\varepsilon \, d\mu $$
Первый интеграл строго меньше $\varepsilon$. 
Так как $g_\varepsilon$ — простая функция, она ограничена некоторой константой $C_\varepsilon$. Тогда второй интеграл: $\int\limits_e g_\varepsilon \, d\mu \le \int\limits_e C_\varepsilon \, d\mu = C_\varepsilon \cdot \mu(e)$.
Потребуем, чтобы $C_\varepsilon \cdot \mu(e) < \varepsilon$. Отсюда находим $\delta = \frac{\varepsilon}{C_\varepsilon}$.
Тогда для любого $e$ с мерой $\mu(e) < \delta$ будет выполнено требуемое неравенство.
\end{Proof}

\begin{Theorem}{О равномерной сходимости}{}
Пусть $f_n \in L^1(X, \mu)$, последовательность равномерно сходится к функции $f$ ($f_n \xrightarrow{\rightrightarrows} f$ на $X$), и мера пространства конечна ($\mu(X) < +\infty$). Тогда $f \in L^1(X, \mu)$ и:
$$ \lim\limits_{n \to \infty} \int\limits_X f_n \, d\mu = \int\limits_X f \, d\mu $$
\end{Theorem}

\begin{Proof}
Из равномерной сходимости следует, что существует такой номер $n_0$, что для всех $n > n_0$ и для всех $x \in X$ выполнено $|f_n(x) - f(x)| < 1$. \\
Проинтегрируем это неравенство: $\int\limits_X |f_n - f| \, d\mu \le \int\limits_X 1 \, d\mu = \mu(X) < +\infty$.
Значит, разность $|f_n - f| \in L^1(X, \mu)$. 
Представим $|f| = |f - f_n + f_n| \le |f - f_n| + |f_n|$. Так как оба слагаемых суммируемы, то $f \in L^1(X, \mu)$.
Теперь оценим разность интегралов:
$$ \left| \int\limits_X f \, d\mu - \int\limits_X f_n \, d\mu \right| \le \int\limits_X |f - f_n| \, d\mu \le \int\limits_X \left( \sup\limits_{x \in X} |f(x) - f_n(x)| \right) \, d\mu = \sup\limits_{x \in X} |f - f_n| \cdot \mu(X) $$
Так как сходимость равномерная, супремум стремится к $0$ при $n \to \infty$. Значит, предел разности интегралов равен нулю.
\end{Proof}

\begin{Theorem}{Леви (Вторая формулировка)}{}
Пусть $f_n$ — последовательность неотрицательных измеримых функций, и для почти всех $x \in X$ выполнено: $f_n(x) \xrightarrow{\text{п.в.}} f(x)$ и $f_n(x) \le f_{n+1}(x)$. Тогда:
$$ \lim\limits_{n \to \infty} \int\limits_X f_n \, d\mu = \int\limits_X f \, d\mu $$
\end{Theorem}

\begin{Theorem}{Лебега о мажорируемой сходимости}{}
Пусть $f_n$ — последовательность измеримых функций, $f_n \xrightarrow{\mu} f$ (сходится по мере). Если:
\begin{enumerate}
    \item Для всех $n$ почти всюду выполнено $|f_n| \le g$.
    \item Существует мажоранта $g \in L^1(X, \mu)$.
\end{enumerate}
Тогда $f_n$ и предельная функция $f \in L^1(X, \mu)$, и справедлив предельный переход:
$$ \lim\limits_{n \to \infty} \int\limits_X f_n \, d\mu = \int\limits_X f \, d\mu $$
\end{Theorem}

\begin{Proof}
\textbf{1) Суммируемость предела:} Переходя к пределу, получаем $|f| \le g$ п.в. Поскольку мажоранта $g$ суммируема, функция $f$ также суммируема.

\textbf{2) Случай конечной меры $\mu(X) < +\infty$:} 
Зададим $\varepsilon > 0$. Введём множество $X_n(\varepsilon) = \left\{ x \in X : |f(x) - f_n(x)| \ge \varepsilon \right\}$.
Оценим модуль разности интегралов:
\begin{align*}
    \left| \int\limits_X f \, d\mu - \int\limits_X f_n \, d\mu \right| &\le \int\limits_X |f - f_n| \, d\mu = \int\limits_{X_n(\varepsilon)} |f - f_n| \, d\mu + \int\limits_{X \setminus X_n(\varepsilon)} |f - f_n| \, d\mu \\
    &\le \int\limits_{X_n(\varepsilon)} 2g \, d\mu + \int\limits_{X \setminus X_n(\varepsilon)} \varepsilon \, d\mu \le \int\limits_{X_n(\varepsilon)} 2g \, d\mu + \varepsilon \cdot \mu(X)
\end{align*}
Так как $f_n \xrightarrow{\mu} f$, мера $\mu(X_n(\varepsilon)) \to 0$ при $n \to \infty$. 
В силу абсолютной непрерывности интеграла Лебега, первый интеграл $\int\limits_{X_n(\varepsilon)} 2g \, d\mu \xrightarrow{n \to \infty} 0$. Второе слагаемое $\varepsilon \cdot \mu(X)$ может быть сделано сколь угодно малым. Следовательно, предел равен нулю.

\textbf{3) Случай бесконечной меры $\mu(X) = +\infty$:} 
Так как мажоранта $g \in L^1$, по лемме "о хвостах", для любого $\varepsilon > 0$ существует множество $A$ конечной меры ($\mu(A) < +\infty$), такое что $\int\limits_{X \setminus A} 2g \, d\mu < \varepsilon$.
Тогда:
$$ \int\limits_X |f - f_n| \, d\mu \le \int\limits_A |f - f_n| \, d\mu + \int\limits_{X \setminus A} 2g \, d\mu < \int\limits_A |f - f_n| \, d\mu + \varepsilon $$
Первый интеграл берётся по множеству $A$ конечной меры, и по пункту (2) он стремится к нулю. Устремляя $\varepsilon \to 0$, получаем требуемое равенство.
\end{Proof}

\noindent \textbf{NB:} В теореме Лебега сходимость по мере может быть заменена на сходимость почти везде ($f_n \xrightarrow{\text{п.в.}} f$).

\begin{Definition}{Гамма-функция Эйлера}{}
Классический пример интеграла:
$$ \Gamma(\alpha) = \int\limits_0^{+\infty} x^{\alpha - 1} e^{-x} \, dx $$
Свойства:
\begin{enumerate}
    \item Для $\alpha > 0$ интеграл сходится (существует $\Gamma(\alpha)$).
    \item Формула понижения: $\Gamma(\alpha + 1) = \alpha \Gamma(\alpha)$ при $\alpha > 0$.
    \item Обобщение факториала: $\Gamma(n + 1) = n! \cdot \Gamma(1) = n!$.
    \item Значение: $\Gamma\left(\frac{1}{2}\right) = \sqrt{\pi}$.
\end{enumerate}
\end{Definition}

\begin{Theorem}{Фату (Лемма Фату)}{}
Пусть $f_n$ — последовательность неотрицательных измеримых функций, и $f_n \xrightarrow{\text{п.в.}} f$. Если $\exists C > 0$ такое, что $\int\limits_X f_n \, d\mu \le C$ для всех $n$, то:
$$ \int\limits_X f \, d\mu \le C $$
\end{Theorem}

\begin{Proof}
Определим новую последовательность функций: $g_n(x) = \inf\{f_n, f_{n+1}, \dots\}$. 
Понятно, что:
1) $g_{n+1} \ge g_n$ (последовательность монотонно возрастает).
2) $g_n \xrightarrow{\text{п.в.}} f$.
По теореме Леви к последовательности $g_n$:
$$ \lim\limits_{n \to \infty} \int\limits_X g_n \, d\mu = \int\limits_X f \, d\mu $$
Итого, по построению $g_n \le f_n$, значит:
$$ \int\limits_X g_n \, d\mu \le \int\limits_X f_n \, d\mu \le C $$
Переходя к пределу, получаем $\int\limits_X f \, d\mu \le C$.
\end{Proof}

\section{§ Интегралы, зависящие от параметра}

Рассмотрим интеграл вида:
$$ I(y) = \int\limits_X f(x, y) \, d\mu(x) $$
где $f : X \times Y \to \mathbb{R}$. \\
\textbf{Основные вопросы для исследования $I(y)$:}
\begin{enumerate}
    \item Вычисление предела $\lim\limits_{y \to y_0} I(y)$? (В частности, непрерывна ли функция $I(y)$?)
    \item Существует ли производная $I'(y)$ и можно ли дифференцировать под знаком интеграла?
    \item Существует ли повторный интеграл: $\int\limits_Y I(y) \, d\nu(y) = \int\limits_Y \left( \int\limits_X f(x,y) \, d\mu(x) \right) \, d\nu(y)$? (Решается теоремой Фубини / Тонелли).
\end{enumerate}

\begin{Definition}{Интегралы Эйлера (Гамма и Бета функции)}{}
Классические примеры интегралов, зависящих от параметра:
\begin{itemize}
    \item \textbf{Интеграл Эйлера 2-го рода (Гамма-функция):} 
    $$ \Gamma(\alpha) = \int\limits_0^{+\infty} x^{\alpha - 1} e^{-x} \, dx $$
    \item \textbf{Интеграл Эйлера 1-го рода (Бета-функция):} 
    $$ B(\alpha, \beta) = \int\limits_0^1 x^{\alpha-1} (1-x)^{\beta-1} \, dx, \quad \text{при } \alpha > 0, \beta > 0 $$
\end{itemize}
Функции тесно связаны между собой фундаментальным соотношением:
$$ B(\alpha, \beta) = \frac{\Gamma(\alpha)\Gamma(\beta)}{\Gamma(\alpha+\beta)} $$
\end{Definition}


\subsection*{1. Равномерная сходимость семейства функций}

\begin{Definition}{Равномерная сходимость по параметру}{}
Пусть задана функция $f: X \times Y \to \mathbb{R}$, и $x_0$ — предельная (или внутренняя) точка $X$.
Говорят, что $f(x, y)$ \textbf{равномерно (по $y$)} сходится к $g(y)$ при $x \to x_0$, если:
$$ \forall \varepsilon > 0 \quad \exists \delta > 0 \quad \forall x \in X: 0 < |x - x_0| < \delta \quad \forall y \in Y \implies |f(x,y) - g(y)| < \varepsilon $$
Обозначение: $f(x,y) \rightrightarrows_{x \to x_0} g(y)$.
\end{Definition}

\begin{Example}[ Равномерной и неравномерной сходимости]
Рассмотрим предел при $x \to +\infty$ для $y \in \mathbb{R}$:
\begin{enumerate}
    \item $f(x, y) = \frac{\sin(xy)}{1+x}$. Поточечный предел $\lim\limits_{x \to \infty} f(x, y) = 0$. Сходимость \textbf{равномерная} по $y$, так как мы можем оценить модуль функцией, не зависящей от $y$:
    $$ \left| \frac{\sin(xy)}{1+x} \right| \le \frac{1}{1+x} \xrightarrow{x \to \infty} 0 $$
    \item $f(x, y) = \frac{\sin(xy)}{1+xy}$. Поточечно при $x \to +\infty$ предел также равен $0$. Однако сходимость \textbf{не является равномерной}. 
    Действительно, если двигаться по траектории $y = \frac{1}{x}$ (при $x \to \infty$), то $f\left(x, \frac{1}{x}\right) = \frac{\sin(1)}{2} \ne 0$, что нарушает определение равномерности.
\end{enumerate}
\end{Example}

\begin{Theorem}{О равномерном пределе под знаком интеграла}{}
Пусть $f: X \times Y \to \mathbb{R}$ равномерно сходится к $g(x)$ при $y \to y_0$ ($f \underset{y\to y_0}{\rightrightarrows} g$). Если $f(x, y)$ суммируема по $X$ для всех $y \in Y$, и мера пространства конечна ($\mu(X) < +\infty$), то:
$$ \lim\limits_{y \to y_0} \int\limits_X f(x,y) \, d\mu(x) = \int\limits_X g(x) \, d\mu(x) $$
\end{Theorem}

\begin{Proof}
По определению равномерной сходимости, существует такое $\delta > 0$, что для всех $y$, удовлетворяющих $0 < |y - y_0| < \delta$, и для всех $x \in X$ выполнено: $|f(x,y) - g(x)| < 1$. \\
Оценим модуль предельной функции $g(x)$:
$$ |g(x)| = |g(x) - f(x,y) + f(x,y)| \le |g(x) - f(x,y)| + |f(x,y)| \le 1 + |f(x,y)| $$
Так как мера $\mu(X)$ конечна, константа $1$ интегрируема. Функция $f(x,y)$ суммируема по условию. Значит, мажоранта суммируема, следовательно, $g(x)$ — суммируемая функция.

Оценим разность интегралов:
\begin{align*}
    \left| \int\limits_X f(x,y) \, d\mu(x) - \int\limits_X g(x) \, d\mu(x) \right| &\le \int\limits_X |f(x,y) - g(x)| \, d\mu(x) \\
    &\le \sup\limits_{x \in X} |f(x,y) - g(x)| \cdot \mu(X)
\end{align*}
Поскольку сходимость равномерная, супремум разности стремится к $0$ при $y \to y_0$. Так как $\mu(X) < +\infty$, вся правая часть стремится к нулю. Теорема доказана.
\end{Proof}

\begin{Theorem}{Следствие (Непрерывность интеграла)}{}
Если в условиях предыдущей теоремы функция $f(x, y)$ непрерывна в точке $y_0$ для любого $x \in X$, то интеграл $I(y)$ непрерывен в точке $y_0$.
\end{Theorem}

\begin{Proof}
Из непрерывности $f$ по $y$ следует, что $\lim\limits_{y \to y_0} f(x, y) = f(x, y_0)$. Применяя доказанную теорему:
$$ \lim\limits_{y \to y_0} I(y) = \lim\limits_{y \to y_0} \int\limits_X f(x,y) \, d\mu(x) = \int\limits_X f(x, y_0) \, d\mu(x) = I(y_0) $$
Что и означает непрерывность функции $I(y)$.
\end{Proof}

\subsection*{2. Случай суммируемой мажоранты}

Требование конечности меры $\mu(X)$ и строгой равномерной сходимости часто бывает слишком жёстким. Мощной альтернативой выступает применение теоремы Лебега.

\begin{Theorem}{О предельном переходе (через суммируемую мажоранту)}{}
Пусть $f: X \times Y \to \mathbb{R}$, и существует поточечный предел $f(x,y) \xrightarrow{y \to y_0} g(x)$ при почти всех $x \in X$. \\
Если существуют окрестность $\mathring{U}(y_0)$ и функция $h(x)$, такие что:
\begin{enumerate}
    \item $|f(x,y)| \le h(x)$ для всех $y \in \mathring{U}(y_0)$ и для почти всех $x \in X$.
    \item Мажоранта $h(x)$ суммируема ($h \in L^1(X, \mu)$).
\end{enumerate}
Тогда предельная функция $g(x)$ суммируема, и возможен предельный переход:
$$ \lim\limits_{y \to y_0} \int\limits_X f(x,y) \, d\mu(x) = \int\limits_X g(x) \, d\mu(x) $$
\end{Theorem}

\begin{Proof}
\textit{(Упражнение с лекции)} \\
Докажем по определению предела по Гейне. Возьмём произвольную последовательность $y_n \to y_0$ (где $y_n \in \mathring{U}(y_0)$). Определим функциональную последовательность $f_n(x) = f(x, y_n)$.
По условию, $f_n(x) \xrightarrow{\text{п.в.}} g(x)$.
Кроме того, для всех $n$ почти всюду выполнено $|f_n(x)| \le h(x) \in L^1$.
Это в точности условия классической \textbf{теоремы Лебега о мажорируемой сходимости} для последовательностей. Применяя её, получаем:
$$ \lim\limits_{n \to \infty} \int\limits_X f_n(x) \, d\mu(x) = \int\limits_X g(x) \, d\mu(x) $$
Так как это верно для любой последовательности $y_n \to y_0$, то предел функции равен интегралу от предела.
\end{Proof}

\begin{Theorem}{Следствие}{}
Если функция $f(x,y) \xrightarrow{y \to y_0} f(x, y_0)$ при почти всех $x$ (то есть непрерывна по $y$), ограничена ($|f| \le C$) и мера пространства конечна ($\mu(X) < +\infty$), то $I(y)$ непрерывна в $y_0$.
\end{Theorem}

\begin{Proof}
Если $|f| \le C$ и $\mu(X) < +\infty$, то константа $C$ является суммируемой мажорантой. Применяем предыдущую теорему о предельном переходе.
\end{Proof}

\noindent \textbf{NB (Связь с числовыми рядами):} Рассмотрим пространство $(\mathbb{R}, 2^{\mathbb{R}}, \#)$, где $\#(A) = \text{card}(A \cap \mathbb{N})$ — считающая мера. 
Тогда обычная сумма ряда — это просто интеграл Лебега по считающей мере: 
$$ \sum\limits_{n=1}^{\infty} a_n = \int\limits_{\mathbb{R}} f \, d\#, \quad \text{где } f(n) = a_n $$
В этом контексте знаменитый \textbf{признак Вейерштрасса} равномерной сходимости рядов (наличие сходящегося мажорирующего числового ряда) является прямым следствием \textbf{теоремы Лебега о мажорируемой сходимости}!

\begin{Example}[ Непрерывность интегралов с параметром]
\begin{enumerate}
    \item \textbf{Гамма-функция:} $\Gamma(\alpha) = \int\limits_0^{+\infty} x^{\alpha-1} e^{-x} \, dx$. Применяя локальные суммируемые мажоранты, можно строго доказать, что $\Gamma(\alpha) \in C(0, +\infty)$ (является непрерывной на всём положительном луче).
    \item \textbf{Интеграл с косинусом:} Докажем непрерывность функции $I(\alpha) = \int\limits_0^{+\infty} \frac{\cos(\alpha x)}{1+x^2} \, d\lambda$ на всей оси $\mathbb{R}$. \\
    Подынтегральная функция непрерывна по параметру $\alpha$. Построим суммируемую мажоранту, не зависящую от $\alpha$:
    $$ \left| \frac{\cos(\alpha x)}{1+x^2} \right| \le \frac{1}{1+x^2} $$
    Функция $h(x) = \frac{1}{1+x^2}$ интегрируема на $[0, +\infty)$ (интеграл равен $\pi/2$). Следовательно, по теореме о мажорируемом предельном переходе, интеграл $I(\alpha) \in C(\mathbb{R})$.
\end{enumerate}
\end{Example}



\begin{Theorem}{О дифференцировании интеграла по параметру}{}
Пусть задан интеграл $I(y) = \int\limits_X f(x, y) \, d\mu(x)$, где $y \in Y$ (некоторый промежуток). Предположим, что:
\begin{enumerate}
    \item Для всех $y \in Y$ функция $f(x, y)$ суммируема по $x$ на $X$.
    \item Для почти всех $x \in X$ и для всех $y \in Y$ существует частная производная $f'_y(x,y)$.
    \item Существует суммируемая функция $h(x) \in L^1(X, \mu)$, мажорирующая производную: $|f'_y(x, y)| \le h(x)$ для всех $y$.
\end{enumerate}
Тогда интеграл $I(y)$ дифференцируем на $Y$, и производную можно вносить под знак интеграла:
$$ I'(y_0) = \int\limits_X f'_y(x, y_0) \, d\mu(x) $$
\end{Theorem}

\begin{Proof}
Рассмотрим разностное отношение для интеграла:
$$ \frac{I(y_0 + h) - I(y_0)}{h} = \int\limits_X \frac{f(x, y_0 + h) - f(x, y_0)}{h} \, d\mu(x) $$
По теореме Лагранжа о конечных приращениях для функции $f$ по переменной $y$, существует такое число $\theta \in (0, 1)$, что:
$$ \frac{f(x, y_0 + h) - f(x, y_0)}{h} = f'_y(x, y_0 + \theta h) $$
По условию теоремы, производная ограничена суммируемой мажорантой $h(x)$. Следовательно, подынтегральное выражение в разностном отношении ограничено интегрируемой функцией, не зависящей от $h$. 
Остаётся применить теорему Лебега о мажорируемой сходимости при предельном переходе $h \to 0$. Предел интеграла равен интегралу от предела:
$$ \lim\limits_{h \to 0} \int\limits_X f'_y(x, y_0 + \theta h) \, d\mu(x) = \int\limits_X f'_y(x, y_0) \, d\mu(x) $$
\end{Proof}

\begin{Example}[ Вычисление интеграла Лапласа]
Рассмотрим интеграл Лапласа:
$$ I(y) = \int\limits_0^{+\infty} e^{-x^2} \cos(yx) \, dx $$
Здесь $f(x,y) = e^{-x^2} \cos(yx)$. Найдём её частную производную:
$$ f'_y(x,y) = -e^{-x^2} x \sin(yx) $$
Производная непрерывна и мажорируется суммируемой функцией $|f'_y(x,y)| \le x e^{-x^2} \in L^1([0, +\infty))$. Значит, можно дифференцировать под знаком интеграла:
$$ I'(y) = \int\limits_0^{+\infty} -e^{-x^2} x \sin(yx) \, dx $$
Проинтегрируем по частям. Пусть $u = \sin(yx) \implies du = y \cos(yx) \, dx$. \\
И $dv = -x e^{-x^2} \, dx \implies v = \frac{1}{2} e^{-x^2}$.
\begin{align*}
    I'(y) &= \left. \frac{1}{2} e^{-x^2} \sin(yx) \right|_0^{+\infty} - \frac{y}{2} \int\limits_0^{+\infty} e^{-x^2} \cos(yx) \, dx \\
    &= 0 - \frac{y}{2} I(y)
\end{align*}
Мы получили линейное дифференциальное уравнение первого порядка: $I'(y) = -\frac{y}{2} I(y)$. \\
Решаем его разделением переменных: $\frac{dI}{I} = -\frac{y}{2} dy \implies \ln|I| = -\frac{y^2}{4} + C \implies I(y) = \tilde{C} e^{-y^2/4}$. \\
Константу $\tilde{C}$ найдём, подставив $y = 0$:
$$ I(0) = \int\limits_0^{+\infty} e^{-x^2} \, dx = \frac{\sqrt{\pi}}{2} \quad \text{(интеграл Эйлера-Пуассона)} $$
Следовательно, точное значение интеграла: $I(y) = \frac{\sqrt{\pi}}{2} e^{-y^2/4}$.
\end{Example}


\section{§ Интеграл по произведению мер}

Пусть заданы два пространства с мерами: $(X, \mathcal{U}, \mu)$ и $(Y, \mathcal{V}, \nu)$. \\
Рассмотрим прямое произведение пространств $X \times Y$. Множество $\mathcal{U} \otimes \mathcal{V} = \{ A \times B \mid A \in \mathcal{U}, B \in \mathcal{V} \}$ является полукольцом (мы доказывали это ранее).

\begin{Theorem}{Мера на произведении пространств}{}
Пусть $\mu, \nu$ — \textbf{$\sigma$-конечные} меры. Тогда на полукольце измеримых прямоугольников $\mathcal{U} \otimes \mathcal{V}$ корректно определена $\sigma$-конечная мера:
$$ (\mu \times \nu)(A \times B) = \mu(A) \cdot \nu(B) $$
\end{Theorem}

\begin{Proof}
Мы знаем, что $\mathcal{U} \otimes \mathcal{V}$ — полукольцо, а заданная функция является конечно аддитивным объёмом (это доказывается стандартной разбивкой прямоугольников). Докажем \textbf{счётную аддитивность} объёма. \\
Пусть прямоугольник разбивается в счётное дизъюнктное объединение: $A \times B = \bigsqcup\limits_{i=1}^{\infty} (A_i \times B_i)$. \\
Перейдём к индикаторным функциям (характеристическим функциям множеств). Для прямоугольника индикатор распадается в произведение: $\chi_{A \times B}(x,y) = \chi_A(x) \cdot \chi_B(y)$.
Из дизъюнктности разбиения следует поточечное равенство функций:
$$ \chi_A(x) \cdot \chi_B(y) = \sum\limits_{i=1}^{\infty} \chi_{A_i}(x) \cdot \chi_{B_i}(y) = \lim\limits_{N \to \infty} \sum\limits_{i=1}^N \chi_{A_i}(x) \cdot \chi_{B_i}(y) $$
Зафиксируем произвольный $y$. Ряд состоит из неотрицательных измеримых (по $x$) функций. Применим \textbf{теорему Леви} к интегрированию по мере $\mu$:
$$ \int\limits_X \chi_A(x) \chi_B(y) \, d\mu(x) = \lim\limits_{N \to \infty} \sum\limits_{i=1}^N \int\limits_X \chi_{A_i}(x) \chi_{B_i}(y) \, d\mu(x) $$
$$ \mu(A) \chi_B(y) = \sum\limits_{i=1}^{\infty} \mu(A_i) \chi_{B_i}(y) $$
Мы получили новое поточечное равенство неотрицательных функций, зависящих только от $y$. Вновь применим теорему Леви, интегрируя обе части по мере $\nu$:
$$ \int\limits_Y \mu(A) \chi_B(y) \, d\nu(y) = \sum\limits_{i=1}^{\infty} \int\limits_Y \mu(A_i) \chi_{B_i}(y) \, d\nu(y) $$
$$ \mu(A) \cdot \nu(B) = \sum\limits_{i=1}^{\infty} \mu(A_i) \cdot \nu(B_i) $$
Счётная аддитивность доказана. Мера $\sigma$-конечна, так как $X$ и $Y$ можно разбить на счётное объединение кусков конечной меры, значит $X \times Y$ разобьётся на счётное число их попарных произведений конечной меры.
\end{Proof}

\begin{Theorem}{Тонелли и Фубини (О связи кратного и повторных интегралов)}{}
Пусть $E \subset X \times Y$. Обозначим его сечения:
$E_x = \{ y \in Y : (x,y) \in E \}$ (сечение при фиксированном $x$) и 
$E^y = \{ x \in X : (x,y) \in E \}$. \\
Площадь множества можно вычислять через интеграл от мер сечений:
$$ (\mu \times \nu)(E) = \int\limits_X \nu(E_x) \, d\mu(x) = \int\limits_Y \mu(E^y) \, d\nu(y) $$

Обобщение на функции:
\begin{itemize}
    \item \textbf{Теорема Тонелли:} Если $f(x,y) \ge 0$ измерима на произведении пространств, то:
    1) Интеграл $\mathcal{I}(x) = \int_Y f(x,y) \, d\nu(y)$ измерим для почти всех $x$.
    2) Интеграл $\mathcal{J}(y) = \int_X f(x,y) \, d\mu(x)$ измерим для почти всех $y$.
    При этом справедлива формула сведения кратного интеграла к повторным:
    $$ \int\limits_{X \times Y} f(x,y) \, d(\mu \times \nu) = \int\limits_X \left( \int\limits_Y f(x,y) d\nu(y) \right) d\mu(x) = \int\limits_Y \left( \int\limits_X f(x,y) d\mu(x) \right) d\nu(y) $$
    
    \item \textbf{Теорема Фубини:} Если функция $f$ \textbf{суммируема} (абсолютно интегрируема) на произведении пространств ($f \in L^1(X \times Y)$), то теорема работает аналогично без требования неотрицательности $f$.
\end{itemize}
\end{Theorem}

\begin{center}
\begin{tikzpicture}[scale=1.8, >=stealth, line join=round, line cap=round]

\definecolor{darktext}{HTML}{1A3A5C}
\definecolor{surfacefill}{HTML}{E8F4FD}
\definecolor{surfaceedge}{HTML}{5B9BD5}
\definecolor{sectionfill}{HTML}{2E75B6}
\definecolor{sectionline}{HTML}{1F4E79}
\definecolor{gridcolor}{HTML}{B4C7E7}
\definecolor{axiscolor}{HTML}{2C3E50}

% --- Оси 3D: y вправо, z вверх, x вглубь (вниз-влево) ---
\draw[->, thick, axiscolor] (0,0,0) -- (3.2,0,0) node[right, font=\small] {$y$};
\draw[->, thick, axiscolor] (0,0,0) -- (0,2.8,0) node[above, font=\small] {$z = f(x,y)$};
\draw[->, thick, axiscolor] (0,0,0) -- (0,0,3.2) node[below left, font=\small] {$x$};

% --- Основание XY (плоскость z=0) ---
\fill[gridcolor, opacity=0.12] (0,0,0) -- (2.5,0,0) -- (2.5,0,2.5) -- (0,0,2.5) -- cycle;
\draw[axiscolor, dashed, opacity=0.4] (0,0,2.5) -- (2.5,0,2.5) -- (2.5,0,0);
\foreach \t in {0,0.5,...,2.5} {
    \draw[axiscolor, dashed, opacity=0.15] (\t,0,0) -- (\t,0,2.5);
    \draw[axiscolor, dashed, opacity=0.15] (0,0,\t) -- (2.5,0,\t);
}
\node[axiscolor, opacity=0.7] at (1.25, -0.25, 1.25) {$X \times Y$};

% ============================================
% ПОВЕРХНОСТЬ z = f(x,y)
% Координаты в TikZ: (y, z, x)
% ============================================

% Передняя граница x = 0 (ближе к зрителю)
\coordinate (F0) at (0.00, 1.20, 0.00);
\coordinate (F1) at (0.31, 1.05, 0.00);
\coordinate (F2) at (0.62, 0.95, 0.00);
\coordinate (F3) at (0.94, 0.90, 0.00);
\coordinate (F4) at (1.25, 1.00, 0.00);
\coordinate (F5) at (1.56, 1.15, 0.00);
\coordinate (F6) at (1.88, 1.35, 0.00);
\coordinate (F7) at (2.19, 1.55, 0.00);
\coordinate (F8) at (2.50, 1.70, 0.00);

% Задняя граница x = 2.5 (дальше от зрителя)
\coordinate (B0) at (0.00, 1.60, 2.50);
\coordinate (B1) at (0.31, 1.40, 2.50);
\coordinate (B2) at (0.62, 1.35, 2.50);
\coordinate (B3) at (0.94, 1.30, 2.50);
\coordinate (B4) at (1.25, 1.40, 2.50);
\coordinate (B5) at (1.56, 1.55, 2.50);
\coordinate (B6) at (1.88, 1.75, 2.50);
\coordinate (B7) at (2.19, 1.95, 2.50);
\coordinate (B8) at (2.50, 2.10, 2.50);

% Сечение x = 1.25 — точно на поверхности (линейная интерполяция между F и B)
\coordinate (C0) at (0.00, 1.40, 1.25);
\coordinate (C1) at (0.31, 1.225, 1.25);
\coordinate (C2) at (0.62, 1.15, 1.25);
\coordinate (C3) at (0.94, 1.10, 1.25);
\coordinate (C4) at (1.25, 1.20, 1.25);
\coordinate (C5) at (1.56, 1.35, 1.25);
\coordinate (C6) at (1.88, 1.55, 1.25);
\coordinate (C7) at (2.19, 1.75, 1.25);
\coordinate (C8) at (2.50, 1.90, 1.25);

% Заливка поверхности (прозрачная, чтобы видно было сечение)
\fill[surfacefill, opacity=0.25] 
    plot[smooth, tension=0.8] coordinates {(F0)(F1)(F2)(F3)(F4)(F5)(F6)(F7)(F8)}
    -- plot[smooth, tension=0.8] coordinates {(B8)(B7)(B6)(B5)(B4)(B3)(B2)(B1)(B0)}
    -- cycle;

% Линии поверхности по направлению x (от передней к задней границе)
\foreach \i in {0,...,8} {
    \draw[surfaceedge, opacity=0.3, very thin] (F\i) -- (B\i);
}

% Контур поверхности
\draw[surfaceedge, thick] plot[smooth, tension=0.8] coordinates {(F0)(F1)(F2)(F3)(F4)(F5)(F6)(F7)(F8)};
\draw[surfaceedge, thick] plot[smooth, tension=0.8] coordinates {(B0)(B1)(B2)(B3)(B4)(B5)(B6)(B7)(B8)};
\draw[surfaceedge, thick] (F0) -- (B0);
\draw[surfaceedge, thick] (F8) -- (B8);

% ============================================
% СЕЧЕНИЕ при x = x_0 — параллельно плоскости рисунка
% ============================================

% Заливка от основания z=0 до верхней кривой поверхности
\fill[sectionfill, opacity=0.35] 
    (0,0,1.25) -- 
    plot[smooth, tension=0.8] coordinates {(C0)(C1)(C2)(C3)(C4)(C5)(C6)(C7)(C8)} 
    -- (2.5,0,1.25) -- cycle;

% Штриховка (вертикальные линии внутри сечения, обрезанные по контуру)
\begin{scope}
    \clip (0,0,1.25) -- 
          plot[smooth, tension=0.8] coordinates {(C0)(C1)(C2)(C3)(C4)(C5)(C6)(C7)(C8)} 
          -- (2.5,0,1.25) -- cycle;
    \foreach \y in {0.1,0.2,...,2.4} {
        \draw[sectionfill, opacity=0.5, very thin] (\y,0,1.25) -- (\y,2.2,1.25);
    }
\end{scope}

% Верхняя граница — часть самой поверхности, ультра-жирная и контрастная
\draw[sectionline, ultra thick] plot[smooth, tension=0.8] coordinates {(C0)(C1)(C2)(C3)(C4)(C5)(C6)(C7)(C8)};

% Боковые и нижняя границы сечения
\draw[sectionline, thick] (0,0,1.25) -- (C0);
\draw[sectionline, thick] (2.5,0,1.25) -- (C8);
\draw[sectionline, thick] (0,0,1.25) -- (2.5,0,1.25);

% ============================================
% Оформление
% ============================================

% Точка x_0 на оси x (в точке начала сечения)
\fill[sectionline] (0,0,1.25) circle (0.05);
\draw[axiscolor, thick, dashed] (0,0,1.25) -- (-0.15,0,1.25) 
    node[left, font=\small\bfseries] {$x_0$};

% Подписи
\node[darktext, align=left, anchor=west, font=\small] at (3.0, 2.2) {
    \textbf{\large Теорема Фубини} \\[4pt]
    Объём под поверхностью $z = f(x,y)$ \\ 
    вычисляется интегрированием площадей \\ 
    сечений $x = x_0 = \text{const}$ по мере $\nu$: \\[6pt]
    $\displaystyle \iint_{X \times Y} f(x,y)\, d(\mu \times \nu) = 
    \int_X \underbrace{\left( \int_Y f(x,y)\, d\nu(y) \right)}_{\mathcal{I}(x)} d\mu(x)$
};

\draw[sectionfill, thick, ->] (3.8, 0.5) -- (2, 1.0, 1.7) 
    node[midway, above, font=\tiny, sectionfill] {сечение $\mathcal{I}(x_0)$};

\end{tikzpicture}
\end{center}

\begin{Proof}[Доказательство теорем Тонелли и Фубини]
Доказательство проводится стандартным в теории меры поэтапным методом (от простых множеств к произвольным функциям).

\textbf{Этап 1: Индикатор измеримого прямоугольника.} \\
Пусть $f(x,y) = \chi_{A \times B}(x,y) = \chi_A(x) \chi_B(y)$, где $A \in \mathcal{U}, B \in \mathcal{V}$. 
Тогда интеграл по $Y$ от сечения равен:
$$ \mathcal{I}(x) = \int\limits_Y \chi_A(x) \chi_B(y) \, d\nu(y) = \chi_A(x) \int\limits_Y \chi_B(y) \, d\nu(y) = \chi_A(x) \cdot \nu(B) $$
Эта функция измерима по $x$. Интегрируем её по $X$:
$$ \int\limits_X \mathcal{I}(x) \, d\mu(x) = \int\limits_X \chi_A(x) \nu(B) \, d\mu(x) = \mu(A) \cdot \nu(B) $$
С другой стороны, по определению произведения мер, двойной интеграл равен:
$$ \int\limits_{X \times Y} \chi_{A \times B} \, d(\mu \times \nu) = (\mu \times \nu)(A \times B) = \mu(A) \cdot \nu(B) $$
Равенство выполняется. (Аналогично для интеграла $\mathcal{J}(y)$).

\textbf{Этап 2: Индикатор произвольного измеримого множества.} \\
Пусть $f(x,y) = \chi_E(x,y)$, где $E \in \mathcal{U} \otimes \mathcal{V}$. Заметим, что интеграл от индикатора по слою равен мере сечения: $\int\limits_Y \chi_E(x,y) \, d\nu(y) = \nu(E_x)$. \\
В силу $\sigma$-конечности мер, класс множеств, для которых выполняется равенство $(\mu \times \nu)(E) = \int_X \nu(E_x) d\mu(x)$, содержит полукольцо измеримых прямоугольников и замкнут относительно счетных дизъюнктных объединений и монотонных пределов. По лемме о монотонных классах, равенство распространяется на всю $\sigma$-алгебру измеримых множеств.

\textbf{Этап 3: Простая неотрицательная функция.} \\
Всякая простая неотрицательная функция представима в виде конечной линейной комбинации индикаторов измеримых множеств: $f(x,y) = \sum_{i=1}^n c_i \chi_{E_i}(x,y)$, где $c_i \ge 0$. Теорема выполняется в силу линейности интеграла Лебега как для кратного интеграла, так и для повторных.

\textbf{Этап 4: Произвольная неотрицательная функция (Теорема Тонелли).} \\
Пусть $f(x,y) \ge 0$ — измеримая функция на $X \times Y$. По теореме об аппроксимации существует монотонно возрастающая последовательность простых неотрицательных функций $f_n(x,y) \nearrow f(x,y)$.
Для каждой $f_n$ равенство кратного и повторного интегралов уже доказано:
$$ \int\limits_{X \times Y} f_n \, d(\mu \times \nu) = \int\limits_X \left( \int\limits_Y f_n \, d\nu \right) d\mu $$
Так как $f_n \nearrow f$, то интегралы по сечениям также монотонно возрастают $\int\limits_Y f_n \, d\nu \nearrow \int\limits_Y f \, d\nu$. 
Дважды применяя теорему Леви (о монотонной сходимости) к обеим частям равенства, получаем предельную формулу для функции $f$. Теорема Тонелли доказана.

\textbf{Этап 5: Произвольная суммируемая функция (Теорема Фубини).} \\
Пусть $f \in L^1(X \times Y, \mu \times \nu)$. Разложим функцию на положительную и отрицательную части: $f = f_+ - f_-$, где $f_+, f_- \ge 0$. \\
Так как $f$ суммируема, обе функции $f_+$ и $f_-$ суммируемы (их интегралы конечны). Применим доказанную теорему Тонелли к каждой из них по отдельности. 
Поскольку интегралы конечны, мы можем воспользоваться линейностью интеграла и вычесть уравнения друг из друга без риска получить неопределенность $(\infty - \infty)$:
\begin{align*}
    \int\limits_{X \times Y} f \, d(\mu \times \nu) &= \int\limits_{X \times Y} f_+ \, d(\mu \times \nu) - \int\limits_{X \times Y} f_- \, d(\mu \times \nu) \\
    &= \int\limits_X \left( \int\limits_Y f_+ \, d\nu \right) d\mu - \int\limits_X \left( \int\limits_Y f_- \, d\nu \right) d\mu \\
    &= \int\limits_X \left( \int\limits_Y (f_+ - f_-) \, d\nu \right) d\mu = \int\limits_X \left( \int\limits_Y f \, d\nu \right) d\mu
\end{align*}
Теорема Фубини доказана.
\end{Proof}

\section{§ Интеграл Лебега и интеграл Римана}

\begin{Theorem}{Связь интегралов Римана и Лебега}{}
Если функция $f(x)$ интегрируема по Риману на отрезке $[a, b]$ (то есть $f \in R[a,b]$), то она суммируема по мере Лебега на этом отрезке ($f \in L^1([a,b], \lambda)$), и их значения совпадают:
$$ (R) \int\limits_a^b f(x) \, dx = (L) \int\limits_{[a, b]} f(x) \, d\lambda $$
\end{Theorem}

\begin{Proof}
Рассмотрим разбиение отрезка $x_k = a + \frac{k}{n}(b-a)$ для $k \in \{0, \dots, n\}$.
Построим нижние и верхние суммы Дарбу:
$$ \underline{S}_n = \sum\limits_{i=1}^n m_i \Delta x_i, \quad \text{где } m_i = \inf\limits_{\Delta_i} f(x) $$
$$ \overline{S}_n = \sum\limits_{i=1}^n M_i \Delta x_i, \quad \text{где } M_i = \sup\limits_{\Delta_i} f(x) $$
Определим две ступенчатые (простые) функции: $\underline{f}_n(x) = m_i$ и $\overline{f}_n(x) = M_i$ на интервалах $\Delta_i$. 
\begin{center}
\begin{tikzpicture}[scale=1.2, >=latex]
    \definecolor{darktext}{HTML}{154360}
    \definecolor{colorUnder}{HTML}{7DCEA0}
    \definecolor{colorOver}{HTML}{F5B041}
    
    % Оси
    \draw[->, thick, darktext] (-0.5, 0) -- (5, 0) node[right] {$x$};
    \draw[->, thick, darktext] (0, -0.5) -- (0, 3) node[above] {$y$};
    
    % Кривая f(x)
    \draw[ultra thick, darktext] plot [smooth, tension=0.7] coordinates { (0.5, 1.5) (1.5, 2.5) (2.5, 1.0) (3.5, 2.0) (4.5, 0.5) };
    \node[darktext] at (1.5, 2.8) {$f(x)$};
    
    % Нижняя сумма Дарбу (Зеленая)
    \filldraw[fill=colorUnder, fill opacity=0.3, draw=colorUnder, thick] (0.5, 0) rectangle (1.5, 1.5);
    \filldraw[fill=colorUnder, fill opacity=0.3, draw=colorUnder, thick] (1.5, 0) rectangle (2.5, 1.0);
    \filldraw[fill=colorUnder, fill opacity=0.3, draw=colorUnder, thick] (2.5, 0) rectangle (3.5, 1.0);
    \filldraw[fill=colorUnder, fill opacity=0.3, draw=colorUnder, thick] (3.5, 0) rectangle (4.5, 0.5);
    
    % Верхняя сумма Дарбу (Оранжевая - пустые контуры)
    \draw[colorOver, thick, dashed] (0.5, 0) rectangle (1.5, 2.5);
    \draw[colorOver, thick, dashed] (1.5, 0) rectangle (2.5, 2.5);
    \draw[colorOver, thick, dashed] (2.5, 0) rectangle (3.5, 2.0);
    \draw[colorOver, thick, dashed] (3.5, 0) rectangle (4.5, 2.0);
    
    % Подписи
    \node[colorUnder, font=\bfseries] at (2.5, 0.5) {$\underline{f}_n \le f$};
    \node[colorOver, font=\bfseries] at (3.4, 2.3) {$\overline{f}_n \ge f$};
\end{tikzpicture}
\end{center}

Тогда интегралы Лебега от этих простых функций в точности равны суммам Дарбу:
$$ \underline{S}_n = \int\limits_{[a,b]} \underline{f}_n \, d\lambda \quad \text{и} \quad \overline{S}_n = \int\limits_{[a,b]} \overline{f}_n \, d\lambda $$
При увеличении $n$ (измельчении разбиения) функции $\underline{f}_n$ монотонно не убывают, а $\overline{f}_n$ монотонно не возрастают. Все они ограничены функцией $f$.
Значит, существуют поточечные пределы: $\underline{f}_n \nearrow \underline{f}$ и $\overline{f}_n \searrow \overline{f}$. Причём $\underline{f} \le f \le \overline{f}$ почти всюду.
По теореме Леви, пределы интегралов равны интегралам от пределов:
$$ \lim\limits_{n \to \infty} \underline{S}_n = \int\limits_{[a,b]} \underline{f} \, d\lambda \quad \text{и} \quad \lim\limits_{n \to \infty} \overline{S}_n = \int\limits_{[a,b]} \overline{f} \, d\lambda $$
Так как функция $f$ интегрируема по Риману, пределы сумм Дарбу совпадают. Значит:
$$ \int\limits_{[a,b]} \overline{f} \, d\lambda = \int\limits_{[a,b]} \underline{f} \, d\lambda \implies \int\limits_{[a,b]} (\overline{f} - \underline{f}) \, d\lambda = 0 $$
Так как подынтегральная функция неотрицательна, интеграл равен нулю только если $\overline{f} = \underline{f}$ почти всюду. Из неравенства $\underline{f} \le f \le \overline{f}$ следует, что $f$ совпадает с $\underline{f}$ почти всюду.
Так как $\underline{f}$ является измеримой (как предел простых), то и $f$ измерима, и её интеграл Лебега совпадает с общим пределом сумм Дарбу (интегралом Римана).
\end{Proof}

\begin{Theorem}{Следствие (Абсолютная сходимость) }{}
Если несобственный интеграл Римана сходится абсолютно, то функция интегрируема по Лебегу, и интегралы равны:
$$ (R) \int\limits_a^b f \, dx = (L) \int\limits_{[a, b)} f \, d\lambda $$
\end{Theorem}

\subsection*{Несобственные интегралы Лебега}

\begin{Definition}{Несобственный интеграл Лебега}{}
Говорят, что $f \in L^1_{loc}([a,b))$ (локально суммируема), если $f \in L^1([a, w])$ для любого $w < b$. 
Несобственный интеграл Лебега определяется как предел:
$$ \int\limits_{[a, b)} f \, d\lambda \stackrel{\text{def}}{=} \lim\limits_{w \to b-0} \int\limits_{[a, w]} f \, d\lambda $$
\end{Definition}

\begin{Definition}{Равномерная сходимость несобственных интегралов}{}
Пусть $I(y) = \int\limits_{[a, b)} f(x, y) \, d\lambda$ сходится для всех $y \in Y$. Говорят, что этот интеграл сходится \textbf{равномерно по параметру $y$}, если интегралы по урезанным отрезкам равномерно сходятся к $I(y)$:
$$ \int\limits_{[a, w]} f(x, y) \, d\lambda \rightrightarrows_{w \to b-0} I(y) $$
Эквивалентное определение через "хвост" интеграла:
$$ \forall \varepsilon > 0 \quad \exists \Delta > 0 \quad \forall \delta \in (b-\Delta, b) \quad \forall y \in Y \implies \left| \int\limits_{[\delta, b)} f(x,y) \, d\lambda \right| < \varepsilon $$
\end{Definition}

\begin{Example}
Рассмотрим интеграл $\int\limits_{[0, +\infty)} e^{-xy} \, d\lambda$.
Вычислим его "хвост" на отрезке $[\delta, +\infty)$:
$$ \int\limits_{[\delta, +\infty)} e^{-xy} \, dx = \left. -\frac{1}{y} e^{-xy} \right|_{\delta}^{+\infty} = \frac{e^{-\delta y}}{y} $$
\begin{itemize}
    \item Пусть $y \in [1, +\infty)$. Тогда $\frac{e^{-\delta y}}{y} \le e^{-\delta y} \le e^{-\delta}$. Выбирая достаточно большое $\delta$, мы сделаем хвост меньше $\varepsilon$ независимо от $y$. Значит, интеграл \textbf{сходится равномерно}.
    \item Пусть $y \in (0, +\infty)$. На этом множестве равномерной сходимости нет. Если мы возьмём малую точку $y = \frac{1}{\delta}$, хвост будет равен:
    $$ \frac{e^{-\delta \cdot (1/\delta)}}{1/\delta} = \delta e^{-1} $$
    Для любого $\varepsilon$ можно найти такое маленькое $\delta$, что хвост станет сколь угодно большим, что нарушает определение равномерной сходимости.
\end{itemize}
\end{Example}



\section{§ Признаки равномерной сходимости несобственных интегралов}

\begin{Theorem}{Признак Вейерштрасса}{}
Пусть для функции $f(x,y)$ существует мажоранта $g(x) \in L^1([a,b), \lambda)$, такая что $|f(x,y)| \le g(x)$ для почти всех $x$ и для всех $y \in Y$. \\
Тогда несобственный интеграл $\int\limits_{[a,b)} f(x,y) \, d\lambda$ сходится равномерно на множестве $Y$.
\end{Theorem}
\begin{Proof}
Очевидно вытекает из оценки "хвоста" интеграла: $\left| \int\limits_{[\delta, b)} f(x,y) \, d\lambda \right| \le \int\limits_{[\delta, b)} g(x) \, d\lambda$. Так как $g$ суммируема, хвост стремится к нулю независимо от $y$.
\end{Proof}

\begin{Theorem}{Критерий Коши}{}
Интеграл $\int\limits_{[a,b)} f(x,y) \, d\lambda$ сходится равномерно на $Y$ тогда и только тогда, когда:
$$ \forall \varepsilon > 0 \quad \exists \Delta > 0 \quad \forall \delta', \delta'' \in (b-\Delta, b) \quad \forall y \in Y \implies \left| \int\limits_{[\delta', \delta'')} f(x,y) \, d\lambda \right| < \varepsilon $$
\end{Theorem}

\begin{Theorem}{Признаки Абеля и Дирихле}{}
Пусть функции $f(x,y)$ и $g(x,y)$ непрерывны по $x$ на $[a, b)$ при каждом $y \in Y$, причём функция $g(x,y)$ \textbf{монотонна} по $x$ при каждом фиксированном $y$.
Интеграл $\int\limits_{[a,b)} f(x,y)g(x,y) \, d\lambda$ сходится равномерно на $Y$, если выполняется одно из двух условий:

\textbf{1. Признак Дирихле:}
\begin{itemize}
    \item Интеграл от $f$ равномерно ограничен: $\exists C > 0 : \left| \int\limits_{[a,\xi)} f(x,y) \, d\lambda \right| \le C \quad \forall \xi \in [a, b), \forall y \in Y$.
    \item Функция $g(x,y) \rightrightarrows 0$ при $x \to b-0$ равномерно по $y$.
\end{itemize}

\textbf{2. Признак Абеля:}
\begin{itemize}
    \item Интеграл $\int\limits_{[a,b)} f(x,y) \, d\lambda$ сходится равномерно на $Y$.
    \item Функция $g(x,y)$ равномерно ограничена по совокупности переменных: $\exists M > 0 : |g(x,y)| \le M$.
\end{itemize}
\end{Theorem}

\begin{Proof}[Доказательство признака Абеля]
Покажем равномерную сходимость с помощью критерия Коши. Зададим $\varepsilon > 0$. \\
В силу непрерывности $f$ и $g$, а также монотонности $g$ по $x$, мы можем применить ко второй интегральной теореме о среднем (теореме Бонне). Для любых $\delta', \delta''$ существует точка $\xi \in [\delta', \delta'']$, такая что:
$$ \int\limits_{\delta'}^{\delta''} f(x,y) g(x,y) \, dx = g(\delta', y) \int\limits_{\delta'}^{\xi} f(x,y) \, dx + g(\delta'', y) \int\limits_{\xi}^{\delta''} f(x,y) \, dx $$
Так как интеграл от $f$ сходится равномерно, по критерию Коши существует такое $\Delta > 0$, что для любых отрезков внутри $(b-\Delta, b)$ модули интегралов от $f$ строго меньше $\frac{\varepsilon}{2M}$. \\
Тогда, используя равномерную ограниченность $|g| \le M$:
$$ \left| \int\limits_{\delta'}^{\delta''} f(x,y) g(x,y) \, dx \right| \le M \cdot \frac{\varepsilon}{2M} + M \cdot \frac{\varepsilon}{2M} = \varepsilon $$
Равномерная сходимость по критерию Коши доказана.
\end{Proof}


\section{§ Теоремы о перестановке пределов}

\begin{Theorem}{О перестановке двух пределов}{}
Пусть задана функция $f: X \times Y \to \mathbb{R}$, где $X, Y$ — метрические пространства. Пусть существуют пределы:
\begin{enumerate}
    \item $\forall y \in Y \quad \exists \lim\limits_{x \to x_0} f(x,y) = h(y)$
    \item $\forall x \in X \quad \exists \lim\limits_{y \to y_0} f(x,y) = \xi(x)$
\end{enumerate}
Если \textbf{хотя бы один} из этих двух пределов является \textbf{равномерным} (по другой переменной), то двойные пределы совпадают:
$$ \lim\limits_{y \to y_0} h(y) = \lim\limits_{x \to x_0} \xi(x) $$
\textit{(То есть предельные переходы можно менять местами: $\lim\limits_{y \to y_0} \lim\limits_{x \to x_0} f(x,y) = \lim\limits_{x \to x_0} \lim\limits_{y \to y_0} f(x,y)$).}
\end{Theorem}

\begin{Proof}
Пусть второй предел является равномерным: $f(x,y) \rightrightarrows \xi(x)$ при $y \to y_0$ равномерно по $x$.

\textbf{А) Докажем, что существует предел $\lim\limits_{y \to y_0} h(y) = A$.} \\
Зададим $\varepsilon > 0$. Из равномерной сходимости по критерию Коши следует, что существует проколотая окрестность $\mathring{U}_\delta(y_0)$, такая что для любых $y', y'' \in \mathring{U}_\delta(y_0)$ и для всех $x \in X$:
$$ |f(x, y') - f(x, y'')| < \varepsilon $$
Так как это неравенство верно для любого $x$, перейдём в нём к пределу при $x \to x_0$. Поскольку $\lim_{x \to x_0} f(x,y) = h(y)$, предельный переход сохраняет нестрогое неравенство:
$$ |h(y') - h(y'')| \le \varepsilon $$
В силу критерия Коши для функции $h(y)$, предел существует. Обозначим $\lim_{y \to y_0} h(y) = A$.

\textbf{Б) Докажем, что $\lim\limits_{x \to x_0} \xi(x) = A$.} \\
Рассмотрим модуль разности и применим неравенство треугольника:
$$ |\xi(x) - A| = |\xi(x) - f(x,y) + f(x,y) - h(y) + h(y) - A| \le |\xi(x) - f(x,y)| + |f(x,y) - h(y)| + |h(y) - A| $$
Нам нужно сделать все три слагаемых меньше $\frac{\varepsilon}{3}$. 
1) Так как $\lim_{y \to y_0} h(y) = A$, существует окрестность $U_1(y_0)$, где $|h(y) - A| < \frac{\varepsilon}{3}$.
2) Так как $f(x,y) \rightrightarrows \xi(x)$, существует окрестность $U_2(y_0)$, где $|f(x,y) - \xi(x)| < \frac{\varepsilon}{3}$ для \textbf{всех} $x$. \\
Зафиксируем какую-нибудь одну точку $\bar{y} \in U_1(y_0) \cap U_2(y_0)$. Первое и третье слагаемые для неё уже меньше $\frac{\varepsilon}{3}$.
3) Для этой фиксированной $\bar{y}$ существует $\lim_{x \to x_0} f(x,\bar{y}) = h(\bar{y})$. Следовательно, существует окрестность $U(x_0)$, такая что для всех $x \in U(x_0)$ второе слагаемое $|f(x,\bar{y}) - h(\bar{y})| < \frac{\varepsilon}{3}$.
Складывая оценки, для всех $x \in U(x_0)$ получаем $|\xi(x) - A| < \varepsilon$. Предел доказан.
\end{Proof}

\begin{Theorem}{Предельный переход под знаком несобственного интеграла}{}
Пусть $Y$ — метрическое пространство, $f: [a, b) \times Y \to \mathbb{R}$. Если:
\begin{enumerate}
    \item Для почти всех $x \in [a, b)$ существует $\lim\limits_{y \to y_0} f(x,y) = g(x)$.
    \item Функция $g \in L^1_{loc}([a,b))$, и для любого $t \in [a, b)$ предел интеграла равен интегралу предела: $\lim\limits_{y \to y_0} \int\limits_{[a,t]} f(x,y) \, d\lambda = \int\limits_{[a,t]} g(x) \, d\lambda$.
    \item Интеграл $\int\limits_{[a,b)} f(x,y) \, d\lambda$ сходится \textbf{равномерно} на $Y$.
\end{enumerate}
Тогда предельный переход правомерен:
$$ \lim\limits_{y \to y_0} \int\limits_{[a,b)} f(x,y) \, d\lambda = \int\limits_{[a,b)} g(x) \, d\lambda $$
\end{Theorem}

\begin{Proof}
Рассмотрим функцию $H_1(t,y) = \int\limits_{[a,t]} f(x,y) \, d\lambda$. 
По условию 2, существует предел $\lim\limits_{y \to y_0} H_1(t,y) = \int\limits_{[a,t]} g(x) \, d\lambda \equiv H_2(t)$.
По условию 3 (равномерная сходимость интеграла), существует предел $\lim\limits_{t \to b-0} H_1(t,y) = \int\limits_{[a,b)} f(x,y) \, d\lambda \equiv H_3(y)$, причём этот предел \textbf{равномерный} по параметру $y$.
Мы находимся в условиях теоремы о перестановке пределов. Так как один из пределов ($H_3$) равномерный, мы имеем право их переставить:
$$ \lim\limits_{y \to y_0} H_3(y) = \lim\limits_{t \to b-0} H_2(t) \implies \lim\limits_{y \to y_0} \int\limits_{[a,b)} f(x,y) \, d\lambda = \lim\limits_{t \to b-0} \int\limits_{[a,t]} g(x) \, d\lambda = \int\limits_{[a,b)} g(x) \, d\lambda $$
\end{Proof}


\section{§ Вычисление интеграла Дирихле}

Продемонстрируем мощь этих теорем, вычислив знаменитый интеграл Дирихле.
Рассмотрим функцию с параметром $y \ge 0$:
$$ I(y) = \int\limits_0^{+\infty} e^{-xy} \frac{\sin x}{x} \, dx $$
Очевидно, что исходный интеграл Дирихле равен $I(0) = \lim\limits_{y \to +0} I(y)$. 
Чтобы внести предел под знак интеграла, проверим признак Абеля:
\begin{itemize}
    \item Интеграл $\int\limits_0^{+\infty} \frac{\sin x}{x} \, dx$ сходится (это классический факт). Он не зависит от $y$, значит сходится "равномерно".
    \item Функция $e^{-xy}$ при $y \ge 0$ монотонно убывает по $x$ и равномерно ограничена: $|e^{-xy}| \le 1$.
\end{itemize}
По признаку Абеля интеграл $I(y)$ сходится \textbf{равномерно} при $y \ge 0$. Следовательно, функция $I(y)$ непрерывна в нуле, и мы имеем право написать: $I(0) = \lim\limits_{y \to +0} I(y)$.

Вычислим производную $I'(y)$. Формально дифференцируя по параметру $y$:
$$ I'(y) = \int\limits_0^{+\infty} \left(e^{-xy} \frac{\sin x}{x}\right)'_y \, dx = -\int\limits_0^{+\infty} e^{-xy} \sin x \, dx $$
Правомерно ли это? Для любого отрезка $[\delta, +\infty)$, где $\delta > 0$, производная мажорируется: $|e^{-xy} \sin x| \le e^{-xy} \le e^{-\delta x} \in L^1(0, +\infty)$. Значит, при $y > 0$ дифференцировать можно.

Вычисляем полученный интеграл $I'(y)$ двукратным интегрированием по частям:
$$ I'(y) = -\frac{1}{1+y^2} $$
Интегрируем обратно по $y$, чтобы найти $I(y)$:
$$ I(y) = -\operatorname{arctg}(y) + C, \quad y \ge 0 $$
Чтобы найти константу $C$, устремим $y \to +\infty$. Оценим исходный интеграл:
$$ |I(y)| \le \int\limits_0^{+\infty} \left| e^{-xy} \frac{\sin x}{x} \right| \, dx \le \int\limits_0^{+\infty} e^{-xy} \, dx = \left. -\frac{1}{y} e^{-xy} \right|_0^{+\infty} = \frac{1}{y} $$
При $y \to +\infty$ оценка $\frac{1}{y} \to 0$, следовательно $\lim\limits_{y \to +\infty} I(y) = 0$.
Подставим это в наше уравнение:
$$ 0 = -\operatorname{arctg}(+\infty) + C \implies 0 = -\frac{\pi}{2} + C \implies C = \frac{\pi}{2} $$
Таким образом, $I(y) = \frac{\pi}{2} - \operatorname{arctg}(y)$.
Искомый интеграл Дирихле:
$$ \int\limits_0^{+\infty} \frac{\sin x}{x} \, dx = I(0) = \frac{\pi}{2} - 0 = \frac{\pi}{2} $$


\section{§ Интегрирование несобственного интеграла по параметру}

\begin{Theorem}{Об интегрировании интеграла с параметром}{}
Пусть $f: [a, b) \times Y \to \mathbb{R}$, $f \in L^1_{loc}([a,b))$ для любого $y \in Y$. Пусть $(Y, \mathcal{U}, \nu)$ — измеримое пространство с мерой $\nu(Y) < +\infty$.
Обозначим:
$$ I(x) = \int\limits_Y f(x,y) \, d\nu(y) \quad \text{и} \quad \mathcal{J}(y) = \int\limits_{[a,b)} f(x,y) \, dx $$
Предположим, что функция $f$ суммируема по $Y$ для любого $x$, и несобственный интеграл $\mathcal{J}(y)$ сходится \textbf{равномерно} на $Y$. Тогда интеграл $I(x)$ интегрируем на $[a,b)$, и порядки интегрирования можно менять:
$$ \int\limits_{[a,b)} I(x) \, dx = \int\limits_Y \mathcal{J}(y) \, d\nu(y) \quad \left( \text{т.е. } \int\limits_{[a,b)} \int\limits_Y f(x,y) \, d\nu \, dx = \int\limits_Y \int\limits_{[a,b)} f(x,y) \, dx \, d\nu \right) $$
\end{Theorem}

\begin{Proof}[Доказательство]
Обозначим интегралы по урезанным отрезкам: $J_w(y) = \int\limits_a^w f(x,y) \, dx$. \\
Из определения равномерной сходимости несобственного интеграла следует, что $J_w(y) \rightrightarrows \mathcal{J}(y)$ на множестве $Y$ при $w \to b-0$. \\
Так как отрезок $[a, w]$ конечен, мы можем применить классическую теорему Фубини для перестановки конечных интегралов:
$$ \int\limits_a^w I(x) \, dx = \int\limits_a^w \left( \int\limits_Y f(x,y) \, d\nu \right) dx = \int\limits_Y \left( \int\limits_a^w f(x,y) \, dx \right) d\nu = \int\limits_Y J_w(y) \, d\nu $$
Перейдём к пределу в обеих частях равенства при $w \to b-0$.
С левой стороны получаем несобственный интеграл по определению:
$$ \lim\limits_{w \to b-0} \int\limits_a^w I(x) \, dx = \int\limits_{[a,b)} I(x) \, dx $$
С правой стороны стоит предел: $\lim\limits_{w \to b-0} \int\limits_Y J_w(y) \, d\nu$. Так как последовательность функций $J_w(y)$ сходится равномерно к $\mathcal{J}(y)$, и мера пространства конечна ($\nu(Y) < +\infty$), мы имеем право внести предел под знак интеграла по теореме о равномерном предельном переходе:
$$ \lim\limits_{w \to b-0} \int\limits_Y J_w(y) \, d\nu = \int\limits_Y \left( \lim\limits_{w \to b-0} J_w(y) \right) d\nu = \int\limits_Y \mathcal{J}(y) \, d\nu $$
Приравнивая пределы, получаем требуемое равенство.
\end{Proof}



\section{§ Примеры применения интегралов с параметром}

\begin{Example}[ Вычисление интеграла Лапласа (окончание)]
Рассмотрим интеграл $K(y) = \int\limits_0^{+\infty} \frac{\cos(yx)}{1+x^2} \, dx$.
В предыдущих лекциях мы показали, что $K(y)$ непрерывна на $\mathbb{R}$, и $K(0) = \frac{\pi}{2}$.

\textbf{1. Дифференцирование по параметру:}
Формально продифференцируем по $y$:
$$ K'(y) = -\int\limits_0^{+\infty} \frac{x \sin(yx)}{1+x^2} \, dx $$
Обоснуем законность дифференцирования. Подынтегральная функция не имеет интегрируемой мажоранты на всей оси $(0, +\infty)$, не зависящей от $y$. Однако воспользуемся признаком Дирихле:
$$ \left| \int\limits_a^A \sin(yx) \, dx \right| = \left| \frac{-\cos(yx)}{y} \right|_a^A \le \frac{2}{y} \le \frac{2}{y_0} $$
Для любого $y \in [y_0, +\infty)$, где $y_0 > 0$, эта величина равномерно ограничена. Функция $\frac{x}{1+x^2}$ монотонно стремится к нулю. Значит, по признаку Дирихле, интеграл сходится \textbf{равномерно локально} на луче $[y_0, +\infty)$. Это позволяет дифференцировать под знаком интеграла при $y > 0$.

\textbf{2. Составление дифференциального уравнения:}
Преобразуем $K'(y)$ алгебраически (прибавим и вычтем $x$ в числителе):
$$ K'(y) = -\int\limits_0^{+\infty} \frac{x^2 \sin(yx)}{x(1+x^2)} \, dx = -\int\limits_0^{+\infty} \frac{(x^2+1-1) \sin(yx)}{x(1+x^2)} \, dx = -\int\limits_0^{+\infty} \frac{\sin(yx)}{x} \, dx + \int\limits_0^{+\infty} \frac{\sin(yx)}{x(1+x^2)} \, dx $$
Первый интеграл — это интеграл Дирихле, равный $\frac{\pi}{2}$ (при $y > 0$).
Продифференцируем полученное выражение ещё раз по $y$:
$$ K''(y) = 0 + \int\limits_0^{+\infty} \frac{x \cos(yx)}{x(1+x^2)} \, dx = \int\limits_0^{+\infty} \frac{\cos(yx)}{1+x^2} \, dx = K(y) $$
Мы получили линейное дифференциальное уравнение второго порядка: $K''(y) - K(y) = 0$.

\textbf{3. Решение:}
Фундаментальная система решений: $K(y) = C_1 e^y + C_2 e^{-y}$.
Так как исходный интеграл ограничен: $|K(y)| \le \int\limits_0^{+\infty} \frac{1}{1+x^2} dx = \frac{\pi}{2}$, то при $y \to +\infty$ должно выполняться $C_1 = 0$.
Так как $K(0) = \frac{\pi}{2}$, получаем $C_2 = \frac{\pi}{2}$.
Итого, $K(y) = \frac{\pi}{2} e^{-y}$ для $y \ge 0$. В силу чётности косинуса, функция $K(y)$ чётная, поэтому окончательный ответ на всей оси:
$$ \int\limits_0^{+\infty} \frac{\cos(yx)}{1+x^2} \, dx = \frac{\pi}{2} e^{-|y|} $$
\end{Example}


\begin{Theorem}{Интегралы Френеля}{}
Справедливо следующее равенство для несобственных интегралов:
$$ \int\limits_0^{+\infty} \cos(x^2) \, dx = \int\limits_0^{+\infty} \sin(x^2) \, dx = \frac{1}{2}\sqrt{\frac{\pi}{2}} $$
\end{Theorem}

\begin{center}
\begin{tikzpicture}[scale=1.2, >=latex]
    \definecolor{darktext}{HTML}{154360}
    \definecolor{colorCos}{HTML}{2874A6}
    \definecolor{colorSin}{HTML}{E67E22}

    % Оси
    \draw[->, thick, darktext] (-0.5, 0) -- (6, 0) node[right] {$x$};
    \draw[->, thick, darktext] (0, -1.5) -- (0, 1.5) node[above] {$y$};

    % Графики функций (пересчет в градусы для PGF)
    \draw[thick, colorCos, samples=200, domain=0:5.5] plot (\x, {cos(\x*\x * 180/pi)});
    \draw[thick, colorSin, samples=200, domain=0:5.5] plot (\x, {sin(\x*\x * 180/pi)});

    % Подписи
    \node[colorCos, font=\bfseries] at (2.5, 1.2) {$\cos(x^2)$};
    \node[colorSin, font=\bfseries] at (1.5, -1.2) {$\sin(x^2)$};

    % Пояснение
    \node[darktext, align=left, anchor=west] at (0, -2.4) {
        \small Осцилляции становятся всё чаще при $x \to \infty$. \\
        \small Площади соседних «полуволн» почти уничтожают друг друга, \\
        \small что и обеспечивает сходимость несобственного интеграла.
    };
\end{tikzpicture}
\end{center}

\begin{Proof}
Вычислим первый интеграл. Сделаем замену переменной $x^2 = t \implies 2x \, dx = dt \implies dx = \frac{dt}{2\sqrt{t}}$.
Тогда интеграл примет вид:
$$ I_c = \int\limits_0^{+\infty} \cos(x^2) \, dx = \frac{1}{2} \int\limits_0^{+\infty} \frac{\cos t}{\sqrt{t}} \, dt $$
Воспользуемся известным тождеством (из интеграла Эйлера-Пуассона):
$$ \frac{1}{\sqrt{t}} = \frac{2}{\sqrt{\pi}} \int\limits_0^{+\infty} e^{-tx^2} \, dx $$
Подставим это в наш интеграл:
$$ I_c = \frac{1}{2} \int\limits_0^{+\infty} \cos t \left( \frac{2}{\sqrt{\pi}} \int\limits_0^{+\infty} e^{-tx^2} \, dx \right) dt = \frac{1}{\sqrt{\pi}} \int\limits_0^{+\infty} \int\limits_0^{+\infty} e^{-tx^2} \cos t \, dx \, dt $$
Поменяем порядок интегрирования (по теореме Фубини/Тонелли):
$$ I_c = \frac{1}{\sqrt{\pi}} \int\limits_0^{+\infty} dx \left( \int\limits_0^{+\infty} e^{-tx^2} \cos t \, dt \right) $$
Внутренний интеграл — это табличный интеграл Лапласа (вычисляется двукратным интегрированием по частям):
$$ \int\limits_0^{+\infty} e^{-tx^2} \cos t \, dt = \frac{x^2}{(x^2)^2 + 1^2} = \frac{x^2}{x^4 + 1} $$
Возвращаемся к внешнему интегралу:
$$ I_c = \frac{1}{\sqrt{\pi}} \int\limits_0^{+\infty} \frac{x^2}{x^4 + 1} \, dx $$
Для вычисления этого рационального интеграла применим изящный трюк. Заметим, что замена $x = \frac{1}{u} \implies dx = -\frac{du}{u^2}$ даёт:
$$ \int\limits_0^{+\infty} \frac{x^2}{x^4 + 1} \, dx = \int\limits_{+\infty}^0 \frac{1/u^2}{1/u^4 + 1} \left( -\frac{du}{u^2} \right) = \int\limits_0^{+\infty} \frac{1}{u^4 + 1} \, du $$
Тогда удвоенный интеграл $I_c$ можно записать как сумму:
$$ 2 \sqrt{\pi} I_c = \int\limits_0^{+\infty} \frac{x^2}{x^4+1} dx + \int\limits_0^{+\infty} \frac{1}{x^4+1} dx = \int\limits_0^{+\infty} \frac{x^2 + 1}{x^4 + 1} \, dx = \int\limits_0^{+\infty} \frac{1 + 1/x^2}{x^2 + 1/x^2} \, dx $$
Заметим, что числитель является дифференциалом функции $\left(x - \frac{1}{x}\right)$, так как $d\left(x - \frac{1}{x}\right) = \left(1 + \frac{1}{x^2}\right) dx$. А знаменатель можно свернуть как полный квадрат: $x^2 + \frac{1}{x^2} = \left(x - \frac{1}{x}\right)^2 + 2$. 

Сделаем замену $z = x - \frac{1}{x}$. При $x \to +0$ имеем $z \to -\infty$, а при $x \to +\infty$ имеем $z \to +\infty$:
$$ 2 \sqrt{\pi} I_c = \int\limits_{-\infty}^{+\infty} \frac{dz}{z^2 + 2} = \left. \frac{1}{\sqrt{2}} \operatorname{arctg}\left(\frac{z}{\sqrt{2}}\right) \right|_{-\infty}^{+\infty} = \frac{1}{\sqrt{2}} \left( \frac{\pi}{2} - \left(-\frac{\pi}{2}\right) \right) = \frac{\pi}{\sqrt{2}} $$
Отсюда находим окончательный ответ:
$$ I_c = \frac{1}{2\sqrt{\pi}} \frac{\pi}{\sqrt{2}} = \frac{1}{2} \sqrt{\frac{\pi}{2}} $$
Интеграл от $\sin(x^2)$ вычисляется абсолютно аналогично, только внутренний интеграл даст $\frac{1}{x^4+1}$, что алгебраически приведет к точно такому же ответу.
\end{Proof}

\end{document}